% !TeX root = ../Book3.tex
\BookPart{椭圆函数、模形式与 \texorpdfstring{\(\theta\)}{\AnalysisPdfTheta} 函数}{b3:part:05}
% 本编导读（章级标题，不计入编号）
\BookPartGuide

双周期函数把复平面压缩到一个紧的环面；变动这个环面的格，又产生模形式；theta 级数则把格点计数写成可计算的 Fourier 展开。本编从具体函数出发，把复分析、代数曲线和整数表示连接起来。
\[
\begin{tikzcd}[column sep=large,row sep=large]
\text{第二十二章：椭圆函数}\arrow[r]\arrow[dr]&
\text{第二十三章：模形式}\arrow[d]\\
&\text{第二十四章：theta 函数与平方和}
\end{tikzcd}
\]
第二十二章提供 Weierstrass 函数、半周期与三次曲线。第二十三章研究格基变换和全模群，解释权与尖点的来源。第二十四章的四平方证明直接使用椭圆函数的半周期常数，同时以模形式观点解释生成函数的结构。Gaussian 的 Fourier 变换使 Poisson 求和成为 theta 参数反演的来源。
\ProseParagraphBreak
本编有两项容易混淆的正规化。模形式通常取 \(q=e^{2\pi\mathrm{i}\tau}\)，经典 theta 函数取 \(q=e^{\pi\mathrm{i}\tau}\)；两者都合理，但系数比较前必须换算。Weierstrass 方程采用 \(y^2=4x^3-g_2x-g_3\)，其判别式与正规化模判别式 \(\Delta\) 也相差明确的常数。
\ProseParagraphBreak
计算 \(j\)、\(E_4\)、\(\theta\) 时，可以反复追踪三个联系：周期条件怎样限制极点，变换律怎样限制 Fourier 系数，这些限制又怎样导出平方和公式。附录 E 从 Hecke 算子继续模形式主线，附录 F 从周期积分继续椭圆函数主线。

% !TeX root = ../../Book3.tex
\chapter{椭圆函数}
\label{b3:ch:22}
\BookChapterNumber{b3:ch:22}
% 本章摘要（不计入编号）
\begin{chapterabstract}
双周期性将复平面的局部亚纯理论放到紧环面上。本章先用对边积分约束零点、极点和留数，再通过正规收敛格和构造 Weierstrass 函数。其微分方程给出光滑三次曲线，函数的次数与半周期分析进一步证明环面和该曲线双全纯，并解释弦切群律。
\end{chapterabstract}

% 本章目标（不计入编号）
\begin{learninggoals}
能够验证格和的收敛和周期性，计算椭圆函数的零极点总数及除数限制；掌握 \(\wp\) 的微分方程、分歧点和缩放规律；能从复参数说明三次曲线的群律与函数域。
\end{learninggoals}

一个整函数若同时沿两条独立方向重复，就会被困在一个紧平行四边形中的模界内。要得到新的函数，必须允许极点；但这些极点又受到留数和为零的限制。Weierstrass 的构造从一个无留数的二阶极点出发，将它复制到全格，给出这些限制的最小实例。所需基础为留数定理、辐角原理和正规级数；Chandrasekharan 的《Elliptic Functions》\cite{Chandrasekharan1985Elliptic}可用于继续比较格、椭圆函数与模函数三种表述。

% !TeX root = ../../Book3.tex
\section{格与双周期函数}
\label{b3:sec:2201}

指数函数把一个平移周期变成单值函数，两个实线性无关的周期则把平面压成紧曲面。紧性迫使非常值函数具有极点，因此双周期理论自然属于亚纯函数，而非整函数。

\subsection{\texorpdfstring{\(\mathbb C\)}{C} 中的格}
\label{b3:subsec:220101}

\begin{definition}\label{b3:def:22-lattice}
若 \(\omega_1,\omega_2\in\mathbb C\) 在实数域上线性无关，则离散加法子群
\(\Lambda=\mathbb Z\omega_1+\mathbb Z\omega_2\) 称为复平面中的\term[ge]{格}[lattice]。有序生成元称为格的一组基；通常调整方向使 \(\operatorname{Im}(\omega_2/\omega_1)>0\)。
\end{definition}
实线性映射 \((x,y)\mapsto x\omega_1+y\omega_2\) 可逆，故存在 \(c,C>0\)，使
\(c\sqrt{m^2+n^2}\le|m\omega_1+n\omega_2|\le C\sqrt{m^2+n^2}\)。因此格点没有有限聚点，半径 \(R\) 内有 \(O(1+R^2)\) 个格点，且
\[
\sum_{0\ne\omega\in\Lambda}|\omega|^{-p}<\infty\quad(p>2).
\]
最后结论可按整数方框壳 \(\max(|m|,|n|)=j\) 求和；该壳有 \(8j\) 个点，故归结于 \(\sum j^{1-p}\)。这条估计将负责椭圆级数的正规收敛。

\subsection{基本平行四边形}
\label{b3:subsec:220102}

集合 \(P=\{s\omega_1+t\omega_2:0\le s,t<1\}\) 是\term[jiben-pingxing-sibianxing]{基本平行四边形}[fundamental parallelogram]。每个复数模 \(\Lambda\) 有唯一一个代表落在 \(P\)；闭平行四边形上对边点须按平移等同。

\ProseParagraphBreak

积分时常将整个平行四边形平移一个小量，使边界不经过给定亚纯函数的零极点。这样的平移确实存在：在固定有界范围内只有有限多个零极点，需要避开的平移参数是一族有限线段，其并集不含开集。计算最终只取决于模格的点，不取决于平行四边形的起点。

\subsection{双周期函数}
\label{b3:subsec:220103}

若亚纯函数 \(f\) 满足 \(f(z+\omega)=f(z)\) 对全部 \(\omega\in\Lambda\) 成立，就称它关于 \(\Lambda\) 双周期。检验两组生成元足够，因正向平移、反向平移和有限次复合给出全部格平移。

\ProseParagraphBreak

周期群可能比指定的格更大。例如相对于 \(\Lambda\) 的函数也可相对于子格 \(2\Lambda\) 讨论，但后一基本区域中极点数变成原来的四倍。凡说“次数”或“模格的零点”，必须先固定采用哪个格；不能无说明地用函数的全部周期替换给定格。

\subsection{椭圆函数}
\label{b3:subsec:220104}

上述亚纯双周期函数称为\term[tuoyuan-hanshu]{椭圆函数}[elliptic function]。常函数也包含在这个定义中。
\begin{proposition}\label{b3:prop:22-entire-constant}
没有非常值的整椭圆函数。
\end{proposition}
\begin{proof}
整函数在闭基本平行四边形上有界；周期性把这个界传到全平面。Liouville 定理使它为常数。
\end{proof}
仅有一个实周期时结论失败：\(e^{2\pi\mathrm{i}z}\) 是周期为一的非常值整函数，其商空间仍有无穷远端。两个独立周期使商空间紧致，才造成上述刚性。

\ProseParagraphBreak

离散性也能从周期本身反向检验。设非常值亚纯函数有非零周期 \(\omega_j\to0\)。在任一没有极点的小圆盘内，差商
\(\bigl(f(z+\omega_j)-f(z)\bigr)/\omega_j\) 恒为零，极限给 \(f'(z)=0\)。因此函数在一个开集上常值，亚纯恒等定理使它处处常值，矛盾。所以非常值亚纯函数的周期群一定离散。特别地，三个实线性相关而产生稠密整数加法群的周期，不会产生比椭圆函数更自由的非常值函数。
\ProseParagraphBreak
所谓两个周期“独立”，必须是在实数域上独立。若两个非零实周期的比为无理数，它们生成的群在实轴稠密，上一论证同样迫使常数；若比为有理数，它们只生成一维离散周期群。只有二维实独立的情形给紧的基本区域。在判定题目中的函数是否真正双周期时，应先检查这个线性代数条件，而不是数出公式中出现了两个平移数值。
\ProseParagraphBreak
格和的收敛指数二也有几何含义。一个大圆中格点数不仅是 \(O(R^2)\)，还与面积同阶：用同样大小的基本平行四边形铺平面，圆内完全包含的格胞覆盖一个稍小圆，与圆相交的格胞则落在一个稍大圆。格胞面积固定，故计数介于 \(cR^2\) 与 \(CR^2\) 之间。于是未作相消的 \(\sum|\omega|^{-2}\) 在二维格上发散，而 \(\sum|\omega|^{-3}\) 收敛。Weierstrass 级数中的减项正是跨过这个临界指数所必需的修正。

% !TeX root = ../../Book3.tex
\section{椭圆函数的基本结构}
\label{b3:sec:2202}

平行四边形的对边积分相消，把局部极点资料转成全局限制。以下各积分均选取无零极点的边界，并把零极点按模格的重数计数。

\subsection{零点与极点}
\label{b3:subsec:220201}

\begin{proposition}\label{b3:prop:22-zero-pole-count}
非恒零椭圆函数在一个基本区域内的零点总重数等于极点总阶数。
\end{proposition}
\begin{proof}
对数导数 \(f'/f\) 仍双周期，对边方向相反，所以边界积分为零。辐角原理将积分除以 \(2\pi\mathrm{i}\) 识别为零点数减极点数。商空间紧致保证两类点都有限；等式不包含恒零函数。
\end{proof}
对任意有限值 \(a\)，同一结论应用于 \(f-a\)，说明取值 \(a\) 的原像按重数计数等于极点总阶数。这个整数对 \(a\) 不变，给出椭圆函数作为覆盖映射的次数。

\subsection{留数和}
\label{b3:subsec:220202}

\begin{proposition}\label{b3:prop:22-residue-sum}
椭圆函数在一个基本区域内的留数之和为零。
\end{proposition}
\begin{proof}
将上一证明的被积函数改成 \(f\)，用留数定理而不是辐角原理；对边仍因周期性相消。
\end{proof}
因此一个基本区域中不能恰有一个一阶极点，因为一阶极点的留数非零。但一个二阶极点可以出现，其留数可以为零。下一节的 \(\wp\) 函数正是次数二的最简单实例。两个一阶极点的留数必须相反；这个条件将允许用两个平移的对数导数相减来构造它。

\subsection{除数关系}
\label{b3:subsec:220203}

\begin{proposition}\label{b3:prop:22-divisor-necessary}
设非恒零椭圆函数的零点代表为 \(a_j\)，极点代表为 \(b_k\)，均按重数重复。则
\(\sum_j a_j-\sum_k b_k\in\Lambda\)。
\end{proposition}
\begin{proof}
令 \(h=f'/f\)。留数定理给差值等于 \((2\pi\mathrm{i})^{-1}\int_{\partial P}zh(z)\,\mathrm{d}z\)。将对边配对，得到两个基周期乘以沿另两条边的 \(h\) 积分，并附相应正负号。每一条这样的边，其端点处 \(f\) 值相同；沿边连续选对数，终点与起点对数相差 \(2\pi\mathrm{i}\) 的整数倍。因此所得差是 \(\omega_1,\omega_2\) 的整数线性组合。换代表只改变格元素，不影响结论。
\end{proof}
这一条件比零极点数相等更强。例如指定单零点 \(a\) 和单极点 \(b\) 会要求 \(a-b\in\Lambda\)，这时零极点在商空间上相同并相互抵消，所以次数一不可能出现。给定除数的存在性则还需要构造函数，不能从必要条件倒推；下文先建立足以进行这种构造的 \(\wp\) 及其积分。

\subsection{椭圆函数的次数}
\label{b3:subsec:220204}

若非常值椭圆函数在基本区域有总阶数 \(d\) 的极点，就称 \(d\) 为它相对于 \(\Lambda\) 的\term[tuoyuan-hanshu-cishu]{次数}[degree of an elliptic function]。局部展开 \(f(z)-a=c(z-z_0)^m+\cdots\) 中的 \(m\) 是该纤维点的重数；极点处改用 \(1/f\) 计算。

\ProseParagraphBreak

在非分歧值附近，各个局部逆分支连续存在，故原像个数恰为 \(d\)。分歧值上不同逆分支可能合并，但总重数不变。这样，辐角原理的整数不只是记账工具，而是刻画了映射的每条纤维。次数二的偶函数预示了一对原像 \(z,-z\)，它们只在半周期点合并。

% !TeX root = ../../Book3.tex
\section{Weierstrass \texorpdfstring{\(\wp\)}{℘} 函数}
\label{b3:sec:2203}

直接相加 \((z-\omega)^{-2}\) 的常数尾项只有临界的二次衰减。减去 \(\omega^{-2}\) 后，衰减提高一次，得到可在紧集上求导的格和。

\subsection{定义与收敛性}
\label{b3:subsec:220301}

\begin{definition}\label{b3:def:22-wp}
对格 \(\Lambda\)，定义\term[weierstrass-wp-hanshu]{Weierstrass 椭圆函数}[Weierstrass elliptic function]
\[
\wp(z)=\frac1{z^2}+\sum_{0\ne\omega\in\Lambda}
\left(\frac1{(z-\omega)^2}-\frac1{\omega^2}\right).
\]
\end{definition}
\symidx{\(\wp(z;\Lambda)\)：Weierstrass 椭圆函数}
若 \(|z|\le R\)、\(|\omega|>2R\)，括号内的模不超过 \(C_R|\omega|^{-3}\)。因此避开格点的紧集上正规收敛。在每个格点恰有一个二阶极点，主部为到该格点距离的倒数平方。把 \(\omega\) 与 \(-\omega\) 配对证明 \(\wp\) 为偶函数，且原点 Laurent 展开为
\[
\wp(z)=z^{-2}+3G_4z^2+5G_6z^4+7G_8z^6+\cdots,
\quad G_{2k}=\sum_{\omega\ne0}\omega^{-2k}.
\]
此式先在最短非零周期以内的圆盘成立，由几何级数展开和绝对收敛换序得到。

\subsection{\texorpdfstring{\(\wp'\)}{\AnalysisPdfWeierstrassP'}}
\label{b3:subsec:220302}

逐项求导后，\(\wp'(z)=-2\sum_{\omega\in\Lambda}(z-\omega)^{-3}\)。这个级数绝对正规收敛，平移后可直接重排，因此 \(\wp'\) 双周期。对每个基周期 \(\omega_j\)，差 \(\wp(z+\omega_j)-\wp(z)\) 为常数；取 \(z=-\omega_j/2\)，两点均非格点，偶性使常数为零。两基周期生成全部格，故 \(\wp\) 也双周期。

\ProseParagraphBreak

这一论证解释为何不直接在原级数中重排两个各自发散的二次格和。正规收敛只对扣除后的整项成立；先求导得到绝对收敛的三次格和，再积分回常数，是避免错误重排的关键。

\subsection{微分方程}
\label{b3:subsec:220303}

定义 \(g_2=60G_4\)、\(g_3=140G_6\)。
\begin{theorem}[Weierstrass 微分方程]\label{b3:thm:22-wp-equation}
函数 \(\wp\) 满足 \((\wp')^2=4\wp^3-g_2\wp-g_3\)。
\end{theorem}
\begin{proof}
写 \(\wp=z^{-2}+az^2+bz^4+O(z^6)\)，其中 \(a=3G_4\)、\(b=5G_6\)。展开得
\((\wp')^2-4\wp^3=-20az^{-2}-28b+O(z^2)\)。因此
\((\wp')^2-4\wp^3+g_2\wp+g_3\) 在原点可去且值为零。周期性使它在每个格点可去，其他点本来全纯；整椭圆函数必常数，原点值又使常数为零。
\end{proof}
再次求导并先在 \(\wp'\ne0\) 处约去，再用恒等定理，得到 \(\wp''=6\wp^2-g_2/2\)。它与初始主部一起递推确定全部 Laurent 系数，也给出下一节代数曲线的切线斜率。

\subsection{零点与极点}
\label{b3:subsec:220304}

模 \(\Lambda\) 的三个非零半周期为 \(\omega_1/2,\omega_2/2,(\omega_1+\omega_2)/2\)。奇性与周期性使 \(\wp'\) 在这三点为零。由于 \(\wp'\) 的唯一极点阶数为三，零点计数说明三点均为单零点，没有其他零点。

\ProseParagraphBreak

记它们的 \(\wp\) 值为 \(e_1,e_2,e_3\)。这些值两两不同：若两点具有同一值，则 \(\wp-e_j\) 在各点至少二阶消失，零点总数至少四，与其次数二矛盾。因此
\[
4X^3-g_2X-g_3=4(X-e_1)(X-e_2)(X-e_3),
\quad e_1+e_2+e_3=0.
\]
一般值的两个原像恰为 \(z,-z\)。若 \(\wp(z)=\wp(w)\)，\(\wp-\wp(w)\) 的次数二已经被 \(w,-w\) 耗尽；在半周期时这两个点合为一个二重点。这是后面证明解析表示单射的依据。

\ProseParagraphBreak

把 \(-\wp\) 积分一次会产生\term[weierstrass-zeta-hanshu]{Weierstrass zeta 函数}[Weierstrass zeta function]。为避免与后面的 Riemann zeta 函数混淆，记
\[
 Z_\Lambda(z)=\frac1z+\sum_{\omega\in\Lambda\setminus\{0\}}
 \left(\frac1{z-\omega}+\frac1\omega+\frac z{\omega^2}\right).
\]
在任意避开格点的紧集上，括号内为 \(O_K(|\omega|^{-3})\)，所以级数正规收敛且可逐项微分，得到 \(Z_\Lambda'=-\wp\)。对称配对 \(\omega,-\omega\) 说明它是奇函数。由于导数周期，
\(Z_\Lambda(z+\omega_j)-Z_\Lambda(z)=\eta_j\) 为常数，称为相应准周期。这两个准周期不可能同时为零，否则会得到只有一个单极点的椭圆函数，违反留数和。
\ProseParagraphBreak
对不经过格点、内部恰含一个格点的基本平行四边形积分 \(Z_\Lambda\)，留数定理给 \(2\pi\mathrm{i}\)。把上边与下边、右边与左边分别配对，准周期差给
\begin{equation}\label{b3:eq:22-legendre-relation}
 \eta_1\omega_2-\eta_2\omega_1=2\pi\mathrm{i}.
\end{equation}
这里基的方向为正。反转基顺序会同时反转围道方向和左侧符号。这一常数不是任意正规化，而是一个单位留数沿基本区域累计的结果。
\ProseParagraphBreak
再定义整函数
\[
 \Sigma_\Lambda(z)=z\prod_{\omega\ne0}
 \left(1-\frac z\omega\right)
 \exp\left(\frac z\omega+\frac{z^2}{2\omega^2}\right).
\]
每个尾因子的对数从三次项开始，所以格点三次幂和保证正规收敛。零点恰为格点且均单，原点导数为一，直接微分得 \(\Sigma_\Lambda'/\Sigma_\Lambda=Z_\Lambda\)。奇性也由 \(Z_\Lambda\) 的奇性确定：\(-\Sigma_\Lambda(-z)\) 与 \(\Sigma_\Lambda(z)\) 具有相同对数导数及原点导数，故相等。这就是\term[weierstrass-sigma-hanshu]{Weierstrass sigma 函数}[Weierstrass sigma function]，此处用大写符号避免与算术函数的约数和记号冲突。
\ProseParagraphBreak
对数导数的准周期性给商
\(\Sigma_\Lambda(z+\omega_j)/(e^{\eta_jz}\Sigma_\Lambda(z))\) 为非零常数。取 \(z=-\omega_j/2\)，用奇性及半周期不是格点，确定该常数为 \(-e^{\eta_j\omega_j/2}\)。于是
\[
 \Sigma_\Lambda(z+\omega_j)
 =-e^{\eta_j(z+\omega_j/2)}\Sigma_\Lambda(z).
\]
这条公式同时保留了指数和负号。若仅比较导数，会遗漏乘法常数；后面构造双周期商时，漏掉负号就可能产生错误的半周期。

\begin{theorem}[椭圆除数的实现]\label{b3:thm:22-divisor-realization}
设有限点列 \(a_1,\ldots,a_d\)、\(b_1,\ldots,b_d\) 按重数重复，且两列模格不相交。存在零点恰为前一列、极点恰为后一列的椭圆函数，当且仅当
\(\sum a_j-\sum b_j\in\Lambda\)。该函数相差一个非零常数唯一。
\end{theorem}
\begin{proof}
必要性已由基本区域的除数关系证明。对充分性，写 \(D=\sum a_j-\sum b_j=m\omega_1+n\omega_2\)，取 \(c=m\eta_1+n\eta_2\)，令
\[
 f(z)=e^{cz}\prod_{j=1}^d\frac{\Sigma_\Lambda(z-a_j)}{\Sigma_\Lambda(z-b_j)}.
\]
在平移 \(\omega_i\) 时，分子分母各有 \(d\) 个负号，恰好相消；总乘子为 \(e^{c\omega_i-\eta_iD}\)。由式~\eqref{b3:eq:22-legendre-relation}，指数在 \(i=1\) 时为 \(-2\pi\mathrm{i}n\)，在 \(i=2\) 时为 \(2\pi\mathrm{i}m\)，所以两乘子都为一。零极点由整乘积的单零点完全确定。若有两个实现函数，其商无零点、无极点且双周期，故为非零常数。
\end{proof}
构造中改变某个点的格代表，会改变指数因子与前面的常数，却不改变最终除数和按常数唯一的函数。给定零极点的总和条件因此是商环面上的条件，不能依赖选取平行四边形时的偶然代表。

% !TeX root = ../../Book3.tex
\section{椭圆曲线的解析表示}
\label{b3:sec:2204}

微分方程把一对亚纯函数放在一条代数曲线上。但“满足方程”尚不等于得到曲面的参数化，还需要验证无穷远点、局部坐标、单射和满射。

\subsection{三次方程}
\label{b3:subsec:220401}

把仿射方程 \(y^2=4x^3-g_2x-g_3\) 齐次化为
\[
E_\Lambda:\quad Y^2Z=4X^3-g_2XZ^2-g_3Z^3.
\]
在 \(Z=0\) 上只有一点 \(O=[0:1:0]\)。有限点的奇点须同时满足 \(y=0\)、\(12x^2-g_2=0\) 及曲线方程，等价于右侧三次多项式有重根。前节已经排除重根；在 \(O\) 处，对 \(Z\) 的偏导为一，也非奇点。因此这是一个光滑紧复曲线。

\subsection{判别式}
\label{b3:subsec:220402}

多项式 \(4X^3-g_2X-g_3\) 的无重根条件等价于
\(D=g_2^3-27g_3^2\ne0\)。这里把 \(D\) 称为该 Weierstrass 方程的\term[tuoyuan-quxian-panbieshi]{判别式}[discriminant]；有的代数规范另带固定常数，使用具体公式时应核对归一化。

\ProseParagraphBreak

换格为 \(a\Lambda\) 时，由绝对收敛格和可得
\[
\wp(az;a\Lambda)=a^{-2}\wp(z;\Lambda),\quad
g_2(a\Lambda)=a^{-4}g_2(\Lambda),\quad g_3(a\Lambda)=a^{-6}g_3(\Lambda).
\]
因此判别式的缩放次数为十二。这个整数将成为下一章模形式权十二的来源，并不是任意选择的指数。

\subsection{\texorpdfstring{\(\mathbb C/\Lambda\)}{C/\AnalysisPdfCapitalLambda}}
\label{b3:subsec:220403}

商空间 \(\mathbb C/\Lambda\) 由平行四边形粘合对边得到。每个小于半个最短周期的圆盘都与其非平凡格平移不相交，其投影给出复坐标。坐标变换为平移，因而全纯；闭平行四边形覆盖商空间，故它紧致。

\ProseParagraphBreak

对 \(z\notin\Lambda\) 定义 \(\Phi(z)=[\wp(z):\wp'(z):1]\)，对格点定义 \(\Phi(z)=O\)。在原点附近除以 \(\wp'\)，坐标为
\([\wp/\wp':1:1/\wp']=[-z/2+O(z^5):1:-z^3/2+O(z^7)]\)，所以该赋值全纯，且在原点为浸入。它按周期下降为商曲面上的全纯映射。

\subsection{解析统一化}
\label{b3:subsec:220404}

\begin{theorem}\label{b3:thm:22-uniformization}
映射 \(\Phi:\mathbb C/\Lambda\to E_\Lambda\) 是双全纯映射。
\end{theorem}
\begin{proof}
若两个有限点具有同一 \(x\) 坐标，则对应参数为 \(z\) 与 \(-z\)。它们的 \(y\) 坐标相反；若还相等，则 \(\wp'(z)=0\)，于是 \(z\) 为半周期，\(z\equiv-z\pmod\Lambda\)。无穷远点只来自格点，故单射。对任意 \((x,y)\in E_\Lambda\)，次数二保证 \(\wp\) 取到 \(x\)；微分方程给导数值为两个平方根之一，换 \(z\) 为 \(-z\) 即得到指定 \(y\)，所以满射。
在 \(\wp'\ne0\) 处可用 \(x\) 作局部坐标；在半周期处，\(\wp'\) 的零点为单零点，故 \(y\) 是局部坐标；在 \(O\) 处用上面 \(\wp/\wp'\) 的非零一阶项。逆函数定理在三种情形都给全纯局部逆，故全局逆也全纯。
\end{proof}
这里证明的是由格构造的曲线的解析表示。反过来从任意光滑复三次曲线恢复周期格，需要积分全纯微分并核对周期；这一一般方向可结合\appref{b3:app:A}的 Riemann 曲面理论继续学习，不能仅由上述显式代入推出。

\subsection{群律的解释}
\label{b3:subsec:220405}

通过 \(\Phi\) 将商群加法搬到曲线上，\(O\) 为单位元，\((x,y)\) 的逆为 \((x,-y)\)。其几何意义可由除数关系读出：一般直线与三次曲线交于三点，按重数记为 \(P,Q,R\)，相应线性函数在环面上只有原点三阶极点，因此三个参数之和模格为零。故 \(P+Q\) 是第三交点关于 \(x\) 轴的对称点。

\ProseParagraphBreak

当两点的横坐标不同时，直线斜率为 \(m=(y_2-y_1)/(x_2-x_1)\)，代入三次方程并比较根的和得
\[
x(P+Q)=m^2/4-x_1-x_2.
\]
对应纵坐标取第三交点的负值。两点合并时斜率改为 \((12x_1^2-g_2)/(2y_1)\)；纵坐标为零的切线情形给 \(2P=O\)。这些公式来自已经存在的群结构，结合律由参数相加自动保证，无需再对繁长有理式作三重展开。

\ProseParagraphBreak

还可确定整个椭圆函数域。若 \(f\) 为偶椭圆函数，它在 \(\wp\) 的每条纤维上取同一值，因此下降为球面上的亚纯函数，必为 \(\wp\) 的有理函数。在半周期处，局部参数变换 \(u\mapsto-u\) 使偶 Laurent 展开只含 \(u^{2k}\)，保证下降函数确实亚纯。对一般 \(f\)，偶部 \((f(z)+f(-z))/2\) 为 \(R(\wp)\)，奇部除以 \(\wp'\) 后也为偶函数，故
\[
f(z)=R(\wp(z))+\wp'(z)S(\wp(z))
\]
其中 \(R,S\) 有理。这说明曲线方程已经包含全部椭圆函数，而不只是两个偶然选择的函数。

\ProseParagraphBreak

三次曲线上的微分也可以通过参数直接理解。在非分歧点，\(x=\wp(z)\)、\(y=\wp'(z)\)，所以
\[
 \frac{\mathrm{d}x}{y}=\mathrm{d}z.
\]
在半周期处，两边仍全纯：若 \(z=z_0+t\)，则 \(x=e_j+\tfrac12\wp''(z_0)t^2+O(t^3)\)、\(y=\wp''(z_0)t+O(t^2)\)，分子与分母的一阶因子相消，商趋于 \(\mathrm{d}t\)。在无穷远，\(x=z^{-2}+O(z^2)\)、\(y=-2z^{-3}+O(z)\)，又得 \(\mathrm{d}x/y=(1+O(z^4))\mathrm{d}z\)。因此这个看似含分母的微分在整个紧曲线上无极点、无零点。
\ProseParagraphBreak
环面上的任意全纯微分提升到平面后写成 \(h(z)\mathrm{d}z\)。对格平移不变要求 \(h\) 是整椭圆函数，故常数。因此全纯微分空间恰为一维，由 \(\mathrm{d}z\) 生成。沿两条基本闭曲线积分得到 \(\omega_1,\omega_2\)，周期格遂从微分中恢复出来。这个结论是在已知格环面上直接计算；一般紧 Riemann 曲面的微分空间维数需要附录 A 的进一步理论。
\ProseParagraphBreak
二重映射 \(\wp\) 的四个分歧点也由局部式看得清楚。三个非零半周期在有限值 \(e_1,e_2,e_3\) 上有局部参数平方；原点的二阶极点在目标无穷远坐标 \(1/\wp\) 下同样是平方。其他点导数非零，局部为双全纯。因而三次曲线写成 \(y^2=4(x-e_1)(x-e_2)(x-e_3)\) 时，选择一个 \(x\) 通常有两种符号的 \(y\)，绕一个单根则交换这两种选择。代数方程的平方根分支与环面上的半周期是同一现象的两种坐标描述。
\ProseParagraphBreak
矩形格能进一步给出实图像。设 \(\omega_1>0\)、\(\omega_2\) 纯虚且虚部正，格关于共轭不变，由级数有 \(\overline{\wp(z)}=\wp(\bar z)\)。实轴上 \(\wp\) 实值，从原点附近的正无穷开始下降；在 \((0,\omega_1/2)\) 上导数没有零点，所以直到半周期才达到第一个临界值。它必须是三个实根中的最大者，否则下降途中先遇到别的根就会先有导数零点。沿虚轴，\(\wp(\mathrm{i}t)\sim-t^{-2}\)，从负无穷上升，首次临界值同理是最小根。第三个半周期给中间根。三个值确为实数，因为共轭半周期与原半周期相差格点。由此不用数值画图，就能确定矩形格的三个分歧值顺序及实椭圆曲线的两支。
\ProseParagraphBreak
上述顺序依赖矩形格假设。对一般倾斜格，三个分歧值可以全为非实复数，谈“最大根”便没有意义；但无重根、四个二阶分歧点和全纯微分的结论仍成立。几何图像有助于选择积分路径，却不能替代复参数情况下的局部代数检验。

% 本章小结（不计入编号）
\begin{chaptersummary}
格的二维点数估计保证扣除后的级数收敛，对边相消则把局部 Laurent 数据变成全局整数关系。\(\wp\) 与 \(\wp'\) 的次数分别为二与三，它们的三次方程给出完整的环面解析表示。缩放格改变 \(g_2,g_3\) 的权，而保持相应曲线的复结构；下一章将研究这种变化。
\end{chaptersummary}


\BookExercises

% 八题自拟：补充格的对称性、主部构造及曲线计算。
\begin{exercise}\label{b3:ex:22-1}
证明格的另一组正向基恰由 \(\mathrm{SL}_2(\mathbb Z)\) 的整系数基变换给出。
\end{exercise}
\begin{solution}
两组基互为整系数线性组合，故变换矩阵及其逆均整，行列式为 \(\pm1\)。保持方向使行列式为正，即一。反之这样的矩阵有整逆，生成同一格并保持方向。
\end{solution}
\begin{exercise}\label{b3:ex:22-2}
设 \(f\) 只有模格原点处一个至多二阶极点。证明 \(f=A\wp+B\)。
\end{exercise}
\begin{solution}
留数和为零排除主部中的一阶项。取 \(A\) 为二阶项系数，则 \(f-A\wp\) 无极点，整椭圆函数定理使其为常数。
\end{solution}
\begin{exercise}\label{b3:ex:22-3}
对方格 \(\Lambda=\mathbb Z+\mathrm{i}\mathbb Z\) 证明 \(g_3=0\)，并求 \(\wp(\mathrm{i}z)\) 与 \(\wp(z)\) 的关系。
\end{exercise}
\begin{solution}
由于 \(\mathrm{i}\Lambda=\Lambda\)，缩放给 \(g_3=\mathrm{i}^{-6}g_3=-g_3\)，故零；同理 \(\wp(\mathrm{i}z)=-\wp(z)\)。判别式非零进一步给 \(g_2\ne0\)。
\end{solution}
\begin{exercise}\label{b3:ex:22-4}
对 \(\Lambda=\mathbb Z+e^{\pi\mathrm{i}/3}\mathbb Z\) 证明 \(g_2=0\)、\(g_3\ne0\)。
\end{exercise}
\begin{solution}
记 \(a=e^{\pi\mathrm{i}/3}\)，有 \(a^2=a-1\)，故 \(a\Lambda=\Lambda\)。于是 \(g_2=a^{-4}g_2\)，而 \(a^{-4}\ne1\)，所以 \(g_2=0\)。判别式非零排除 \(g_3=0\)。
\end{solution}
\begin{exercise}\label{b3:ex:22-5}
证明 \(\wp(z)-\wp(a)\) 的除数为 \((a)+(-a)-2(0)\)，并区分 \(2a\in\Lambda\) 的情形。
\end{exercise}
\begin{solution}
假定 \(a\notin\Lambda\)。偶性给两个零点。若它们不同，次数二使两者均单；若为半周期，两点合并，导数零且二阶导数非零，故是二重零点。唯一极点在原点二阶。
\end{solution}
\begin{exercise}\label{b3:ex:22-6}
设椭圆函数 \(f\) 在原点的主部为 \(az^{-4}+bz^{-3}+cz^{-2}\)，且模格没有其他极点。用 \(\wp,\wp',\wp''\) 表示它。
\end{exercise}
\begin{solution}
三个函数的主部分别为 \(z^{-2},-2z^{-3},6z^{-4}\)，均无一阶项。于是 \(f-a\wp''/6+b\wp'/2-c\wp\) 无极点，为常数。再用微分方程可将 \(\wp''\) 换为 \(6\wp^2-g_2/2\)。
\end{solution}
\begin{exercise}\label{b3:ex:22-7}
求环面上满足 \(nz=0\) 的点数，并说明 \(n=2\) 时它们在三次曲线上的位置。
\end{exercise}
\begin{solution}
点由 \((a\omega_1+b\omega_2)/n\) 表示，\(0\le a,b<n\)，实线性独立保证这些 \(n^2\) 个点互异。\(n=2\) 时为单位元与三个非零半周期，对应 \(O\) 和 \((e_j,0)\)。
\end{solution}
\begin{exercise}\label{b3:ex:22-8}
设 \(a\) 不是格点或半周期。证明
\[
 Z_\Lambda(z+a)-Z_\Lambda(z)-Z_\Lambda(a)
 =\frac12\frac{\wp'(z)-\wp'(a)}{\wp(z)-\wp(a)}.
\]
\end{exercise}
\begin{solution}
两边关于 \(z\) 都是椭圆函数。右侧可能的极点在 \(0,\pm a\)：在 \(a\) 处分子分母同阶消去，在 \(-a\) 处留数为一，在0处主部为 \(-1/z\)。左侧在 \(-a,0\) 的留数也为一、负一，故两边之差无极点而为常数。令 \(z\to0\)，左侧加 \(1/z\) 后趋零；右侧利用 \(\wp=z^{-2}+O(z^2)\)、\(\wp'=-2z^{-3}+O(z)\) 展开，加 \(1/z\) 后也趋零。常数因而为零。半周期情形可经参数极限另外解释，但题目假设避免同时处理合并极点。
\end{solution}

% !TeX root = ../../Book3.tex
\chapter{Eisenstein 级数与模形式}
\label{b3:ch:23}
\BookChapterNumber{b3:ch:23}
% 本章摘要（不计入编号）
\begin{chapterabstract}
格的换基给出模群作用，格和给出带权变换的全纯函数。本章建立全模群的基本域、Eisenstein 级数和模形式，证明价公式并由此求出模形式环。判别式与 \(j\)-不变量把这些解析计算解释为复环面及其参数空间的几何。
\end{chapterabstract}

% 本章目标（不计入编号）
\begin{learninggoals}
能够推导模群作用与格比例的关系，检查尖点全纯性，计算 Eisenstein 级数的 Fourier 系数；能够在价公式中正确计入固定点的权重，用消失阶判断恒等式，并解释 \(j\) 对复环面的分类。
\end{learninggoals}

同一个格换一组基，椭圆函数代表的曲线并没有变化，但 \(\tau\)、\(g_2\)、\(g_3\) 的数值会按不同规律改变。哪些量只反映坐标，哪些量真正区别复结构，是本章要回答的问题。前章的格缩放、留数计算和三次曲线表示构成直接基础。Apostol 的《Modular Functions and Dirichlet Series in Number Theory》\cite{Apostol1990Modular}提供格和与模函数的系统叙述；本文价公式及其代数后果也可与 Zagier 的《Elliptic Modular Forms and Their Applications》\cite{Zagier2008Elliptic}相互核对。

% !TeX root = ../../Book3.tex
\section{模群}
\label{b3:sec:2301}

前章把一个格写成 \(\Lambda=\mathbb Z\omega_1+\mathbb Z\omega_2\)，并由它构造复环面与椭圆函数。若选取正向基，再除以 \(\omega_1\)，格就化为
\[
 \Lambda_\tau=\mathbb Z+\mathbb Z\tau,
 \quad \tau=\frac{\omega_2}{\omega_1}\in\mathfrak H,
 \quad \mathfrak H=\{\tau\in\mathbb C:\operatorname{Im}\tau>0\}.
\]
这里的上半平面与前文的 \(\mathbb H\) 是同一个区域；本章用 \(\mathfrak H\) 突出它的模参数角色。参数不是唯一的：改变同一个格的基，就会改变 \(\tau\)。模群先回答哪些参数表示同一比例格，再为后面的函数变换律提供准确的换基因子。

\subsecref{b3:subsec:110404}已从覆盖的角度讨论模群。本节恢复所需的矩阵与基本域事实，并进一步确定固定点和尖点的局部坐标。特别要记住，模群含有固定点，商映射不能直接套用自由群作用的覆盖结论。

\subsection{\texorpdfstring{\(\mathrm{SL}_2(\mathbb Z)\)}{SL2(Z)}}
\label{b3:subsec:230101}

\Needspace{6\baselineskip}
\begin{definition}\label{b3:def:23-modular-group}
行列式为一的二阶整数矩阵组成群
\[
 \operatorname{SL}_2(\mathbb Z)
 =\left\{\begin{pmatrix}a&b\\c&d\end{pmatrix}:
       a,b,c,d\in\mathbb Z,\ ad-bc=1\right\},
\]
群运算为矩阵乘法。本章在记录带权变换律时使用这个矩阵群；在只讨论其分式线性作用时，模群指它的有效作用群
\(\operatorname{PSL}_2(\mathbb Z)=\operatorname{SL}_2(\mathbb Z)/\{\pm I\}\)，与\subsecref{b3:subsec:110404}的约定一致。
\end{definition}
\termidx[moqun]{模群}\foreignidx{modular group}

逆矩阵为 \(\left(\begin{smallmatrix}d&-b\\-c&a\end{smallmatrix}\right)\)，所以整数条件在取逆后仍保持。行列式一还给一个将反复使用的事实：底行 \((c,d)\) 的两个整数互素。反过来，若 \(c,d\) 互素，Bézout 恒等式便给整数 \(a,b\) 使 \(ad-bc=1\)，从而把这行补成一个群元素。

记
\[
 T=\begin{pmatrix}1&1\\0&1\end{pmatrix},
 \quad S=\begin{pmatrix}0&-1\\1&0\end{pmatrix}.
\]
\cref{b3:prop:11-modular-generators}的 Euclid 算法也适用于这里的矩阵版本。若第一列为 \((a,c)^{\mathsf T}\)、\(c\ne0\)，作带余除法 \(a=qc+r\)、\(|r|<|c|\)，左乘 \(T^{-q}\) 和 \(S\) 将它依次变成
\[
 (a,c)^{\mathsf T}\longmapsto(r,c)^{\mathsf T}
 \longmapsto(-c,r)^{\mathsf T}.
\]
下左元素的绝对值严格下降，有限步后为零。所得矩阵的两个对角元素只能同为 \(1\) 或同为 \(-1\)，故为某个 \(\pm T^n\)。由于 \(-I=S^2\)，逆向读取操作就把原矩阵写成 \(S,T\) 及其逆的有限乘积。这个生成结论以后能把“对所有整数矩阵检查变换律”缩成对两个基本变换的检查。

\subsection{\texorpdfstring{\(\mathrm{PSL}_2(\mathbb Z)\)}{PSL2(Z)}}
\label{b3:subsec:230102}

矩阵 \(\gamma=\left(\begin{smallmatrix}a&b\\c&d\end{smallmatrix}\right)\) 决定分式线性映射
\[
 \gamma\tau=\frac{a\tau+b}{c\tau+d}.
\]
同时改变四个系数的符号不改变映射，所以 \(\gamma\) 和 \(-\gamma\) 作用相同。这个重复恰好是全部重复：若 \(\gamma\tau=\tau\) 在上半平面恒成立，则
\[
 c\tau^2+(d-a)\tau-b=0
\]
是一个恒为零的多项式，故 \(c=b=0\)、\(a=d\)；再由行列式条件得 \(a=d=\pm1\)。两个矩阵作用相同，便可对它们的商使用这个结论，说明它们只差整体负号。

因而 \(\operatorname{PSL}_2(\mathbb Z)\) 的每个元素是一个二元集合 \(\{\gamma,-\gamma\}\)，且它在上半平面上的作用忠实。这里取的是矩阵群的商；稍后还要对上半平面中的点按轨道取商，两种操作的对象不同。

矩阵等式为
\[
 S^2=-I,\quad (ST)^3=-I,
\]
所以在射影群中，\(S\) 的阶为二，\(ST\) 的阶为三。乘积 \(ST\) 表示先作 \(T\)、后作 \(S\)，其分式为 \(-1/(\tau+1)\)。在带权函数的变换律中仍保留矩阵 \(-I\)：虽然它不移动 \(\tau\)，它的权因子未必是一，这将直接排除本章的奇数权形式。

\subsection{在 \texorpdfstring{\(\mathfrak H\)}{上半平面} 上的作用}
\label{b3:subsec:230103}

先核对作用确实留在所需区域内。若 \(c\ne0\)，分母的零点 \(-d/c\) 在实轴上；若 \(c=0\)，则 \(d=\pm1\)。因此分母在 \(\mathfrak H\) 不为零。将分式与其共轭相减，并使用 \(ad-bc=1\)，得到
\begin{equation}\label{b3:eq:23-imaginary-transformation}
 \operatorname{Im}(\gamma\tau)
 =\frac{\operatorname{Im}\tau}{|c\tau+d|^2},
 \quad \gamma'(\tau)=\frac1{(c\tau+d)^2}.
\end{equation}
前一式说明虚部仍为正，后一式说明它处处局部可逆；由整数逆矩阵给出的逆映射保证它是上半平面的全纯自同构。

现在计算换基与参数变换之间的对应。前章\cref{b3:prop:22-lattice-bases}证明，同一格的两组正向基恰由整数行列式一矩阵联系。为与上述分式的系数顺序一致，把新基写成
\begin{equation}\label{b3:eq:23-basis-change}
 \eta_1=d\omega_1+c\omega_2,
 \quad \eta_2=b\omega_1+a\omega_2.
\end{equation}
相应换基矩阵是 \(\left(\begin{smallmatrix}d&b\\c&a\end{smallmatrix}\right)\)，其行列式仍为一。于是
\[
 \frac{\eta_2}{\eta_1}
 =\frac{a\tau+b}{c\tau+d},
 \quad \eta_1=\omega_1(c\tau+d).
\]
后一个尺度因子与前一个参数变换同样重要：模形式的权将记录它怎样作用于格上的量。

\Needspace{6\baselineskip}
\begin{proposition}[参数轨道与比例格]\label{b3:prop:23-lattice-orbits}
设 \(\tau,\sigma\in\mathfrak H\)，\(\Lambda_\tau=\mathbb Z+\mathbb Z\tau\)、\(\Lambda_\sigma=\mathbb Z+\mathbb Z\sigma\)。存在 \(\lambda\in\mathbb C\setminus\{0\}\) 使 \(\Lambda_\sigma=\lambda\Lambda_\tau\)，当且仅当存在
\(\gamma=\left(\begin{smallmatrix}a&b\\c&d\end{smallmatrix}\right)\in\operatorname{SL}_2(\mathbb Z)\) 使 \(\sigma=\gamma\tau\)。对每个这样的 \(\gamma\)，都有
\begin{equation}\label{b3:eq:23-lattice-transformation}
 \Lambda_{\gamma\tau}=(c\tau+d)^{-1}\Lambda_\tau.
\end{equation}
\end{proposition}
\begin{proof}
整数对 \((d,c)\)、\((b,a)\) 构成行列式一的矩阵，所以 \(c\tau+d\)、\(a\tau+b\) 是 \(\Lambda_\tau\) 的一组基。将这组基同时除以 \(c\tau+d\)，所得格由 \(1\) 和 \(\gamma\tau\) 生成，证明了公式及充分性。

反过来，若 \(\Lambda_\sigma=\lambda\Lambda_\tau\)，则 \((1,\sigma)\) 与 \((\lambda,\lambda\tau)\) 都是该格的正向基。第二组的有向面积为 \(|\lambda|^2\operatorname{Im}\tau>0\)，所以这里不能出现行列式负一的换基。由\cref{b3:prop:22-lattice-bases}，有整数 \(a,b,c,d\) 满足 \(ad-bc=1\) 及
\[
 1=\lambda(d+c\tau),\quad
 \sigma=\lambda(b+a\tau).
\]
相除即得 \(\sigma=\gamma\tau\)。
\end{proof}

例如 \(T\tau=\tau+1\) 只把第二个生成元改为 \(\tau+1\)，格的点集本身没有改变。\(S\tau=-1/\tau\) 则在交换两条基方向后重新把第一条基归一到一，产生因子 \(\tau^{-1}\)。这里分类的是相差非零复比例的格；它与复环面双全纯分类的关系将在\secref{b3:sec:2306}证明。

\subsection{基本域}
\label{b3:subsec:230104}

为在每条参数轨道中选代表，先用 \(T\) 把实部限制在宽一的竖条带内。若所得点还在单位圆内，\(S\) 会增大它的虚部。由此得到候选区域
\[
 \mathcal F=\{\tau\in\mathfrak H:
           |\operatorname{Re}\tau|\le\tfrac12,\ |\tau|\ge1\}.
\]
它包含边界，因此边界上的代表可以重复。以下陈述把唯一性准确限制在内部。

\Needspace{6\baselineskip}
\begin{proposition}\label{b3:prop:23-fundamental-domain}
设 \(\operatorname{SL}_2(\mathbb Z)\) 以上述分式线性变换作用于 \(\mathfrak H\)，令 \(T(\tau)=\tau+1\)、\(S(\tau)=-1/\tau\)，并令
\[
 \mathcal F=\{\tau\in\mathfrak H:
       |\operatorname{Re}\tau|\le1/2,\ |\tau|\ge1\}.
\]
每条轨道与 \(\mathcal F\) 相交；\(\mathcal F\) 内部的不同点不在同一轨道。\(\mathcal F\) 中的边界识别由两类配对生成：两条竖边由 \(T\) 配对，单位圆弧的两半由 \(S\) 配对。
\end{proposition}
\begin{proof}
固定 \(\tau\in\mathfrak H\)。在全部互素整数对 \((c,d)\) 上，\(|c\tau+d|\) 有正的最小值。具体说，候选 \((0,1)\) 已给上界一；在 \(|c\tau+d|\le1\) 时，先由 \(|c|\operatorname{Im}\tau\le1\) 限制 \(c\)，再由 \(|c\operatorname{Re}\tau+d|\le1\) 限制 \(d\)，所以只有有限个候选，而且没有一个值为零。把取得最小值的底行补成群矩阵，由\eqref{b3:eq:23-imaginary-transformation}，其像具有整条轨道中的最大虚部。再作整数平移使实部进入 \([-1/2,1/2]\)。若此时模小于一，作用 \(S\) 会严格增大虚部，矛盾。因此所得点属于 \(\mathcal F\)。

为比较两个代表，先证明对 \(z\in\mathcal F\) 和任意互素整数对 \((c,d)\)，有
\[
 |cz+d|\ge1.
\]
若 \(c=0\)，互素性给 \(|d|=1\)。若 \(|c|\ge2\)，由 \(\operatorname{Im}z\ge\sqrt3/2\) 得 \(|cz+d|\ge\sqrt3>1\)。剩下 \(|c|=1\)，改变整体符号后只须看 \(|z+d|\)。当 \(d=0\) 时用 \(|z|\ge1\)；当 \(|d|\ge1\) 时，
\[
 |z+d|^2=|z|^2+2d\operatorname{Re}z+d^2
 \ge1-|d|+d^2\ge1.
\]
若 \(z\) 在内部，\(c\ne0\) 的各情形都严格大于一。

若 \(z,w\in\mathcal F\)、\(w=\gamma z\)，分别对 \(\gamma\) 和 \(\gamma^{-1}\) 使用这个估计，得到两点虚部相等及 \(|cz+d|=1\)。在内部只可能 \(c=0\)，此时变换为整数平移；两个实部都在 \((-1/2,1/2)\)，平移量只能为零。

在边界，\(c=0\) 的非零平移只能把一条竖边送到另一条。若 \(|c|=1,d=0\)，取等要求 \(|z|=1\)，相应变换是 \(S\) 后接整数平移；除端点外，这给圆弧两半的配对。\(|c|=|d|=1\) 还能取等的点只有实部为 \(\pm1/2\) 的圆弧端点；它们的识别已由这两类配对生成。其余情形均不取等，所以没有遗漏的边界识别。
\end{proof}

为把这幅基本域图变成商空间上的局部坐标，还需要说明远处的群元素不会无限次挤入同一个小邻域。这是\cref{b3:prop:11-discrete-halfplane-action}的紧集有限相交性质；在当前整数群中可以直接看到它的原因。

\Needspace{6\baselineskip}
\begin{proposition}[模群作用的紧集有限性]\label{b3:prop:23-proper-action}
设 \(\Gamma=\operatorname{PSL}_2(\mathbb Z)\) 作用于 \(\mathfrak H\)。对任意紧集 \(K,L\subset\mathfrak H\)，满足 \(\gamma K\cap L\ne\varnothing\) 的 \(\gamma\in\Gamma\) 只有有限个。特别地，每点的稳定子 \(\Gamma_p=\{\gamma\in\Gamma:\gamma p=p\}\) 都是有限群。
\end{proposition}
\begin{proof}
只须讨论两个紧集均非空的情形。若 \(\tau\in K\)、\(\gamma\tau\in L\)，两点的虚部都有统一正下界和有限上界。虚部变换公式先给 \(|c\tau+d|\) 的上界，随后由 \(\operatorname{Im}\tau\) 的正下界限制整数 \(c\)，再由 \(\operatorname{Re}\tau\) 的有界性限制整数 \(d\)。因此底行只有有限种可能。

固定一个互素底行 \((c,d)\) 和一个上排解 \((a_0,b_0)\) 后，全部上排解恰为 \((a_0+nc,b_0+nd)\)、\(n\in\mathbb Z\)。对应映射是固定映射 \(\gamma_0\) 后接平移 \(n\)。紧集 \(\gamma_0K\) 与 \(L\) 的实部均有界，故能使两者相交的 \(n\) 只有有限个。取 \(K=L=\{p\}\) 即得稳定子有限。
\end{proof}

令 \(\rho=e^{2\pi\mathrm{i}/3}=(-1+\sqrt3\,\mathrm{i})/2\)。基本域中真正有非平凡稳定子的轨道只有 \(\mathrm{i}\) 与 \(\rho\)。对于前者，固定点方程给 \(a=d\)、\(b=-c\)、\(a^2+c^2=1\)，在射影群中只剩 \(I,S\)。对于后者，利用 \(\rho^2+\rho+1=0\)，固定点方程给
\[
 b=-c,\quad d=a+c,\quad a^2+ac+c^2=1.
\]
等价地 \((2a+c)^2+3c^2=4\)，故只能有 \(c=0,\pm1\)；逐个代入得六个矩阵，即 \(\pm I,\pm ST,\pm(ST)^2\)。除去正负重复后，稳定子阶为三。基本域中其他点由上一证明的取等分类可排除非恒等稳定子：竖边上的平移不固定点，圆弧上 \(S\) 只有固定点 \(\mathrm{i}\)，额外的端点情形正是 \(\rho\) 的轨道。

\Needspace{6\baselineskip}
\begin{proposition}[固定点附近的商坐标]\label{b3:prop:23-stabilizer-chart}
设 \(\Gamma=\operatorname{PSL}_2(\mathbb Z)\)、\(p\in\mathfrak H\)，令 \(e=|\Gamma_p|\)。若 \(p\) 分别等价于 \(\mathrm{i}\)、\(\rho=e^{2\pi\mathrm{i}/3}\)，则 \(e\) 分别为二、三；其余点有 \(e=1\)。取
\[
 u_p(\tau)=\frac{\tau-p}{\tau-\overline p}.
\]
存在 \(r>0\)，使 \(U=\{\tau:|u_p(\tau)|<r\}\) 中两点模群等价，当且仅当它们由 \(\Gamma_p\) 的元素联系。这个局部商由
\[
 w=u_p(\tau)^e
\]
同胚到圆盘 \(|w|<r^e\)。在 \(U\) 上对 \(\Gamma_p\) 不变的全纯函数，恰能唯一写成 \(w\) 的全纯函数；亚纯函数有同样的下降性质。
\end{proposition}
\begin{proof}
稳定子阶数已在基本域中算出；其他代表的稳定子由共轭得到，阶数不变。映射 \(u_p\) 把上半平面双全纯映到单位圆盘，并把 \(p\) 送到零。每个稳定子元素在这个坐标中是固定零的圆盘自同构，因而由 Schwarz 引理及其逆向应用成为旋转。一个 \(e\) 阶有限旋转群恰由全部 \(e\) 次单位根组成：取其中最小正转角，反复相减说明所有转角都是它的整数倍，再用群的阶确定该转角为 \(2\pi/e\)。

取 \(p\) 附近一个紧闭小盘。由\cref{b3:prop:23-proper-action}，只有有限多个群元素能使它与自身相交。对于其中不固定 \(p\) 的元素，缩小 \(p\) 的邻域便能使该邻域与它的像不交；有限次缩小后，再取一个更小的 \(u_p\) 坐标圆盘 \(U\)。它被稳定子保持，而所有其他群元素的像均与它不交。因此 \(U\) 内的等价关系正是 \(u\sim\zeta u\)、\(\zeta^e=1\)。

映射 \(u\mapsto u^e\) 的纤维恰是这些旋转轨道，且它是开映射，所以诱导局部商与圆盘之间的同胚。若全纯函数的 Taylor 级数为 \(\sum_{n\ge0}a_nu^n\)，旋转不变性迫使 \(a_n=0\) 除非 \(e\mid n\)。保留的项组成关于 \(w=u^e\) 的收敛幂级数。亚纯情形在零点附近用 Laurent 级数作相同论证；其他位置局部可逆，故同样下降。
\end{proof}

在 \(p=\mathrm{i}\) 时，直接代入可得 \(u_p(S\tau)=-u_p(\tau)\)；在 \(p=\rho\) 时，\(ST\) 在 \(u_p\) 坐标中的旋转因子为 \(e^{-2\pi\mathrm{i}/3}\)。局部商坐标分别是平方和立方。对于不变亚纯函数，上半平面上的整数零点阶数与商坐标中的阶数因此可能不同；价公式中的 \(1/2\)、\(1/3\) 将准确记录这件事。

\subsection{尖点与基本域约化}
\label{b3:subsec:230105}

模群除了在内部识别点，还识别一些边界方向。扩充有理轴 \(\mathbb Q\cup\{\infty\}\) 中的点称为\term[jiandian]{尖点}[cusp]。全模群的尖点只有一条轨道：既约分数 \(a/c\) 可以补成行列式一矩阵的第一列，该矩阵把 \(\infty\) 送到 \(a/c\)；反过来，整数分式线性变换把 \(\infty\) 送到有理数或 \(\infty\)。这里说的是有理边界点，不能把整个实轴都加进来当作同一种局部端。

固定 \(\infty\) 的射影群元素恰为整数平移。确实，固定无穷远要求 \(c=0\)，于是 \(a=d=\pm1\)，分式成为 \(\tau\mapsto\tau+n\)。因此自然尝试用
\[
 q=e^{2\pi\mathrm{i}\tau}
\]
消去这个平移重复。要确认它也是完整模群商的局部坐标，还须排除其他群元素在高处额外识别点。

\Needspace{6\baselineskip}
\begin{proposition}[尖点附近的穿孔圆盘]\label{b3:prop:23-cusp-chart}
设 \(\Gamma=\operatorname{PSL}_2(\mathbb Z)\)、\(Y>1\)，令 \(\mathfrak H_Y=\{\tau:\operatorname{Im}\tau>Y\}\)。若 \(\tau,\sigma\in\mathfrak H_Y\) 模群等价，则 \(\sigma-\tau\in\mathbb Z\)。因此 \(\mathfrak H_Y\) 在轨道商 \(\mathfrak H/\Gamma\) 中的像由
\[
 q=e^{2\pi\mathrm{i}\tau}
\]
同胚到穿孔圆盘 \(0<|q|<e^{-2\pi Y}\)，局部反函数由复对数给出。
\end{proposition}
\begin{proof}
若 \(\sigma=\gamma\tau\) 且 \(c\ne0\)，写 \(y=\operatorname{Im}\tau>Y\)，则
\[
 \operatorname{Im}\sigma
 =\frac{y}{|c\tau+d|^2}
 \le\frac1{c^2y}<\frac1Y<Y,
\]
与 \(\sigma\in\mathfrak H_Y\) 矛盾。所以只能 \(c=0\)，即整数平移。指数映射在 \(\mathfrak H_Y\) 上的像正是所述穿孔圆盘，两个原像的差为整数；它局部双全纯且是开映射，故诱导同胚。选一个小圆盘上的对数分支，\(\tau=(2\pi\mathrm{i})^{-1}\log q\) 就给局部逆；更换分支只会把 \(\tau\) 改变一个整数。
\end{proof}

在这个穿孔圆盘中加入 \(q=0\)，便把唯一的高端补成一个点。\secref{b3:sec:2306}将把它与刚才的内部坐标拼成紧 Riemann 曲面。“在尖点全纯”所要求的是关于 \(q\) 的函数在整个零点邻域全纯，而不是只沿某一条竖直线有有限极限。

\ProseParagraphBreak
基本域的证明也给出一个可执行的约化过程：先作整数平移，把实部送进 \([-1/2,1/2]\)；若模小于一，再作 \(S\)；重复这两步。每次真正使用 \(S\) 都严格增大虚部，而且不会无限重复。若初始点为 \(\tau_0\)、虚部为 \(y_0\)，每个途中点 \(\gamma\tau_0\) 的虚部至少为 \(y_0\)，于是
\[
 |c\tau_0+d|^2
 =\frac{y_0}{\operatorname{Im}(\gamma\tau_0)}\le1.
\]
这样的整数底行只有有限多个，所以途中可能出现的高度只有有限多个。严格增高的步骤只能有限次，最后便落入 \(\mathcal F\)。

例如 \(\tau=(1+2\mathrm{i})/10\) 已在竖条带内，却位于单位圆内。计算得
\[
 S\tau=-2+4\mathrm{i},\quad T^2S\tau=4\mathrm{i},
 \quad T^2S=\begin{pmatrix}2&-1\\1&0\end{pmatrix}.
\]
因此它与 \(4\mathrm{i}\) 表示同一比例格。换基公式给
\(\Lambda_{4\mathrm{i}}=\tau^{-1}\Lambda_\tau\)，而 \(\tau^{-1}=2-4\mathrm{i}\)，所以右边由 \(2-4\mathrm{i}\) 和一生成，确实等于 \(\mathbb Z+4\mathrm{i}\mathbb Z\)。约化改变的是格的归一化描述。内部点具有唯一的基本域代表；遇到边界时，按已证明的配对识别即可。

% !TeX root = ../../Book3.tex
\section{Eisenstein 级数}
\label{b3:sec:2302}

前章的格不变量随基变换产生确定的幂次因子。绝对收敛是把这种格对称性转成函数变换律的条件；权二处会失去它。

\subsection{格和}
\label{b3:subsec:230201}

对整数 \(k\ge2\)，定义\term[eisenstein-jishu]{Eisenstein 级数}[Eisenstein series]
\[
G_{2k}(\tau)\coloneqq\sum_{(m,n)\in\mathbb Z^2\setminus\{(0,0)\}}(m\tau+n)^{-2k}.
\]
这就是格 \(\Lambda_\tau\) 的 \(2k\) 次倒数和。它使用全部非零格向量；若只对互素 \((m,n)\) 求和，结果须乘以 \(\sum_{r\ge1}r^{-2k}\)，不能混用两种归一化。

\subsection{收敛}
\label{b3:subsec:230202}

在任意紧集 \(K\subset\mathfrak H\) 上，实线性变换 \((m,n)\mapsto m\tau+n\) 的最小奇异值有共同正下界，所以 \(|m\tau+n|\ge c_K\sqrt{m^2+n^2}\)。前章格壳估计给正规绝对收敛，故 \(G_{2k}\) 全纯。求导后的项含因子 \(m\)，但总衰减仍至少为 \((m^2+n^2)^{-k}\)，同一紧集估计允许任意固定次求导。

\ProseParagraphBreak

若 \(k=1\)，壳估计只给调和级数 \(\sum j^{-1}\)，绝对收敛失败。规定先对 \(n\) 再对 \(m\) 求和会得到有用函数，但它的变换律多出一项，不能把下面的重排论证照搬过去。

\subsection{\texorpdfstring{\(G_{2k}\)}{G2k}}
\label{b3:subsec:230203}

\begin{proposition}\label{b3:prop:23-eisenstein-transformation}
若 \(k\ge2\)、\(\gamma=\left(\begin{smallmatrix}a&b\\c&d\end{smallmatrix}\right)\in\mathrm{SL}_2(\mathbb Z)\)，则
\(G_{2k}(\gamma\tau)=(c\tau+d)^{2k}G_{2k}(\tau)\)。
\end{proposition}
\begin{proof}
每项分母为 \((m(a\tau+b)+n(c\tau+d))/(c\tau+d)\)。变换 \((m,n)\mapsto(ma+nc,mb+nd)\) 是 \(\mathbb Z^2\) 的双射，且保持非零对。绝对收敛允许依此重排，提出公共因子即得结果。
\end{proof}
令 \(\sigma_r(N)=\sum_{d\mid N}d^r\)。从 \(\pi\cot\pi z=-\pi\mathrm{i}(1+2\sum_{r\ge1}e^{2\pi\mathrm{i}rz})\) 在上半平面的展开出发，求 \(2k-1\) 次导数得到
\[
\sum_{n\in\mathbb Z}(z+n)^{-2k}
=\frac{(2\pi\mathrm{i})^{2k}}{(2k-1)!}\sum_{r\ge1}r^{2k-1}e^{2\pi\mathrm{i}rz}.
\]
先分出 \(m=0\)，将正负 \(m\) 配对，再按 \(N=mr\) 收集系数，得
\[
G_{2k}=2\zeta(2k)+\frac{2(2\pi\mathrm{i})^{2k}}{(2k-1)!}
\sum_{N\ge1}\sigma_{2k-1}(N)q^N.
\]
这里 \(\zeta(2k)\) 只表示收敛实级数。双和的绝对收敛由 \(|q|<1\) 和多项式系数增长保证，因而合并合法。

\subsection{与 \texorpdfstring{\(\wp\)}{\AnalysisPdfWeierstrassP} 的关系}
\label{b3:subsec:230204}

前章的记号现在为 \(g_2(\Lambda_\tau)=60G_4(\tau)\)、\(g_3(\Lambda_\tau)=140G_6(\tau)\)。格缩放律已经预示它们的权分别为四、六。相比之下，\(\wp(z;\Lambda_\tau)\) 还带一个曲线上的变量 \(z\)；改变 \(\tau\) 的基时必须同时缩放 \(z\)，不能将它误当成只依赖模参数的模形式。

% !TeX root = ../../Book3.tex
\section{模形式}
\label{b3:sec:2303}

Eisenstein 级数已经给出一种特殊的全纯函数：改变格基时，它不完全不变，而是乘上一个已知的幂次因子。同时，它在 \(\operatorname{Im}\tau\to\infty\) 的退化方向有有限的常数项。本节把这两项性质提炼为定义，并说明每一项的解析含义。尤其要分清三个层次：整数平移给周期性，整个模群给带权变换律，尖点条件则控制周期函数在无穷远的增长。

\subsection{模形式}
\label{b3:subsec:230301}

由\cref{b3:prop:23-eisenstein-transformation}，\(G_{2r}\) 在变换 \(\tau\mapsto(a\tau+b)/(c\tau+d)\) 下乘以 \((c\tau+d)^{2r}\)。这个因子来自格的尺度变化。若强求每个格和都在换基后数值相同，就会把已经构造出的有用对象排除。因而定义容许一个固定幂次的变化，但要求它对每个群元素都使用同一种规则。

另一方面，变换 \(T\tau=\tau+1\) 的因子为一，故所有这种函数都应当周期一。周期性使 \(f(\tau)\) 能用 \(q=e^{2\pi\mathrm{i}\tau}\) 表示：两个 \(\tau\) 给出同一个 \(q\)，当且仅当它们相差整数。由局部对数构成的函数先定义在穿孔圆盘上；要求它在 \(q=0\) 全纯，是一项额外条件。以下定义将这一要求直接写成一个圆盘全纯函数的存在性，避免把“尖点处全纯”留作未说明的边界用语。

\Needspace{6\baselineskip}
\begin{definition}\label{b3:def:23-modular-form}
设 \(k\in\mathbb Z\)，\(\mathfrak H=\{\tau\in\mathbb C:\operatorname{Im}\tau>0\}\)，\(\mathbb D=\{q\in\mathbb C:|q|<1\}\)。一个全纯函数 \(f:\mathfrak H\to\mathbb C\) 称为全模群 \(\operatorname{SL}_2(\mathbb Z)\) 上权 \(k\) 的\term[mo-xingshi]{模形式}[modular form]，是指它满足：
\begin{enumerate}
\item\label{b3:item:23-modular-transformation}对每个 \(\left(\begin{smallmatrix}a&b\\c&d\end{smallmatrix}\right)\in\operatorname{SL}_2(\mathbb Z)\) 和每个 \(\tau\in\mathfrak H\)，都有
\[
 f\!\left(\frac{a\tau+b}{c\tau+d}\right)
 =(c\tau+d)^k f(\tau);
\]
\item\label{b3:item:23-modular-cusp}存在全纯函数 \(F:\mathbb D\to\mathbb C\)，使
\[
 f(\tau)=F(e^{2\pi\mathrm{i}\tau})
 \quad(\tau\in\mathfrak H).
\]
\end{enumerate}
第二项称为在尖点 \(\infty\) 全纯。所有这样的函数构成复向量空间，记为 \(M_k\) 或 \(M_k(\operatorname{SL}_2(\mathbb Z))\)。本章不另加特征因子。
\end{definition}
\symidx{\(M_k\)：全模群上权 \(k\) 的模形式空间}

这里允许任意整数 \(k\) 是为了统一陈述。哪些权实际有非零形式，要由变换律和后面的零点计数判断；不应在定义时就把这些结论暗中假定。由于 \(c\tau+d\ne0\)，负整数幂也有单值的通常意义，没有复幂分支问题。

定义中的 \(F\) 必唯一，因为指数映射的像包含整个穿孔圆盘，两个可能的 \(F\) 在那里相等，故在零点也相等。空间 \(M_k\) 对复数线性组合封闭：变换律是线性的，圆盘上的全纯延拓也可作相同线性组合。前节的 Fourier 展开已经证明，对每个整数 \(r\ge2\)，都有 \(G_{2r}\in M_{2r}\)；它们不是靠定义猜出的例子，而是绝对收敛格和产生的具体成员。

\subsection{权}
\label{b3:subsec:230302}

定义中的整数 \(k\) 称为\term[mo-xingshi-quan]{权}[weight]。要理解它为何能与群乘法相容，记
\[
 J(\gamma,\tau)=c\tau+d,
 \quad \gamma=\begin{pmatrix}a&b\\c&d\end{pmatrix}.
\]
直接相乘两个矩阵给出
\begin{equation}\label{b3:eq:23-factor-cocycle}
 J(\gamma_1\gamma_2,\tau)
 =J(\gamma_1,\gamma_2\tau)J(\gamma_2,\tau).
\end{equation}
例如，若 \(\gamma_1\) 的底行为 \((c_1,d_1)\)，\(\gamma_2\) 的两行为 \((a_2,b_2)\)、\((c_2,d_2)\)，右边等于
\[
 c_1(a_2\tau+b_2)+d_1(c_2\tau+d_2),
\]
正是乘积矩阵的底行给出的线性式。因此依次作两次变换所积累的权因子，与一次作复合变换的因子一致。

为了把这个运算写得紧凑，对 \(\mathfrak H\) 上的全纯或亚纯函数 \(f\)、整数 \(k\) 和 \(\gamma\in\operatorname{SL}_2(\mathbb Z)\)，记
\[
 (f|_k\gamma)(\tau)
 =J(\gamma,\tau)^{-k}f(\gamma\tau).
\]
这个记号在变换后除去权因子，使带权变换律可以写成不变性条件。公式\eqref{b3:eq:23-factor-cocycle}给
\[
 (f|_k\gamma_1)|_k\gamma_2=f|_k(\gamma_1\gamma_2).
\]
所以变换律正是 \(f|_k\gamma=f\)。它也解释了为什么只核对生成元就足够：满足对 \(S,T\) 的不变性，便依次满足它们的逆和全部有限乘积的不变性。

\Needspace{6\baselineskip}
\begin{proposition}[生成元判别与权的运算]\label{b3:prop:23-weight-operations}
设 \(k\in\mathbb Z\)，\(f\) 为 \(\mathfrak H\) 上全纯函数。它满足全模群的权 \(k\) 变换律，当且仅当
\begin{equation}\label{b3:eq:23-generators-form-law}
 f(\tau+1)=f(\tau),\quad
 f(-1/\tau)=\tau^k f(\tau)
 \quad(\tau\in\mathfrak H).
\end{equation}
奇数权的 \(M_k\) 仅含零函数。若 \(k,\ell\in\mathbb Z\)、\(f\in M_k\)、\(g\in M_\ell\)，则 \(fg\in M_{k+\ell}\)；若 \(g\not\equiv0\)，商 \(f/g\) 在上半平面及尖点亚纯，并满足权 \(k-\ell\) 的变换律。
\end{proposition}
\begin{proof}
必要性是代入 \(S,T\)。充分性由上述带权运算的复合公式及 \(S,T\) 的生成性得到。对矩阵 \(-I\)，变换律给 \(f(\tau)=(-1)^kf(\tau)\)，所以 \(k\) 为奇数时只能 \(f\equiv0\)。

两个变换律相乘便使权相加；圆盘全纯表达式的乘积仍全纯，故乘积在尖点也满足要求。商的变换律同样由相除得到。非恒零全纯函数的零点孤立，局部写成一个整数次幂乘无零全纯函数，所以商在有限点亚纯；在尖点对两个圆盘 Taylor 级数作同样的局部分解，便知商至多有有限阶极点。
\end{proof}

同权形式可以相加，不同权的乘积则进入新的权空间。例如 \(G_4G_6\) 的权为十。若要同时保留所有权，就要把它们组成带权标记的分次空间；不能把不同权相加后仍称为原来某个固定权的形式。同权商的因子相消，得到不变的亚纯函数，这将在\secref{b3:sec:2306}成为模函数的来源。

\ProseParagraphBreak
权的另一个解释来自未归一化的格。对格不变量 \(G_{2r}(\Lambda)=\sum_{\omega\in\Lambda\setminus\{0\}}\omega^{-2r}\)，同时把全部格向量乘以 \(\lambda\) 会使它乘以 \(\lambda^{-2r}\)。下述命题说明一般变换律也是这样一种尺度齐次性的坐标表达。

\Needspace{6\baselineskip}
\begin{proposition}[权与格的尺度]\label{b3:prop:23-lattice-weight}
设 \(k\in\mathbb Z\)，\(f:\mathfrak H\to\mathbb C\) 满足全模群的权 \(k\) 变换律。对复平面中的格 \(\Lambda\)，选取任一正向基 \((\omega_1,\omega_2)\)，定义
\[
 \Phi_f(\Lambda)=\omega_1^{-k}
          f\!\left(\frac{\omega_2}{\omega_1}\right).
\]
此值不依赖正向基的选择，且对每个 \(\lambda\in\mathbb C\setminus\{0\}\) 有
\[
 \Phi_f(\lambda\Lambda)=\lambda^{-k}\Phi_f(\Lambda).
\]
反过来，若一个复值格函数 \(\Phi\) 满足上述尺度律，且 \(f(\tau)=\Phi(\mathbb Z+\mathbb Z\tau)\) 在 \(\mathfrak H\) 全纯，则 \(f\) 满足权 \(k\) 的变换律；它是否是模形式，还须检查尖点全纯性。
\end{proposition}
\begin{proof}
若另一组正向基为 \((\eta_1,\eta_2)\)，由\eqref{b3:eq:23-basis-change}，存在 \(\gamma\in\operatorname{SL}_2(\mathbb Z)\) 使
\[
 \eta_1=\omega_1(c\tau+d),\quad
 \frac{\eta_2}{\eta_1}=\gamma\tau,
 \quad \tau=\frac{\omega_2}{\omega_1}.
\]
因此
\[
 \eta_1^{-k}f(\eta_2/\eta_1)
 =\omega_1^{-k}(c\tau+d)^{-k}f(\gamma\tau)
 =\omega_1^{-k}f(\tau).
\]
乘以非零复数 \(\lambda\) 后，\((\lambda\omega_1,\lambda\omega_2)\) 仍正向，基的比值不变而前因子乘以 \(\lambda^{-k}\)，得到尺度律。

反过来，由\eqref{b3:eq:23-lattice-transformation}及尺度律，
\[
 f(\gamma\tau)
 =\Phi((c\tau+d)^{-1}\Lambda_\tau)
 =(c\tau+d)^k\Phi(\Lambda_\tau)
 =(c\tau+d)^kf(\tau).
\]
这个计算没有使用 \(\operatorname{Im}\tau\to\infty\) 时的任何估计，所以不会自动给出尖点条件。
\end{proof}

格的点集满足 \(-\Lambda=\Lambda\)，所以奇数齐次权也会迫使 \(\Phi_f=0\)，与矩阵 \(-I\) 的判别相同。另一方面，权并不表示 \(f\) 作为变量 \(\tau\) 的普通齐次次数。缩放格与把参数 \(\tau\) 乘一个复数是两种不同操作。

普通导数也不会自动把权提高二。设 \(k\in\mathbb Z\)，\(f\) 在 \(\mathfrak H\) 全纯并满足全模群权 \(k\) 的变换律。对这一变换律求导，使用 \((\gamma\tau)'=(c\tau+d)^{-2}\)，得到
\begin{equation}\label{b3:eq:23-ordinary-derivative-law}
 f'(\gamma\tau)
 =(c\tau+d)^{k+2}f'(\tau)
   +kc(c\tau+d)^{k+1}f(\tau).
\end{equation}
额外的一项来自权因子的导数。\subsecref{b3:subsec:230503}将构造修正项使它相消；这里先把普通导数的真实变换律保留下来。

\subsection{尖点模形式}
\label{b3:subsec:230303}

若模形式的尖点常数项为零，它在退化端比一般模形式衰减得更快。这个性质对线性组合稳定，因而给出一个自然子空间。

\Needspace{6\baselineskip}
\begin{definition}\label{b3:def:23-cusp-form}
设 \(k\in\mathbb Z\)，\(f\in M_k(\operatorname{SL}_2(\mathbb Z))\)，并令 \(F\) 为单位圆盘上的唯一全纯函数，满足 \(f(\tau)=F(e^{2\pi\mathrm{i}\tau})\)。若 \(F(0)=0\)，则称 \(f\) 为全模群上权 \(k\) 的\term[jiandian-mo-xingshi]{尖点模形式}[cusp form]。这类函数构成的复向量空间记为 \(S_k\) 或 \(S_k(\operatorname{SL}_2(\mathbb Z))\)。
\end{definition}
\symidx{\(S_k\)：全模群上权 \(k\) 的尖点模形式空间}

常数项映射
\[
 a_0:M_k\longrightarrow\mathbb C,\quad f\longmapsto F(0)
\]
是线性的，其核恰为 \(S_k\)。因此 \(M_k/S_k\) 的维数至多为一。若已找到常数项非零的权 \(k\) 形式，将它除以常数项就得到 \(a_0=1\) 的形式，从任何同权 \(f\) 减去适当倍数后便进入 \(S_k\)。这一简单操作将成为后面模形式环归纳证明的第一步。

尖点模形式乘以任何全纯模形式仍在尖点消失，所以 \(S_kM_\ell\subset S_{k+\ell}\)。不过，单独的周期性和衰减还不够：\(f(\tau)=e^{2\pi\mathrm{i}\tau}\) 周期一且在高处趋零，但它不是权零模形式。例如
\[
 f(2\mathrm{i})=e^{-4\pi},\quad
 f(S(2\mathrm{i}))=f(\mathrm{i}/2)=e^{-\pi},
\]
不满足权零的 \(S\) 不变性。三个条件必须一起核对，不能由某一端的良好行为代替群变换律。

这里检查一个尖点就够，是因为\subsecref{b3:subsec:230105}已证明全模群的有理尖点只有一条轨道。若用某个 \(\gamma\) 把 \(\infty\) 送到另一尖点，应先用 \(f|_k\gamma\) 表示该处函数；全模群的变换律使它仍等于 \(f\)，所以得到同一尖点条件。对有限指标子群可能出现多条尖点轨道，不能把本章的简化照搬过去。

\subsection{Fourier 展开}
\label{b3:subsec:230304}

现在证明前面使用的 \(q\) 表达式，并把尖点条件变成可以直接检查的增长条件。先暂时不要求 \(S\) 变换律，只保留周期性，这样就能看清 Fourier 展开究竟依赖哪一步。

\Needspace{6\baselineskip}
\begin{proposition}[周期函数的尖点展开]\label{b3:prop:23-q-expansion}
设 \(f:\mathfrak H\to\mathbb C\) 全纯，且 \(f(\tau+1)=f(\tau)\)。则存在唯一全纯函数
\(F:\{q:0<|q|<1\}\to\mathbb C\)，满足 \(f(\tau)=F(e^{2\pi\mathrm{i}\tau})\)。其 Laurent 展开给
\[
 f(\tau)=\sum_{n\in\mathbb Z}a_ne^{2\pi\mathrm{i}n\tau},
\]
在 \(\mathfrak H\) 的紧集上局部一致收敛，并且对任意整数 \(n\) 和任意 \(y>0\)，
\begin{equation}\label{b3:eq:23-fourier-coefficients}
 a_n=e^{2\pi ny}\int_0^1
        f(x+\mathrm{i}y)e^{-2\pi\mathrm{i}nx}\,\mathrm{d}x.
\end{equation}
以下三项等价：\(F\) 在 \(q=0\) 可去；存在 \(Y>0,C<\infty\) 使 \(|f(x+\mathrm{i}y)|\le C\) 对所有 \(x\in\mathbb R,y>Y\) 成立；所有 \(n<0\) 的 \(a_n\) 都为零。条件成立时，级数为 \(\sum_{n\ge0}a_nq^n\)，在整个 \(|q|<1\) 内收敛，且
\[
 \sup_{x\in\mathbb R}|f(x+\mathrm{i}y)-a_0|
 =O(e^{-2\pi y})\quad(y\to\infty).
\]
\end{proposition}
\begin{proof}
给定 \(0<|q|<1\)，任选一个满足 \(q=e^{2\pi\mathrm{i}\tau}\) 的 \(\tau\in\mathfrak H\)，定义 \(F(q)=f(\tau)\)。另一个选择只会将 \(\tau\) 改变整数，周期性保证值不变。在每个不含零的小圆盘内可选取全纯对数，故 \(F(q)=f((2\pi\mathrm{i})^{-1}\log q)\) 局部全纯。这既证明存在性，也由指数映射的满射性证明唯一性。

Laurent 定理在环域 \(0<|q|<1\) 给唯一展开 \(F(q)=\sum_{n\in\mathbb Z}a_nq^n\)，在每个紧子环域上一致收敛。\(\mathfrak H\) 中紧集的指数像是这种环域中的紧集，所以拉回也局部一致。系数公式为
\[
 a_n=\frac1{2\pi\mathrm{i}}
        \int_{|q|=r}F(q)q^{-n-1}\,\mathrm{d}q,
 \quad 0<r<1.
\]
取 \(r=e^{-2\pi y}\)，按 \(q=e^{2\pi\mathrm{i}(x+\mathrm{i}y)}\)、\(0\le x\le1\) 参数化这条正向圆周，便得到\eqref{b3:eq:23-fourier-coefficients}。这个推导也说明结果与 \(y\) 无关：不同高度只是计算同一个 Laurent 系数的不同积分圆周。

若 \(F\) 在零点可去，它在某个闭小圆盘上有界；因为 \(|q|=e^{-2\pi y}\)，对应地 \(f\) 在整个高半平面有界。反过来，高半平面的一致界转成 \(F\) 在某个穿孔小圆盘上的界，由可去奇点定理得到全纯延拓。Laurent 定理又说明，可去性恰等价于全部负次系数为零。

延拓后，\(F(q)=a_0+qB(q)\)，其中 \(B\) 在单位圆盘全纯。在一个固定闭小圆盘上 \(|B|\) 有界，故 \(|F(q)-a_0|\le C|q|\)。代入 \(q=e^{2\pi\mathrm{i}(x+\mathrm{i}y)}\) 就得到最后的一致估计。
\end{proof}

因此，对一个已经满足全模群变换律的全纯函数，验证尖点条件可以改为在 \(|\operatorname{Re}\tau|\le1/2\)、\(\operatorname{Im}\tau>Y\) 上取得统一界；周期性会把这个界延伸到整个高半平面。只沿一条竖线检查极限并不具有同样作用。例如
\[
 f(\tau)=\exp(-e^{-2\pi\mathrm{i}\tau})
\]
全纯且周期一，它对应 \(F(q)=e^{-1/q}\)。沿 \(\tau=\mathrm{i}y\) 有 \(f(\mathrm{i}y)=\exp(-e^{2\pi y})\to0\)，但沿 \(\tau=1/2+\mathrm{i}y\) 则有 \(f(\tau)=\exp(e^{2\pi y})\to\infty\)。一个方向的有限极限不能排除零点处的本性奇点。这个例子只用于检验周期函数的尖点行为，不声称它满足整个模群的变换律。

对模形式，\eqref{b3:eq:23-fourier-coefficients}就是 Fourier 系数的复分析来源。级数在每个固定高度上按 \(x\) 展开成周期函数，又在 \(q\) 圆盘中是一条 Taylor 级数。这里没有调用 \(L^2\) Fourier 级数的完备性：全纯性已经给出比一般实变量周期函数更强的局部一致收敛。特别地，模形式由全部系数唯一确定，而尖点模形式恰是 \(a_0=0\) 的那些形式。

\subsection{在尖点处消失}
\label{b3:subsec:230305}

尖点全纯只说明没有负次幂。若进一步需要计数零点，就应记录第一个非零幂次。设 \(f\in M_k\) 非恒零，写
\[
 f(\tau)=\sum_{n\ge0}a_nq^n,
 \quad q=e^{2\pi\mathrm{i}\tau}.
\]
其尖点消失阶定义为
\[
 v_\infty(f)=\min\{n\ge0:a_n\ne0\}.
\]
非恒零保证这个最小值存在；若所有系数为零，Taylor 唯一性会使 \(f\equiv0\)。可约定零函数的阶数为 \(+\infty\)，但后面涉及对数导数的公式只用于非恒零函数。尖点形式满足 \(v_\infty(f)\ge1\)，而一般非零模形式允许阶数为零。

价公式还要处理在尖点有极点的函数，所以将同一计算保留到 Laurent 展开具有有限主部的情形。下面的一致估计是围道顶边取极限时的直接依据。

\Needspace{6\baselineskip}
\begin{proposition}[尖点阶数与对数导数]\label{b3:prop:23-cusp-order-estimate}
设 \(Y_0\in\mathbb R\)，\(f\) 是半平面 \(\operatorname{Im}\tau>Y_0\) 上非恒零的周期一亚纯函数，其关于 \(q=e^{2\pi\mathrm{i}\tau}\) 的函数在 \(q=0\) 亚纯。设该 Laurent 展开的首项为 \(a_mq^m\)，其中 \(m\in\mathbb Z\)、\(a_m\ne0\)。则存在 \(Y>Y_0\) 和常数 \(C>0\)，使 \(f\) 在 \(\operatorname{Im}\tau>Y\) 没有零点及极点，而且对所有 \(x\in\mathbb R,y>Y\)，都有
\begin{align}
 \left|\frac{f(x+\mathrm{i}y)}{a_m e^{2\pi\mathrm{i}m(x+\mathrm{i}y)}}-1\right|
 &\le Ce^{-2\pi y},\label{b3:eq:23-cusp-leading-term}\\
 \left|\frac{f'(x+\mathrm{i}y)}{f(x+\mathrm{i}y)}-2\pi\mathrm{i}m\right|
 &\le Ce^{-2\pi y}.\label{b3:eq:23-cusp-log-derivative}
\end{align}
这里 \(m=v_\infty(f)\)，极点的阶数按负整数计。
\end{proposition}
\begin{proof}
尖点亚纯性给出分解 \(F(q)=q^mA(q)\)，其中 \(A\) 在零点某邻域全纯，且 \(A(0)=a_m\ne0\)。取足够小的 \(r>0\)，使 \(A\) 在 \(|q|\le r\) 无零点；再略缩小圆盘，\(A'\) 与 \(1/A\) 都有统一界。因此
\[
 A(q)=a_m+O(|q|),\quad
 \frac{A'(q)}{A(q)}=O(1)
\]
对整个小圆盘一致成立。令 \(e^{-2\pi Y}<r\)，\(q\ne0\) 及 \(A(q)\ne0\) 说明 \(f\) 在所述高半平面无零极点。第一式来自 \(F(q)/(a_mq^m)=A(q)/a_m\)。

对第二式，利用 \(q'(\tau)=2\pi\mathrm{i}q\) 直接计算
\[
 \frac{f'(\tau)}{f(\tau)}
 =2\pi\mathrm{i}\left(m+q\frac{A'(q)}{A(q)}\right).
\]
余项绝对值不超过常数乘 \(|q|=e^{-2\pi y}\)，与辐角 \(2\pi x\) 无关，因而估计对全部 \(x\) 一致。
\end{proof}

这比“每条竖线上趋于某值”更强：它同时控制整条水平线，因此可沿长度为一的顶边积分并让高度趋于无穷。对 \(f\in S_k\setminus\{0\}\)，还得到 \(|f(x+\mathrm{i}y)|\le C e^{-2\pi m y}\)、\(m\ge1\)，说明尖点处的消失在上半平面中表现为指数衰减。

由局部分解立即可得，在各函数均非恒零且在尖点亚纯时，
\[
 v_\infty(fg)=v_\infty(f)+v_\infty(g),\quad
 v_\infty(f/g)=v_\infty(f)-v_\infty(g).
\]
若 \(f+g\not\equiv0\)，则 \(v_\infty(f+g)\ge\min(v_\infty(f),v_\infty(g))\)；两阶相等时首项可能相消而使不等式严格。后面从 \(E_4^3\) 与 \(E_6^2\) 相减构造判别式，正是有意制造这样的相消。

最后，\(q=0\) 的阶数是尖点坐标中的阶数，不是把 \(\tau=\infty\) 当作普通有限点后读出的次数。有限的椭圆固定点另有 \(u^2\)、\(u^3\) 商坐标，它们与这里的指数坐标分工不同。\secref{b3:sec:2305}将把这些局部阶数放入同一条零点计数公式，而每一项所用的坐标都会保持明确。

% !TeX root = ../../Book3.tex
\section{\texorpdfstring{\(E_4,E_6\)}{E₄, E₆} 与判别式}
\label{b3:sec:2304}

已经知道 \(G_4,G_6\) 是模形式，但它们的常数项分别含有 \(\zeta(4),\zeta(6)\)。把常数项规范为一，可以更直接地比较系数。随后，\(E_4^3\) 与 \(E_6^2\) 具有相同的权和常数项，二者相减便产生一个尖点模形式。这个构造还将恢复前章三次曲线的判别式，因而能把尖点处的消失与上半平面内部的非退化性联系起来。

\subsection{归一化 Eisenstein 级数}
\label{b3:subsec:230401}

\begin{definition}[归一化 Eisenstein 级数]\label{b3:def:23-normalized-eisenstein}
设 \(k\ge2\) 为整数，\(\tau\in\mathfrak H\)，且
\[
 G_{2k}(\tau)=\sum_{(m,n)\in\mathbb Z^2\setminus\{(0,0)\}}(m\tau+n)^{-2k},
 \quad \zeta(2k)=\sum_{n=1}^{\infty}n^{-2k}.
\]
定义归一化的 Eisenstein 级数
\[
 E_{2k}(\tau)\coloneqq\frac{G_{2k}(\tau)}{2\zeta(2k)}.
\]
\end{definition}
\symidx{\(E_{2k}(\tau)\)：常数项为一的 Eisenstein 级数}

分母是正实数，所以这个定义不会除以零。\cref{b3:prop:23-eisenstein-transformation}与\cref{b3:prop:23-eisenstein-fourier}说明 \(E_{2k}\in M_{2k}\)，并给出
\begin{equation}\label{b3:eq:23-normalized-eisenstein-fourier}
 E_{2k}(\tau)=1+
 \frac{(2\pi\mathrm{i})^{2k}}{(2k-1)!\,\zeta(2k)}
 \sum_{n=1}^{\infty}\sigma_{2k-1}(n)q^n,
 \quad q=e^{2\pi\mathrm{i}\tau}.
\end{equation}
它的常数项等于一，故不是尖点模形式。要把权四、六的系数写成具体整数，还需要算出相应的偶数幂和。这个计算只用本卷已建立的余切展开。

\begin{proposition}[三个偶数幂和的值]\label{b3:prop:23-zeta-even-values}
对实数 \(s>1\)，令 \(\zeta(s)=\sum_{n=1}^{\infty}n^{-s}\)。则
\[
 \zeta(2)=\frac{\pi^2}{6},\quad
 \zeta(4)=\frac{\pi^4}{90},\quad
 \zeta(6)=\frac{\pi^6}{945}.
\]
\end{proposition}
\begin{proof}
由\cref{b3:prop:20-sine-product}，在 \(0<|z|<1\) 内，
\[
 \pi\cot\pi z
 =\frac1z+\sum_{n=1}^{\infty}\frac{2z}{z^2-n^2}.
\]
对每个正整数 \(n\) 展开几何级数，
\[
 \frac{2z}{z^2-n^2}
 =-2\sum_{r=1}^{\infty}\frac{z^{2r-1}}{n^{2r}}.
\]
若 \(|z|\le\rho<1\)，右侧双和的绝对值之和至多为
\[
 2\sum_{n=1}^{\infty}\frac{\rho/n^2}{1-\rho^2/n^2}
 \le\frac{2\rho}{1-\rho^2}\sum_{n=1}^{\infty}n^{-2}<\infty.
\]
所以可以交换两次求和，得到
\begin{equation}\label{b3:eq:23-cotangent-zeta}
 \pi\cot\pi z=\frac1z-2\sum_{r=1}^{\infty}\zeta(2r)z^{2r-1},
 \quad 0<|z|<1.
\end{equation}

现在用正弦、余弦的 Taylor 级数独立计算左端。置 \(t=\pi z\)，并写
\[
 t\cot t=1+At^2+Bt^4+Ct^6+O(t^8).
\]
因为 \(t\cot t=\cos t/(\sin t/t)\)，所以
\[
 \begin{aligned}
 &\left(1-\frac{t^2}{6}+\frac{t^4}{120}
              -\frac{t^6}{5040}+O(t^8)\right)
   \left(1+At^2+Bt^4+Ct^6+O(t^8)\right)\\
 &\hspace{2em}=1-\frac{t^2}{2}+\frac{t^4}{24}
                    -\frac{t^6}{720}+O(t^8).
 \end{aligned}
\]
逐次比较二、四、六次项，得到
\[
 A-\frac16=-\frac12,\quad
 B-\frac A6+\frac1{120}=\frac1{24},\quad
 C-\frac B6+\frac A{120}-\frac1{5040}=-\frac1{720}.
\]
第一式给 \(A=-1/3\)；代入第二式得 \(B=-1/45\)，再代入第三式得 \(C=-2/945\)。换回 \(z\)，
\[
 \pi\cot\pi z=\frac1z-\frac{\pi^2}{3}z
                    -\frac{\pi^4}{45}z^3
                    -\frac{2\pi^6}{945}z^5+O(z^7).
\]
与\eqref{b3:eq:23-cotangent-zeta}比较对应系数，就得到三个数值。
\end{proof}

这里先由正规收敛的部分分式得到幂和系数，再由初等 Taylor 除法算出同一函数的系数。整个证明不需要关于 \(\zeta(s)\) 的函数方程；后面学习解析数论时，可以把这几个数值作为已知的正实轴数据使用。

\subsection{\texorpdfstring{\(E_4,E_6\)}{E4,E6}}
\label{b3:subsec:230402}

\begin{proposition}[权四与权六的 Eisenstein 级数]\label{b3:prop:23-e4-e6}
设 \(\tau\in\mathfrak H\)，\(q=e^{2\pi\mathrm{i}\tau}\)，并按\cref{b3:def:23-normalized-eisenstein}定义 \(E_4,E_6\)。则 \(E_4\in M_4\)、\(E_6\in M_6\)，且
\begin{equation}\label{b3:eq:23-e4-e6-fourier}
 \begin{aligned}
 E_4(\tau)&=1+240\sum_{n=1}^{\infty}\sigma_3(n)q^n,\\
 E_6(\tau)&=1-504\sum_{n=1}^{\infty}\sigma_5(n)q^n.
 \end{aligned}
\end{equation}
这些幂级数在 \(|q|<1\) 内局部正规收敛，并有
\begin{equation}\label{b3:eq:23-e4-e6-initial}
 \begin{aligned}
 E_4&=1+240q+2160q^2+6720q^3+O(q^4),\\
 E_6&=1-504q-16632q^2-122976q^3+O(q^4).
 \end{aligned}
\end{equation}
特别地，二者都非恒零，且尖点消失阶均为零。
\end{proposition}
\begin{proof}
由\eqref{b3:eq:23-normalized-eisenstein-fourier}，权四的系数因子为
\[
 \frac{(2\pi\mathrm{i})^4}{3!\,\zeta(4)}
 =\frac{16\pi^4}{6\cdot(\pi^4/90)}=240,
\]
而权六的系数因子为
\[
 \frac{(2\pi\mathrm{i})^6}{5!\,\zeta(6)}
 =\frac{-64\pi^6}{120\cdot(\pi^6/945)}=-504.
\]
负号来自 \(\mathrm{i}^6=-1\)。一般 Fourier 展开已经保证了变换律与尖点全纯性，因此它们确实属于所述模形式空间。

前三个正整数的约数给出
\[
 \begin{aligned}
 &\sigma_3(1)=1,\quad \sigma_3(2)=1+2^3=9,\quad
   \sigma_3(3)=1+3^3=28,\\
 &\sigma_5(1)=1,\quad \sigma_5(2)=1+2^5=33,\quad
   \sigma_5(3)=1+3^5=244.
 \end{aligned}
\]
代入即可得到首项展开。常数项均为一，故二者不是零函数，也不在尖点消失。
\end{proof}

这些整数系数来自约数和及刚刚计算的常数。一般模形式可以乘任意复常数，所以“是模形式”本身没有强迫系数为整数。\(\tau\) 的函数恒等式也不能仅凭几个初项相同就判定；\secref{b3:sec:2305}的价公式将给出可以比较有限多个系数的精确依据。

\subsection{判别式模形式}
\label{b3:subsec:230403}

寻找尖点模形式时，可以先让常数项相消。但相减的两个模形式还必须具有相同的权，否则变换因子无法共同提出。权四与权六的最小正公倍数为十二，因此首先比较 \(E_4^3\) 与 \(E_6^2\)。二者常数项都是一，一次项却分别是 \(720q\) 与 \(-1008q\)，所以差的首项为 \(1728q\)。除以 \(1728\) 就使第一个非零系数也规范为一。

\begin{definition}[判别式模形式]\label{b3:def:23-delta}
设 \(\tau\in\mathfrak H\)，\(E_4,E_6\) 为\cref{b3:def:23-normalized-eisenstein}中常数项为一的权四、权六 Eisenstein 级数。函数
\[
 \Delta(\tau)\coloneqq
 \frac{E_4(\tau)^3-E_6(\tau)^2}{1728}
\]
称为\term[panbieshi-mo-xingshi]{判别式模形式}[discriminant modular form]。
\end{definition}
\symidx{\(\Delta(\tau)\)：判别式模形式}

\begin{proposition}[判别式的权与尖点阶数]\label{b3:prop:23-delta-cusp}
对上半平面上的归一化 Eisenstein 级数 \(E_4,E_6\)，令 \(\Delta=(E_4^3-E_6^2)/1728\)。则 \(\Delta\) 是全模群的非零权十二尖点模形式，并且
\begin{equation}\label{b3:eq:23-delta-fourier}
 \Delta(\tau)=q-24q^2+252q^3+O(q^4),
 \quad q=e^{2\pi\mathrm{i}\tau}.
\end{equation}
因此 \(v_\infty(\Delta)=1\)。
\end{proposition}
\begin{proof}
乘积的权相加，所以 \(E_4^3,E_6^2\in M_{12}\)，其差也属于 \(M_{12}\)。两者的常数项相等，故差的常数项为零，得到 \(\Delta\in S_{12}\)。

为确定消失阶，不能停在常数项抵消；还需计算下一项。由\eqref{b3:eq:23-e4-e6-initial}，
\[
 \begin{aligned}
 E_4^3&=1+720q+179280q^2+16954560q^3+O(q^4),\\
 E_6^2&=1-1008q+220752q^2+16519104q^3+O(q^4).
 \end{aligned}
\]
其中平方、立方的二次项分别来自
\[
 2(-16632)+(-504)^2=220752,\quad
 3\cdot2160+3\cdot240^2=179280;
\]
三次项分别来自
\[
 \begin{aligned}
 &2(-122976)+2(-504)(-16632)=16519104,\\
 &3\cdot6720+6\cdot240\cdot2160+240^3=16954560.
 \end{aligned}
\]
相减除以 \(1728\)，就得到\eqref{b3:eq:23-delta-fourier}。一次项系数为一，说明函数非恒零，而且在局部参数 \(q\) 中恰有一阶零点。
\end{proof}

消失阶是在 \(q=0\) 的局部坐标中计数，不是在变量 \(\tau\) 的某个有限点求 Taylor 阶数。具体地，\(\Delta/q=1-24q+252q^2+O(q^3)\) 在 \(q=0\) 附近全纯且非零，因此
\[
 \Delta(x+\mathrm{i}y)
 =e^{2\pi\mathrm{i}x}e^{-2\pi y}
   \left(1+O(e^{-2\pi y})\right)
 \quad(y\longrightarrow\infty),
\]
误差对 \(x\in\mathbb R\) 一致。这证明它在足够高的半平面内没有零点；要排除其余内部零点，还需另一个理由。

\subsection{判别式的两种归一化}
\label{b3:subsec:230404}

前章使用的是格的判别式 \(D(\Lambda)=g_2(\Lambda)^3-27g_3(\Lambda)^2\)，本节使用的是常数已经规范化的模形式 \(\Delta(\tau)\)。两者都具有十二次尺度关系，但数值并不相同。先把比例算准，前章的非退化结论才可以直接用于这里。

\begin{proposition}[格判别式与模判别式]\label{b3:prop:23-delta-nonvanishing}
设 \(\tau\in\mathfrak H\)，\(\Lambda_\tau=\mathbb Z+\mathbb Z\tau\)，并定义
\[
 g_2(\Lambda_\tau)=60\sum_{0\ne\omega\in\Lambda_\tau}\omega^{-4},
 \quad
 g_3(\Lambda_\tau)=140\sum_{0\ne\omega\in\Lambda_\tau}\omega^{-6}.
\]
令 \(D(\Lambda_\tau)=g_2(\Lambda_\tau)^3-27g_3(\Lambda_\tau)^2\)，且
\(\Delta(\tau)=(E_4(\tau)^3-E_6(\tau)^2)/1728\)，其中 \(E_4,E_6\) 为常数项一的 Eisenstein 级数。则
\begin{equation}\label{b3:eq:23-discriminant-normalization}
 D(\Lambda_\tau)=(2\pi)^{12}\Delta(\tau)\ne0.
\end{equation}
更一般地，对每个 \(a\in\mathbb C\setminus\{0\}\)，有
\[
 D(a\Lambda_\tau)=a^{-12}(2\pi)^{12}\Delta(\tau).
\]
\end{proposition}
\begin{proof}
由\eqref{b3:eq:23-g-invariants}及偶数幂和的值，
\[
 g_2(\Lambda_\tau)
 =120\zeta(4)E_4=\frac{4\pi^4}{3}E_4,\quad
 g_3(\Lambda_\tau)
 =280\zeta(6)E_6=\frac{8\pi^6}{27}E_6.
\]
因此
\[
 \begin{aligned}
 D(\Lambda_\tau)
 &=\frac{64\pi^{12}}{27}\left(E_4^3-E_6^2\right)\\
 &=\frac{64\cdot1728}{27}\pi^{12}\Delta
  =4096\pi^{12}\Delta=(2\pi)^{12}\Delta.
 \end{aligned}
\]
因为 \(\operatorname{Im}\tau>0\)，\(\Lambda_\tau\) 是真正的秩二格。\cref{b3:prop:22-discriminant-scaling}已证明，对每一个这样的格，三个非零半周期的 \(\wp\) 值 \(e_1,e_2,e_3\) 互异，而且
\[
 D(\Lambda_\tau)=16\prod_{1\le j<\ell\le3}(e_j-e_\ell)^2\ne0.
\]
于是 \(\Delta(\tau)\ne0\)。同一命题的缩放律 \(D(a\Lambda)=a^{-12}D(\Lambda)\) 给出最后一式。
\end{proof}

这里的“判别式”沿用前章的 \(D\) 约定。若使用多项式 \(4X^3-g_2X-g_3\) 的通常多项式判别式，数值还要乘以 \(16\)。固定非零因子不影响是否有重根，却影响规范化后的精确等式。

\ProseParagraphBreak

尖点处消失与内部没有零点并不冲突。沿 \(\tau=\mathrm{i}y\)、\(y\to\infty\)，\(q\to0\)，由 Fourier 展开，
\[
 g_2\longrightarrow\frac{4\pi^4}{3},\quad
 g_3\longrightarrow\frac{8\pi^6}{27},\quad D\longrightarrow0.
\]
极限的三次多项式为
\[
 4X^3-\frac{4\pi^4}{3}X-\frac{8\pi^6}{27}
 =4\left(X-\frac{2\pi^2}{3}\right)
    \left(X+\frac{\pi^2}{3}\right)^2.
\]
它确实出现了重根，但这是趋向尖点的极限，不是某个有限 \(\tau\in\mathfrak H\) 所对应的格曲线。

\ProseParagraphBreak

至此，\(\Delta\) 在上半平面内部处处可除，而在尖点恰消去一次 \(q\) 因子。这两个性质将在\secref{b3:sec:2305}一起使用：把尖点模形式除以 \(\Delta\) 时，内部不会产生极点，尖点处能否全纯则由消失阶决定。

% !TeX root = ../../Book3.tex
\section{价公式与模形式环}
\label{b3:sec:2305}

普通辐角原理在基本区域上仍有效，但圆弧配对时函数出现权因子，椭圆固定点又只贡献部分圆周。价公式就是把这两种修正逐项记录的结果。

\subsection{基本域中的零点}
\label{b3:subsec:230501}

对非恒零亚纯模形式，记 \(v_p(f)\) 为上半平面局部参数中的零点阶数，极点取负数。模群变换中的因子在上半平面非零，故阶数只依赖轨道。取每条普通轨道一个代表；\(\mathrm{i}\) 与 \(\rho\) 单列。尖点阶数按 \(q\) 参数定义。

\ProseParagraphBreak

固定点的阶数还有同余限制。例如 \(f(-1/\tau)=\tau^kf(\tau)\) 在 \(\mathrm{i}\) 附近比较首项，得到 \((-1)^{v_i}=\mathrm{i}^k\)。权为二时必在 \(\mathrm{i}\) 消失，足以排除非零权二全纯形式；一般同余关系也可以精化维数计数。

\subsection{价公式}
\label{b3:subsec:230502}

\begin{theorem}[价公式]\label{b3:thm:23-valence}
若 \(f\) 是非零权 \(k\) 亚纯模形式，在尖点也亚纯，则
\[
v_\infty(f)+\frac12v_{\mathrm{i}}(f)+\frac13v_\rho(f)
+\sum_{p\ne\mathrm{i},\rho}v_p(f)=\frac{k}{12},
\]
其中求和按普通模群轨道计数。
\end{theorem}
\begin{proof}
先假定边界没有零极点。在高度 \(Y\) 截断基本域，取正向边界：底部圆弧从 \(\rho\) 经 \(\mathrm{i}\) 走到 \(\rho+1\)，顶部由右向左。周期性使竖边的 \(f'/f\) 积分相消；顶部积分趋于 \(-2\pi\mathrm{i}v_\infty(f)\)。
把底弧分为左半弧 \(L:\rho\to\mathrm{i}\) 与右半弧 \(R:\mathrm{i}\to\rho+1\)。映射 \(S\) 将 \(L\) 映成反向的 \(R\)。微分变换律给
\(\mathrm{d}\log f(S\tau)=\mathrm{d}\log f(\tau)+k\,\mathrm{d}\tau/\tau\)。故两弧积分之和等于 \(-k\int_L\mathrm{d}\tau/\tau\)。左弧的辐角从 \(2\pi/3\) 减至 \(\pi/2\)，结果为 \(k\pi\mathrm{i}/6\)。辐角原理除以 \(2\pi\mathrm{i}\) 得内部零极点总数等于 \(k/12-v_\infty\)。
若边界有零极点，先删去其小圆邻域，并将成对边界上的删去部分配成映射下的对应段。\(f'/f=v_p/(\tau-p)+O(1)\)，所以小弧积分的极限等于其角度占 \(2\pi\) 的比例乘 \(2\pi\mathrm{i}v_p\)。普通边界点的两份半圆合成整圆；\(\mathrm{i}\) 只占角度 \(\pi\)；\(\rho,\rho+1\) 两角合计 \(2\pi/3\)。于是分别得到系数一、二分之一、三分之一。被删除的小段对权修正积分的贡献趋零，顶部仍用同一 \(q\) 展开估计。取小圆半径趋零即得所示公式。
\end{proof}
这一结果称为\term[jia-gongshi]{价公式}[valence formula]。Zagier 的《Elliptic Modular Forms and Their Applications》\cite[Proposition 2, pp. 9--10]{Zagier2008Elliptic}可用于核对上述固定点权重；本书在围道计算中显式保留了边界方向。

\subsection{维数论证}
\label{b3:subsec:230503}

\begin{theorem}\label{b3:thm:23-modular-ring}
全模群的分次模形式环为 \(M_* =\mathbb C[E_4,E_6]\)，两个生成元代数独立。对偶数 \(k\ge0\)，
\[
\dim M_k=\begin{cases}\lfloor k/12\rfloor,&k\equiv2\pmod{12},\\
\lfloor k/12\rfloor+1,&k\not\equiv2\pmod{12}.
\end{cases}
\]
负权及奇权空间均为零。
\end{theorem}
\begin{proof}
价公式各项非负，故负权无非零形式；奇权已由 \(-I\) 排除。权零形式无零极点，或用紧化基本域上的最大值原理，可知为常数。权二时 \(v_i\ge1\)，其贡献至少 \(1/2>2/12\)，故空间为零。
因 \(\Delta\) 内部无零点且尖点一阶消失，除以 \(\Delta\) 给出 \(S_k=\Delta M_{k-12}\)。每个偶数 \(k\ge4\) 都可写成 \(4a+6b\)，其中 \(a,b\ge0\)：\(k\equiv0\pmod4\) 时取 \(b=0\)，\(k\equiv2\pmod4\) 时取 \(b=1\)。从 \(f\in M_k\) 减去其常数项乘 \(E_4^aE_6^b\)，余项属于 \(S_k\)。对权归纳，得到所有形式均为两生成元的多项式。每次常数项映射满射到 \(\mathbb C\)，故 \(\dim M_k=1+\dim M_{k-12}\)；结合 \(0,2,4,6,8,10\) 六个初值即得公式。
为证代数独立，先把任何多项式关系按权分解。对固定 \(\tau\)，代入无限多个矩阵 \(\left(\begin{smallmatrix}1&0\\n&1\end{smallmatrix}\right)\)，得到不同权部分乘 \((n\tau+1)^k\) 的多项式恒等式；不同次数的系数必须各为零。因此只需排除同权关系。同权单项式约去一个共同因子后，成为 \(E_4^3/E_6^2\) 的多项式关系。这个亚纯函数非常数，因为两者常数项均一而一次系数分别为 \(720,-1008\)。非常值亚纯函数在没有极点的某小圆盘上像含开集，不可能满足非零常系数多项式，故关系为零。
\end{proof}
例如 \(M_8\) 一维，常数项为一的 \(E_8\) 必等于 \(E_4^2\)。比较 Fourier 系数即得到约数和卷积恒等式。有限维数把一个无穷级数恒等式缩减为有限个系数的核对；减去已匹配项后，价公式保证余项必须为零。

\ProseParagraphBreak

模形式环还允许构造一个保持模性的微分运算。普通导数的变换律含有额外项，必须先计算这项，而不能直接把导数称为高二权形式。记 \(D=(2\pi\mathrm{i})^{-1}\mathrm{d}/\mathrm{d}\tau=q\mathrm{d}/\mathrm{d}q\)，并令
\[
 P(\tau)=\frac{D\Delta(\tau)}{\Delta(\tau)}.
\]
判别式在上半平面无零点，故 \(P\) 全纯。尖点处 \(\Delta=q(1-24q+\cdots)\)，所以 \(P=1-24q+O(q^2)\)。对 \(\Delta(\gamma\tau)=(c\tau+d)^{12}\Delta(\tau)\) 求对数导数，链式法则给
\[
 P(\gamma\tau)=(c\tau+d)^2P(\tau)
       +\frac{12c(c\tau+d)}{2\pi\mathrm{i}}.
\]
额外的加项说明它不是权二模形式；这与 \(M_2=0\) 完全一致。它的任务是抵消其他函数导数中的同类加项。
\ProseParagraphBreak
对权 \(k\) 形式 \(f\) 求导，得到
\[
 (Df)(\gamma\tau)=(c\tau+d)^{k+2}Df(\tau)
       +\frac{kc}{2\pi\mathrm{i}}(c\tau+d)^{k+1}f(\tau).
\]
故 \(D_kf=Df-(k/12)Pf\) 满足权 \(k+2\) 的变换律。其 \(q\) 展开没有负幂；若 \(f\) 是尖点形式，两个乘积也仍无常数项。因此 \(D_k\) 把 \(M_k\) 映入 \(M_{k+2}\)，并保持尖点子空间。这是对普通导数的一项具体修正，全部常数都由权十二的判别式决定。
\ProseParagraphBreak
现在用已经证明的低权维数即可得微分恒等式。\(D_4E_4\in M_6\)，其常数项为 \(-1/3\)，所以 \(D_4E_4=-E_6/3\)。同理 \(D_6E_6\in M_8=\mathbb C E_4^2\)，常数项为 \(-1/2\)。于是
\[
 DE_4=\frac{PE_4-E_6}{3},\qquad
 DE_6=\frac{PE_6-E_4^2}{2}.
\]
再对 \(P\) 的变换律求导，可得 \(DP-P^2/12\) 为权四模形式。具体而言，令 \(J=c\tau+d\)，导数变换中的额外两项是 \(2cJ^3P/(2\pi\mathrm{i})\) 与 \(12c^2J^2/(2\pi\mathrm{i})^2\)；展开 \(P(\gamma\tau)^2/12\) 恰给相同两项，因而完全相消。其常数项为 \(-1/12\)，故
\[
 DP=\frac{P^2-E_4}{12}.
\]
这些式子展示有限维性怎样将微分信息转为代数关系。它们不依赖尚未证明的无限乘积公式，也没有从前几项相同就直接宣称两函数相同；关键是先证明差属于一个已知的一维模形式空间。
\ProseParagraphBreak
也可反过来从这组微分方程递推系数。若写 \(P=1+\sum p_nq^n\)、\(E_4=1+\sum a_nq^n\)、\(E_6=1+\sum b_nq^n\)，每个 \(D\) 把第 \(n\) 项乘以 \(n\)，乘积的第 \(n\) 项则是有限卷积。给定前面的系数后，方程约束后续系数，而 Eisenstein 级数的显式约数和满足这些约束。这是可计算的一致性检验，但如果要据此从零定义三个函数，仍须另外证明所得形式级数的收敛；本章已有全纯构造，故无需倒置这一步。

% !TeX root = ../../Book3.tex
\section{模函数与 \texorpdfstring{\(j\)}{j}-不变量}
\label{b3:sec:2306}

\(E_4\) 与 \(E_6\) 随格基的改变而乘上不同因子，不能直接用它们的数值比较两个复环面。若先把权调成相同，再取商，因子便会约去。问题随之变为：能否找到一个这样的商，使它的值不仅在等价的格上相同，而且足以辨认全部格比例类？下面将用价公式回答这个逆向问题，再给所得参数空间建立具体的复坐标。

\subsection{模函数}
\label{b3:subsec:230601}

\begin{definition}\label{b3:def:23-modular-function}
设 \(\mathfrak H=\{\tau\in\mathbb C:\operatorname{Im}\tau>0\}\)，群 \(\mathrm{SL}_2(\mathbb Z)\) 按分式线性变换作用，\(h\) 为 \(\mathfrak H\) 上的亚纯函数。若
\[
 h\left(\frac{a\tau+b}{c\tau+d}\right)=h(\tau)
 \quad\left(\begin{pmatrix}a&b\\c&d\end{pmatrix}\in\mathrm{SL}_2(\mathbb Z)\right),
\]
且周期函数 \(h\) 的 \(q=e^{2\pi\mathrm i\tau}\) 表达式在 \(q=0\) 亚纯，则称 \(h\) 为全模群的\term[mo-hanshu]{模函数}[modular function]。这里的亚纯函数不包括恒取 \(\infty\) 的映射，允许恒零函数。
\end{definition}

尖点条件允许有限个负次幂；它不要求在尖点有界。若 \(h\not\equiv0\)，则在充分高的半平面有
\[
 h(\tau)=q^m\bigl(a_m+a_{m+1}q+\cdots\bigr),
 \qquad m\in\mathbb Z,\quad a_m\ne0.
\]
正的 \(m\) 表示尖点零点，负的 \(m\) 表示尖点极点。由于 \(|q|=e^{-2\pi\operatorname{Im}\tau}\)，一个 \(r\) 阶尖点极点沿高半平面的增长量级为 \(e^{2\pi r\operatorname{Im}\tau}\)。因此“在无穷远亚纯”说的是指数坐标中的有限阶行为，不能照搬平面坐标 \(1/\tau\) 下的极点定义。

\ProseParagraphBreak
若 \(f,g\in M_k\) 且 \(g\not\equiv0\)，则 \(f/g\) 为模函数。上半平面中的变换因子相消；在任一有限点把分子、分母各自提出其零点幂次，商只有有限阶零点或极点。尖点处同样对两个收敛的 \(q\) 幂级数作局部分解。这说明商在有限点和尖点都亚纯，不能仅凭形式上“权相减为零”省略这两处检查。

\ProseParagraphBreak
非零模函数的倒数、两个模函数的和与积仍为模函数，因而它们组成域。本章稍后将把这个域确定为 \(\mathbb C(j)\)。这里 \(\mathbb C(j)\) 指常系数有理函数代入 \(j\) 所得的函数域，而 \(M_*\) 是按权分次的全纯模形式环；前者允许零权函数有极点，后者用非零权容纳丰富的全纯对象，两者承担不同职责。

\subsection{\texorpdfstring{\(j\)}{j}-不变量}
\label{b3:subsec:230602}

判别式 \(\Delta\) 在上半平面没有零点，且在尖点恰好一阶消失。因此用它作分母，可以把全部极点集中到唯一的尖点。分子须有权十二；选 \(E_4^3\) 还使首项系数为一，这固定了一个方便的归一化。

\begin{definition}\label{b3:def:23-j-invariant}
对 \(\tau\in\mathfrak H\)，令 \(E_4,E_6\) 为全模群上常数项归一为一的权四、权六 Eisenstein 级数，并令 \(\Delta=(E_4^3-E_6^2)/1728\)。定义
\term[j-bubianliang]{\(j\)-不变量}[j-invariant]为
\[
 j(\tau)=\frac{E_4(\tau)^3}{\Delta(\tau)}.
\]
\end{definition}
\symidx{\(j(\tau)\)：椭圆模不变量}

两个权十二因子相除，故 \(j\) 对全模群不变。由前节的 Fourier 展开，
\[
 E_4^3=1+720q+179280q^2+O(q^3),\qquad
 \Delta=q\bigl(1-24q+252q^2+O(q^3)\bigr).
\]
先求第二个括号的倒数：若写成 \(1+b_1q+b_2q^2+O(q^3)\)，乘回原式比较系数得到 \(b_1=24\)、\(b_2=24^2-252=324\)。再与分子相乘，得到
\begin{equation}\label{b3:eq:23-j-fourier}
 j=q^{-1}+744+196884q+O(q^2).
\end{equation}
这里的常数项为 \(720+24\)，一次项为 \(179280+720\cdot24+324\)。归一化中的 \(1728\) 使尖点主部恰为 \(q^{-1}\)；若改取 \(j/1728\)，分类性质不变，但特殊值和全部展开系数都须相应改变。

\Needspace{6\baselineskip}
\begin{theorem}[\texorpdfstring{\(j\)}{j} 对模群轨道的分离]\label{b3:thm:23-j-orbits}
设 \(j=E_4^3/\Delta\)，其中 \(E_4,E_6\) 为全模群归一化 Eisenstein 级数，\(\Delta=(E_4^3-E_6^2)/1728\)，并令 \(\rho=e^{2\pi\mathrm i/3}\)。对每个 \(w\in\mathbb C\)，集合
\(\{\tau\in\mathfrak H:j(\tau)=w\}\) 恰为一条 \(\mathrm{PSL}_2(\mathbb Z)\) 轨道。若 \(p\) 在这条轨道上，\(j-w\) 在上半平面参数 \(\tau-p\) 中的零点阶数为
\[
 \begin{cases}
 3,&p\text{ 与 }\rho\text{ 模群等价},\\
 2,&p\text{ 与 }\mathrm i\text{ 模群等价},\\
 1,&p\text{ 为其他点}.
 \end{cases}
\]
此外 \(j(\rho)=0\)、\(j(\mathrm i)=1728\)。
\end{theorem}
\begin{proof}
固定 \(w\)，考察
\[
 F_w=E_4^3-w\Delta\in M_{12}.
\]
它的常数项为一，故非恒零，且 \(v_\infty(F_w)=0\)。由于 \(\Delta\) 内部无零，\(F_w\) 与 \(j-w\) 在上半平面具有相同零点及阶数。\cref{b3:thm:23-valence}于是给出
\[
 \frac12v_{\mathrm i}(F_w)+\frac13v_\rho(F_w)
 +\sum_{p\ne\mathrm i,\rho}v_p(F_w)=1,
\]
普通点仍按轨道取一个代表。

这个等式本身只给加权总数，必须再知道每条非空零点轨道至少贡献一，才能推出唯一性。普通点的零点阶数为正整数，已有这一性质。在椭圆固定点使用 \(j-w\) 的不变性：由\cref{b3:prop:23-elliptic-orders}的权零情形，它在 \(\mathrm i\) 的阶数为二的倍数，在 \(\rho\) 的阶数为三的倍数。因此这两个位置除以稳定子阶数后的贡献也为正整数。

总数为一首先排除了零点集合为空，又只允许一条零点轨道，并强迫其商参数中的阶数为一。换回上半平面参数，就得到所述三种阶数。最后，由\cref{b3:cor:23-eisenstein-zeros}，\(E_4(\rho)=0\)，而 \(E_6(\mathrm i)=0\)。代入
\[
 j=\frac{E_4^3}{\Delta},\qquad
 j-1728=\frac{E_6^2}{\Delta}
\]
便得到两个特殊值。
\end{proof}

\ProseParagraphBreak
此证明同时给出存在和唯一性。不能先假设某个 \(w\) 被取到，再只证明原像唯一；价公式右侧的正数保证每个有限 \(w\) 都确实有原像。它也说明基本域边界上的两个代表不应算成两个复结构：总数始终按群轨道计，而非按闭基本域中互异点的个数计。

\ProseParagraphBreak
例如 \(E_4^3-1728\Delta=E_6^2\) 的零点全部位于 \(\mathrm i\) 的轨道，每处为二阶；\(E_4^3\) 的零点位于 \(\rho\) 的轨道，每处为三阶。若 \(w\notin\{0,1728\}\)，则 \(E_4^3-w\Delta\) 只有一条普通零点轨道，每处为简单零点。三个情形的加权总数都是一。

\subsection{复环面的 \texorpdfstring{\(j\)}{j} 分类}
\label{b3:subsec:230603}

前面的轨道结论是在参数 \(\tau\) 上得到的。要把它解释成复环面的分类，还须排除一种可能：两个格不成比例，却有某种非线性的全纯坐标变化把它们的商曲面同构起来。平面覆盖使这一问题可以转为整函数的导数估计。

\Needspace{6\baselineskip}
\begin{proposition}[复环面之间的全纯映射]\label{b3:prop:23-torus-affine}
设 \(\Lambda,\Lambda'\subset\mathbb C\) 为秩二格，\(T_\Lambda=\mathbb C/\Lambda\)、\(T_{\Lambda'}=\mathbb C/\Lambda'\) 带有平面商坐标所定义的复结构。每个全纯映射 \(F:T_\Lambda\to T_{\Lambda'}\) 都可写为
\[
 F([z]_\Lambda)=[az+b]_{\Lambda'},\qquad
 a,b\in\mathbb C,\quad a\Lambda\subseteq\Lambda'.
\]
反之，每个满足上述包含条件的仿射式都定义全纯映射。映射 \(F\) 双全纯当且仅当 \(a\ne0\) 且 \(a\Lambda=\Lambda'\)；若还要求 \(F([0])=[0]\)，则可取 \(b=0\)。
\end{proposition}
\begin{proof}
记商投影为 \(\pi_\Lambda,\pi_{\Lambda'}\)。\cref{b3:prop:22-torus-charts}已证明这两个投影是全纯覆盖：取足够小的圆盘，其不同格平移互不相交，投影在每一片上双全纯。平面单连通，故\cref{b3:thm:11-map-lifting}应用于 \(F\circ\pi_\Lambda:\mathbb C\to T_{\Lambda'}\)，给出连续提升 \(\widetilde F:\mathbb C\to\mathbb C\)。提升在局部等于 \(\pi_{\Lambda'}\) 的全纯逆分支与 \(F\circ\pi_\Lambda\) 的复合，因此为整函数。

对每个 \(\lambda\in\Lambda\)，两个点 \(z\) 与 \(z+\lambda\) 的投影相同，所以
\[
 \widetilde F(z+\lambda)-\widetilde F(z)\in\Lambda'.
\]
左端随 \(z\) 连续，右端是离散集合；平面连通使这个差对 \(z\) 为常数。对 \(z\) 求导，便知 \(\widetilde F'\) 为整的 \(\Lambda\)-周期函数。它在闭基本平行四边形上有界，再由周期性在整个平面有界，Liouville 定理使 \(\widetilde F'=a\) 为常数。积分得到 \(\widetilde F(z)=az+b\)，原来的差式给 \(a\lambda\in\Lambda'\)。

反向的良定义也正由这个条件保证：改变代表 \(z\mapsto z+\lambda\)，像代表只增加 \(a\lambda\in\Lambda'\)；局部商坐标中是仿射全纯函数。

若 \(F\) 双全纯，便不恒定，所以 \(a\ne0\)。若有 \(\lambda'\in\Lambda'\setminus a\Lambda\)，则 \([\lambda'/a]_\Lambda\ne[0]_\Lambda\)，却与 \([0]_\Lambda\) 映到同一点，违背单射。因此 \(a\Lambda=\Lambda'\)。反之，等式成立时 \(w\mapsto(w-b)/a\) 良定义地下降为逆映射，故双全纯。若 \(F([0])=[0]\)，则 \(b\in\Lambda'\)，在商中可以去掉这一项。
\end{proof}

这里并未预先假设 \(F\) 保持加法。结论说明，先用目标环面的平移把 \(F([0])\) 移回原点，所得全纯映射自动由复线性映射下降，因此自动是群同态。若只要求实线性同构，则任意两组实基都能相互映到，全部秩二格的商在拓扑上都是环面；复线性这一条件才保留共形结构。

\Needspace{6\baselineskip}
\begin{theorem}[复环面的模参数分类]\label{b3:thm:23-torus-j-classification}
设 \(\Lambda=\omega_1(\mathbb Z+\mathbb Z\tau)\)、\(\Lambda'=\omega_1'(\mathbb Z+\mathbb Z\tau')\)，其中 \(\omega_1,\omega_1'\ne0\)，\(\tau,\tau'\in\mathfrak H\)。以下条件等价：
\begin{enumerate}
\item 复环面 \(\mathbb C/\Lambda\) 与 \(\mathbb C/\Lambda'\) 双全纯；
\item 存在 \(a\in\mathbb C^*\) 使 \(a\Lambda=\Lambda'\)；
\item \(\tau'\) 与 \(\tau\) 属于同一 \(\mathrm{PSL}_2(\mathbb Z)\) 轨道；
\item \(j(\tau')=j(\tau)\)。
\end{enumerate}
每个有限复数均为某个复环面的 \(j\) 值。
\end{theorem}
\begin{proof}
前两项等价由\cref{b3:prop:23-torus-affine}给出，其中常数平移不改变格条件。\cref{b3:prop:23-lattice-orbits}已用正向整数基变换证明：归一化格 \(\mathbb Z+\mathbb Z\tau\) 与 \(\mathbb Z+\mathbb Z\tau'\) 成复比例，当且仅当两参数由行列式一的整数矩阵联系；缩放因子 \(\omega_1,\omega_1'\) 不影响这一判别。最后两项等价以及有限值的满射性由\cref{b3:thm:23-j-orbits}得到。
\end{proof}

\ProseParagraphBreak
对格 \(\Lambda_\tau\)，\secref{b3:sec:2304}给出的常数换算使
\[
 j(\tau)=1728\frac{g_2(\Lambda_\tau)^3}
 {g_2(\Lambda_\tau)^3-27g_3(\Lambda_\tau)^2}.
\]
分子、分母在格缩放下同乘十二次逆幂，故也可直接用任意未归一化的格计算这个比值。方格的 \(g_3=0\) 给 \(j=1728\)，六角格的 \(g_2=0\) 给 \(j=0\)。特殊值表示额外的旋转对称，但两类环面仍各自只有一个同构类，分类定理没有在这些位置失效。

\ProseParagraphBreak
若 \(Y,Y'>1\) 且 \(Y\ne Y'\)，则 \(\mathrm iY,\mathrm iY'\) 是基本域内部不同点，因此对应的矩形环面不双全纯，\(j(\mathrm iY)\ne j(\mathrm iY')\)。另一方面，\(\mathrm iY\) 与 \(\mathrm i/Y\) 由 \(S\) 联系，代表同一环面。这个例子说明“参数不同”与“复结构不同”之间还隔着换基。

\ProseParagraphBreak
由\cref{b3:thm:22-uniformization}，每个格环面已经与它的 Weierstrass 光滑射影三次曲线双全纯对应，原点对应无穷远点。因此上述结论也分类这些由格构造的曲线的保原点解析同构。\(j\) 相等意味着相应复结构同构，并不要求两条仿射方程的系数逐项相等。

\subsection{模曲线 \texorpdfstring{\(X(1)\)}{X(1)} 的几何解释}
\label{b3:subsec:230604}

目前 \(j\) 已给出轨道集合与 \(\mathbb C\) 的双射。集合双射还没有说明连续性、局部复坐标和无穷远点。为理解“模参数空间也是一张 Riemann 曲面”，下面把这些结构逐一建立。

令 \(G=\mathrm{PSL}_2(\mathbb Z)\)，\(Y(1)=G\backslash\mathfrak H\) 带商拓扑，\(\pi:\mathfrak H\to Y(1)\) 为投影。\secref{b3:sec:2301}已证明紧集有限相交性质及局部商模型。普通点附近可取投影单射的小圆盘；在 \(p=\mathrm i\) 或 \(\rho\) 附近，先用
\[
 u_p(\tau)=\frac{\tau-p}{\tau-\overline p}
\]
把稳定子作用化为旋转，再用 \(u_p^e\) 作商坐标，分别取 \(e=2,3\)。其他与这两点等价的位置先经群元搬回即可。这是\cref{b3:prop:23-stabilizer-chart}的具体内容；不能在固定点直接把 \(\pi\) 当成无分歧覆盖。

\ProseParagraphBreak
这些图之间的转移是全纯的。在没有固定点的小片上，选定局部提升后，转移由 Möbius 变换及全纯坐标复合给出；椭圆点中心与普通小片不会同时相交，在去心重叠处可选 \(u_p^e\) 的一个根分支。同一椭圆轨道可固定一个上述坐标图；缩小此图或改用等价代表时，群变换与稳定子旋转相容，相应不变全纯函数按幂级数下降到商坐标。因此这些局部图给出相容的复结构。

\Needspace{6\baselineskip}
\begin{proposition}[唯一尖点的紧化]\label{b3:prop:23-modular-compactification}
设 \(G=\mathrm{PSL}_2(\mathbb Z)\) 按分式线性变换作用于 \(\mathfrak H\)。在 \(Y(1)=G\backslash\mathfrak H\) 中加入一个点 \(\infty\)，规定它的邻域基为
\[
 U_Y=\{\infty\}\cup\pi\{\tau\in\mathfrak H:\operatorname{Im}\tau>Y\},
 \qquad Y>1.
\]
用 \(q=e^{2\pi\mathrm i\tau}\) 将 \(U_Y\) 识别为 \(|q|<e^{-2\pi Y}\)，并以 \(q=0\) 对应新增点。它与普通点坐标、椭圆点坐标 \(u_p^e\) 共同定义一个紧连通 Riemann 曲面，记为 \(X(1)\)。
\end{proposition}
\begin{proof}
\cref{b3:prop:23-cusp-chart}保证当 \(Y>1\) 时，高半平面内只有整数平移可能识别两个点。因此 \(q\) 在 \(U_Y\setminus\{\infty\}\) 上良定义且为一一的局部全纯坐标，像恰为去心圆盘。补上零点就得到完整圆盘。在两个尖点图的重叠上坐标相同；与内部图的重叠上，\(q\) 的导数 \(2\pi\mathrm i q\) 非零，局部对数给出全纯逆，故转移相容。

先证明内部轨道商是 Hausdorff 的。取不属于同一轨道的 \(p,q\in\mathfrak H\)，并取它们的紧邻域 \(K,L\)。由\cref{b3:prop:23-proper-action}，满足 \(\gamma K\cap L\ne\varnothing\) 的群元只有有限个。每个这样的群元均有 \(\gamma p\ne q\)，故可缩小 \(p,q\) 的开邻域，使其在这个群元下不相交；有限次取交后，得到 \(U\subset K,V\subset L\)，使所有这些群元均满足 \(\gamma U\cap V=\varnothing\)。候选集合以外的群元本来就不可能相交。因此 \(\pi(U),\pi(V)\) 是不交的开邻域：商投影为开映射，因为一个开集全部群平移的并仍开。

还须把新尖点与任一内部点分开。固定 \(\tau_0\in\mathfrak H\)，选一个含它的小紧邻域 \(K\)，其中 \(\eta\le\operatorname{Im}\tau\le B\)、\(\eta>0\)。对 \(\tau=x+\mathrm i y\in K\) 及任意群元，若 \(c\ne0\)，则
\[
 \operatorname{Im}(\gamma\tau)
 =\frac{y}{(cx+d)^2+c^2y^2}\le\frac1{c^2y}\le\eta^{-1};
\]
若 \(c=0\)，则虚部仍为 \(y\le B\)。取 \(Y>\max\{1,B,\eta^{-1}\}\)，\(\pi(K)\) 便与 \(U_Y\) 不交，从而尖点与内部点也有不交邻域。

将上半平面的可数基用开商投影送到 \(Y(1)\)，可得 \(Y(1)\) 的可数基。再添上整数高度的 \(U_Y\)，得到 \(X(1)\) 的可数基。因此上述局部图给出 Hausdorff、第二可数的一维复流形。\(Y(1)\) 是连通上半平面的连续像，而每个尖点邻域都与它相交，故添点后仍连通。

最后取 \(Y>1\)。截断基本域
\[
 K_Y=\{\tau\in\mathfrak H:|\operatorname{Re}\tau|\le1/2,
 \ |\tau|\ge1,\ \operatorname{Im}\tau\le Y\}
\]
是紧集：它闭有界，且虚部至少为 \(\sqrt3/2\)。它的投影覆盖高度不超过 \(Y\) 的基本域部分。其余部分加上尖点，在一个稍大尖点图中恰由闭圆盘 \(|q|\le e^{-2\pi Y}\) 参数化，也是紧的。这两个紧集的并覆盖 \(X(1)\)，故 \(X(1)\) 紧致。
\end{proof}

这个紧曲面称为全模群的\term[mo-quxian]{模曲线}[modular curve] \(X(1)\)。此处只讨论这一特定群的解析商及其紧化，不需要一般模空间的代数几何定义。单独加入的尖点不对应一个秩二格环面；当归一化基本域中的虚部趋向无穷时，它记录的是离开所有紧参数集的极限端。

\ProseParagraphBreak
模函数恰好下降为 \(X(1)\) 上的亚纯函数。普通点的下降来自群不变性；在阶数为 \(e\) 的椭圆点，局部 Laurent 级数在 \(u\mapsto\zeta_eu\) 下不变，故只含 \(u^{en}\)，改写成 \(u^e\) 的 Laurent 级数后仍只有有限个负幂。尖点处的亚纯性正是定义中的 \(q\) 条件。反过来，曲面上的亚纯函数拉回 \(\mathfrak H\) 后不变，并在尖点给出同样的有限负幂展开。这一步还说明：有限点在上半平面中的阶数，要除以稳定子阶数，才是商曲面上的阶数。

\Needspace{6\baselineskip}
\begin{theorem}[模曲线的球面坐标]\label{b3:thm:23-j-sphere}
令 \(X(1)\) 为全模群的上述紧化商曲面，\(j=E_4^3/\Delta\)，其中 \(E_4,E_6\) 为全模群归一化 Eisenstein 级数，\(\Delta=(E_4^3-E_6^2)/1728\)。令 \(j(\infty)=\infty\)，则 \(j\) 诱导双全纯映射
\[
 j:X(1)\longrightarrow\widehat{\mathbb C\kern0pt}.
\]
因此全模群的模函数域为 \(\mathbb C(j)\)。其中在整个上半平面全纯的模函数恰为 \(j\) 的多项式，若在尖点也全纯，则必为常数。
\end{theorem}
\begin{proof}
在内部，\cref{b3:thm:23-j-orbits}证明每个有限值恰对应一个商点，并在商坐标中为简单取值。换言之，\(j-j(p)\) 在普通坐标或 \(u_p^e\) 中都有非零一次系数，局部逆函数定理给全纯逆。尖点处，由\eqref{b3:eq:23-j-fourier}，
\[
 \frac1j=q\bigl(1-744q+O(q^2)\bigr),
\]
一次系数也非零；目标无穷远坐标是 \(1/j\)，故此处同样有全纯局部逆。内部没有极点，尖点是无穷远值的唯一原像。因此全局双射及各处全纯局部逆拼成所述双全纯映射。

若 \(h\) 是模函数，它下降到 \(X(1)\) 后与 \(j^{-1}\) 复合，是球面上的亚纯函数。\cref{b3:thm:03-sphere-rational}说明这种函数必为有理函数：球面紧致使极点只有有限多个，减去有限点的主要部分后，只剩一个在无穷远至多有极点的整函数，即多项式。因此 \(h=R(j)\)，其中 \(R\in\mathbb C(t)\)。反向把任意有理函数代入这个球面坐标，当然给出模函数。

若 \(h\) 在上半平面无极点，\(R\) 就不能有有限极点，因为每个有限值都被 \(j\) 取到；有理函数没有有限极点，当且仅当它是多项式。若它还在尖点无极点，该多项式在无穷远也全纯，只能为常数。
\end{proof}

\ProseParagraphBreak
例如一个模函数若只在尖点有至多二阶极点，就具有 \(Aj^2+Bj+C\) 的形式。\(j^2=q^{-2}+1488q^{-1}+\cdots\)，因此先用 \(j^2\) 匹配二阶主部，再用 \(j\) 匹配一阶主部，剩余仅是常数。这样的有限消元由函数域定理保证，不是仅凭若干 Fourier 系数作猜测。

\ProseParagraphBreak
尖点亚纯性在这里不可删除。函数 \(e^j\) 在上半平面全纯且不变，但在尖点有本性奇点：因为 \(1/j\) 是该点的一个全纯局部坐标，\(e^j\) 在此坐标中就是 \(e^{1/z}\)。它既不是模函数，也不是 \(j\) 的有理函数。另一方面，\(j\) 自己已经说明“内部全纯且不变”不强迫常值；必须再控制唯一缺失的尖点。

\ProseParagraphBreak
\(j\) 在上半平面有二阶、三阶局部分歧，与它在 \(X(1)\) 上双全纯并不冲突。先商掉稳定子之后，平方或立方参数已经吸收这些重数，\(j\) 在商坐标中只剩一阶变化。至此，格的换基、模形式的权、价公式中的分数权重和复环面的一个复数参数，便由明确的局部计算衔接起来。

% 本章小结（不计入编号）
\begin{chaptersummary}
模形式同时约束上半平面的对称性与尖点增长。价公式把权转成零点总数，进而给出有限维性、\(E_4,E_6\) 的生成性与 \(\Delta\) 的最小消失阶。权零的比值 \(j\) 消去格尺度，将紧化参数空间识别为球面；这些有限维与系数比较的方法将在平方和问题中继续出现。
\end{chaptersummary}


\BookExercises

% 八题自拟：变换律、固定点、有限维性及系数比较的应用。
\begin{exercise}\label{b3:ex:23-1}
求 \(\mathcal F\) 对面积元 \(y^{-2}\,\mathrm{d}x\,\mathrm{d}y\) 的面积。
\end{exercise}
\begin{solution}
先积 \(y\)，得 \(\int_{-1/2}^{1/2}(1-x^2)^{-1/2}\,\mathrm{d}x=\pi/3\)。无穷远端有有限面积，但并非通常 Euclid 面积有限。
\end{solution}
\begin{exercise}\label{b3:ex:23-2}
确定 \(E_4\) 与 \(E_6\) 的全部零点轨道及其局部阶数。
\end{exercise}
\begin{solution}
\(ST\) 固定 \(\rho\)，权四变换因子不为一，故 \(E_4(\rho)=0\)。价公式总数 \(1/3\) 强迫该点为单零，别处无零。\(S\) 在 \(\mathrm{i}\) 的权六因子为 \(-1\)，故 \(E_6\) 在该点消失；总数 \(1/2\) 强迫单零且无其他轨道。
\end{solution}
\begin{exercise}\label{b3:ex:23-3}
证明两个 \(M_k\) 中的形式，只要从常数项起前 \(\lfloor k/12\rfloor+1\) 个系数相同，就完全相同。
\end{exercise}
\begin{solution}
差的尖点阶数大于 \(k/12\)。若差非零，价公式左侧已经仅由这一项超过右侧，矛盾。此为充分界，在 \(k\equiv2\pmod{12}\) 等情形可以改进。
\end{solution}
\begin{exercise}\label{b3:ex:23-4}
求 \(M_{14}\)、\(S_{24}\) 的基。
\end{exercise}
\begin{solution}
维数分别为一、二。前者由常数项一的 \(E_4^2E_6\) 张成。后者为 \(\Delta M_{12}\)，故可取 \(\Delta E_4^3,\Delta^2\)；两者尖点阶分别为一、二，线性独立。
\end{solution}
\begin{exercise}\label{b3:ex:23-5}
用 \(E_4^2=E_8\) 推导约数和卷积恒等式。
\end{exercise}
\begin{solution}
由一般 Eisenstein 展开有 \(E_8=1+480\sum\sigma_7(n)q^n\)，其中常数由 \(\zeta(8)=\pi^8/9450\) 得到，此值也由余切展开比较七次项算出。比较 \(q^n\) 得
\(\sum_{m=1}^{n-1}\sigma_3(m)\sigma_3(n-m)=(\sigma_7(n)-\sigma_3(n))/120\)。空和 \(n=1\) 时为零，右侧也零。
\end{solution}
\begin{exercise}\label{b3:ex:23-6}
证明在上半平面全纯且在尖点有界的模函数必常数，并给出删去尖点条件的反例。
\end{exercise}
\begin{solution}
有界性使 \(q=0\) 可去，函数下降为紧曲面 \(X(1)\) 上的全纯函数，最大值原理使其常数。\(j\) 在上半平面全纯且不变，但在尖点有极点，所以提供反例。
\end{solution}
\begin{exercise}\label{b3:ex:23-7}
若模函数只有尖点处至多 \(n\) 阶极点，证明它是次数至多 \(n\) 的 \(j\) 多项式。
\end{exercise}
\begin{solution}
利用 \(j=q^{-1}+744+\cdots\)，逐次减去 \(j^n,j^{n-1},\ldots,j\) 的适当倍数可消去全部负幂。余下函数在紧曲面全纯，故常数。最高极点阶数同时给多项式次数的上界。
\end{solution}
\begin{exercise}\label{b3:ex:23-8}
证明 \(j-1728=E_6^2/\Delta\)，并确定 \(j\) 在 \(\rho\)、\(j-1728\) 在 \(\mathrm{i}\) 的消失阶。
\end{exercise}
\begin{solution}
第一式由 \(E_4^3-E_6^2=1728\Delta\) 得到。判别式内部无零，且前题给 \(E_4,E_6\) 各自的单零，故两消失阶分别为三、二。在商空间的立方、平方局部参数中都成为一阶，这与 \(j\) 双全纯并不矛盾。
\end{solution}

% !TeX root = ../../Book3.tex
\chapter{\texorpdfstring{\(\theta\)}{\AnalysisPdfTheta} 函数与整数平方和}
\label{b3:ch:24}
\BookChapterNumber{b3:ch:24}
% 本章摘要（不计入编号）
\begin{chapterabstract}
theta 级数把高斯求和、复函数的零点和整数表示数放在同一个函数中。本章先从收敛与指标替换建立四种 theta 函数、热方程及模变换，再以有限恒等式的极限证明 Jacobi 三重积。乘积确定零点和导数常数，椭圆函数的半周期值将这些资料转成四平方和的约数公式。二平方和另由 Gaussian 整数的除法与唯一分解得到。末节说明一般格的反演律包含对偶格和体积，并在偶自对偶条件下回到前章的全模群模形式。
\end{chapterabstract}

% 本章目标（不计入编号）
\begin{learninggoals}
能够核对 \(q=e^{\pi\mathrm{i}\tau}\) 与前章变量的区别，说明级数求导、乘积取极限和系数重排的依据；能够从乘积找出 theta 零点，从表示数的定义重建生成函数；能够说明二平方和的素因子障碍，并把半周期差、热方程、Lambert 级数连接成四平方和证明。选读后还应能辨别一般格的反演公式与全模群模形式的条件，并计算八维偶自对偶格的短向量数。
\end{learninggoals}

把
\[
1+2q+2q^4+2q^9+\cdots
=\sum_{n\in\mathbb Z}q^{n^2}
\]
平方，\(q^m\) 的系数就在数两个平方相加等于 \(m\) 的方式。这里系数二记录正负号；例如 \(q^1\) 来自 \(1^2\) 与 \((-1)^2\)，而常数项只有零的一种选择。这个生成函数容易写出，但把它相乘并不能立刻给出所有系数的公式。需要找到一种表达式，让整数 \(m\) 的约数或素因子能够显露出来。

本章在三个表示之间来回转换：级数便于计数和微分，乘积便于判断零点和取对数导数，变换公式则把一个参数区域与另一个区域联系起来。每次转换都有具体用途。Poisson 求和解释高斯参数的反演；三重积使 theta 常数的非零性变得可见；椭圆函数的主部比较再把一个对数导数变成 theta 常数的四次幂。最后一步展开几何级数，才得到约数求和。

\ProseParagraphBreak
阅读本章时，只需从前面带来少量明确工具：正规收敛允许逐项求导，正规乘积的非零因子给出非零极限；\chapref{b3:ch:22}中无极点的椭圆函数为常数，Weierstrass 函数在原点有展开 \(\wp(z)=z^{-2}+O(z^2)\)。后面真正使用这些结论时会再次指出其作用。所需的高斯求和在本章重建；二平方和需要的 Gaussian 整数分解也完整证明，不预设代数数论背景。末节的偶自对偶格例子才进一步调用\chapref{b3:ch:23}的模形式空间。

\secref{b3:sec:2401}至\secref{b3:sec:2404}构成基本路线；\secref{b3:sec:2405}将同一方法推广到高维格，可作为选读。Whittaker、Watson 的《A Course of Modern Analysis》\cite{Whittaker1996Course}可供查阅 theta 与椭圆函数的不同记号；Zagier 的《Elliptic Modular Forms and Their Applications》\cite[pp. 26--28]{Zagier2008Elliptic}还说明用模形式空间比较平方和系数的方法。阅读外部公式时，应先核对空间变量的周期以及 \(q\) 的定义，再比较数值常数。

% !TeX root = ../../Book3.tex
\section{Jacobi \texorpdfstring{\(\theta\)}{\AnalysisPdfTheta} 函数}
\label{b3:sec:2401}

若要用一个函数同时保存所有平方数，最直接的办法是把平方数放到幂指数上，写成 \(\sum_{n\in\mathbb Z}q^{n^2}\)。当 \(|q|<1\) 时，这个级数迅速收敛；乘起若干份后，系数就在计数平方和。它还有一个计数问题本身不容易看出的对称：令 \(q=e^{-\pi t}\)，把 \(t\) 换成 \(1/t\)，级数只改变一个简单因子。这一对称来自高斯函数的 Fourier 变换。下面加入一个记录线性相位的空间变量 \(z\)，把这项事实建立为可以求导、平移和作参数变换的函数恒等式。

\subsection{\texorpdfstring{\(\theta\)}{\AnalysisPdfTheta} 级数}
\label{b3:subsec:240101}

\begin{definition}\label{b3:def:24-theta-series}
设 \(\tau\in\mathfrak H=\{\,w\in\mathbb C\mid\operatorname{Im}w>0\,\}\)、\(z\in\mathbb C\)，置 \(q=e^{\pi\mathrm i\tau}\)。四个\term[jacobi-theta-hanshu]{Jacobi theta 函数}[Jacobi theta function]定义为
\begin{equation}\label{b3:eq:24-theta-series}
\begin{aligned}
 \theta_1(z,\tau)&=-\mathrm i\sum_{n\in\mathbb Z}(-1)^n
                q^{(n+1/2)^2}e^{(2n+1)\mathrm iz},\\
 \theta_2(z,\tau)&=\sum_{n\in\mathbb Z}
                q^{(n+1/2)^2}e^{(2n+1)\mathrm iz},\\
 \theta_3(z,\tau)&=\sum_{n\in\mathbb Z}q^{n^2}e^{2n\mathrm iz},\\
 \theta_4(z,\tau)&=\sum_{n\in\mathbb Z}(-1)^nq^{n^2}e^{2n\mathrm iz}.
\end{aligned}
\end{equation}
这里对任意实数 \(a\)，记号 \(q^a\) 一律表示 \(e^{\pi\mathrm i a\tau}\)，不另选复数 \(q\) 的幂分支。省略第二个变量时，始终保持同一个 \(\tau\)。
\end{definition}
\symidx{\(\theta_j(z,\tau)\)：Jacobi theta 函数，\(1\le j\le4\)}

本章的 \(q\) 与前章的模形式展开变量不同：前章用 \(e^{2\pi\mathrm i\tau}\)，本章用它的平方根 \(e^{\pi\mathrm i\tau}\)。这一选择使 \(\theta_3(0,\tau)\) 的幂指数直接成为整数平方，而不再带二分之一。对 \(\theta_1,\theta_2\)，指数还含四分之一；它们首先是 \(\tau\) 的单值函数，不能未经说明就把它们看成穿孔圆盘中 \(q\) 的单值函数。
\ProseParagraphBreak
四个函数不是四种互不相关的构造。由整数格改为半整数格，就从 \(\theta_3\) 得到 \(\theta_2\)；再给相邻格点加上交替符号，就出现 \(\theta_4\) 和 \(\theta_1\)。这些变化在 Fourier 变换下会彼此交换，所以把它们一并引入，可以让后面的变换律保持封闭。先把 \(n\) 与 \(-n\)、或 \(n\) 与 \(-n-1\) 配对，得到可直接计算的实三角形式：
\[
\begin{aligned}
 \theta_1(z,\tau)&=2\sum_{n\ge0}(-1)^nq^{(n+1/2)^2}
                                      \sin((2n+1)z),\\
 \theta_2(z,\tau)&=2\sum_{n\ge0}q^{(n+1/2)^2}
                                      \cos((2n+1)z),\\
 \theta_3(z,\tau)&=1+2\sum_{n\ge1}q^{n^2}\cos(2nz),\\
 \theta_4(z,\tau)&=1+2\sum_{n\ge1}(-1)^nq^{n^2}\cos(2nz).
\end{aligned}
\]
因此 \(\theta_1\) 关于 \(z\) 为奇函数，其余三个为偶函数；\(\theta_1(0,\tau)=0\)。这里的配对和随后使用的逐项求导，都由下面的同一估计保证。

\begin{proposition}\label{b3:prop:24-theta-normal}
对 \(j\in\{1,2,3,4\}\)，\cref{b3:def:24-theta-series}的级数及其任意固定阶混合导数
\(\partial_z^r\partial_\tau^s\)，其中 \(r,s\) 为非负整数，都在
\(\mathbb C\times\mathfrak H\) 的每个紧子集上绝对一致收敛。因而 \(\theta_j(z,\tau)\) 对 \(z\) 为整函数，对 \(\tau\in\mathfrak H\) 全纯，所述导数可以逐项计算。
\end{proposition}
\begin{proof}
在一个给定紧子集上，可以选取 \(\eta>0\)、\(R>0\)，使
\(\operatorname{Im}\tau\ge\eta\)、\(|\operatorname{Im}z|\le R\)。四种级数的指标都是 \(u=n\) 或 \(u=n+1/2\)，除去模为一的系数后，各项及导数的模不超过
\[
 C_{r,s}|u|^{r+2s}\exp(-\pi\eta u^2+2R|u|).
\]
当 \(|u|\) 足够大时，\(2R|u|\le\pi\eta u^2/2\)，故尾项被
\(C_{r,s}|u|^{r+2s}e^{-\pi\eta u^2/2}\) 控制。这个数列可求和，因为任意固定次多项式都被高斯衰减吸收。有限个小指标另取上界，Weierstrass 判别法便给出一致收敛；对单个变量应用全纯函数列的逐项求导定理，即得到各阶混合导数。
\end{proof}

上半平面条件确实参与了论证。若 \(\tau\) 为实数、\(z=0\)，则 \(\theta_3\) 的各项模恒为一，连级数收敛所需的项趋于零都不满足。不能把上面的紧集估计直接延到 \(\operatorname{Im}\tau=0\)。
\ProseParagraphBreak
平方指数还带来一个将在四平方和证明中使用的微分关系：对参数求导产生平方，对空间变量求两次导数也产生同一个平方。

\begin{proposition}\label{b3:prop:24-theta-heat}
对 \(z\in\mathbb C\)、\(\tau\in\mathfrak H\) 和 \(1\le j\le4\)，有
\begin{equation}\label{b3:eq:24-theta-heat}
 \frac1{\pi\mathrm i}\partial_\tau\theta_j(z,\tau)
       =-\frac14\partial_z^2\theta_j(z,\tau).
\end{equation}
若以 \(q=e^{\pi\mathrm i\tau}\) 作局部坐标，则左端为
\(q\partial_q\theta_j\)。特别地，函数 \(u_j(x,t)=\theta_j(x,\mathrm it)\)，
\(x\in\mathbb R\)、\(t>0\)，满足
\(\partial_tu_j=(\pi/4)\partial_x^2u_j\)。
\end{proposition}
\begin{proof}
每项都具有常数乘 \(e^{\pi\mathrm i u^2\tau+2\mathrm iuz}\) 的形式，
其中 \(u=n\) 或 \(n+1/2\)。算子 \((\pi\mathrm i)^{-1}\partial_\tau\)
乘出 \(u^2\)，而 \(\partial_z^2\) 乘出 \(-4u^2\)。由
\cref{b3:prop:24-theta-normal}逐项计算并求和，得到恒等式。
由于 \(\mathrm dq/\mathrm d\tau=\pi\mathrm i q\ne0\)，局部坐标中的链式法则给出 \(q\partial_q=(\pi\mathrm i)^{-1}\partial_\tau\)。
最后 \(\partial_tu_j=\mathrm i\partial_\tau\theta_j\)，代入即得实时间的热方程。
\end{proof}

\begin{proposition}\label{b3:prop:24-theta-translations}
设 \(\tau\in\mathfrak H\)、\(z\in\mathbb C\)，并在下式中简记
\(\theta_j(z)=\theta_j(z,\tau)\)、\(q^a=e^{\pi\mathrm i a\tau}\)。
四个 theta 函数满足
\begin{equation}\label{b3:eq:24-theta-quasiperiods}
\begin{aligned}
 \theta_1(z+\pi)&=-\theta_1(z),&
 \theta_1(z+\pi\tau)&=-q^{-1}e^{-2\mathrm iz}\theta_1(z),\\
 \theta_2(z+\pi)&=-\theta_2(z),&
 \theta_2(z+\pi\tau)&=q^{-1}e^{-2\mathrm iz}\theta_2(z),\\
 \theta_3(z+\pi)&=\theta_3(z),&
 \theta_3(z+\pi\tau)&=q^{-1}e^{-2\mathrm iz}\theta_3(z),\\
 \theta_4(z+\pi)&=\theta_4(z),&
 \theta_4(z+\pi\tau)&=-q^{-1}e^{-2\mathrm iz}\theta_4(z).
\end{aligned}
\end{equation}
半周期平移的结果由\cref{b3:tab:24-half-translations}给出。
\end{proposition}

\begin{table}[htbp]
\centering
\caption[theta 函数的半周期平移]{theta 函数的半周期平移。共同参数为
\(\tau\in\mathfrak H\)，\(q=e^{\pi\mathrm i\tau}\)，每一行对全部 \(z\in\mathbb C\) 成立。第三列已乘去共同指数因子，因而只留下函数交换及一个常数。}
\label{b3:tab:24-half-translations}
\begin{tabular}{ccc}
\toprule
 \(j\) & \(\theta_j(z+\pi/2,\tau)\) &
 \(q^{1/4}e^{\mathrm iz}\theta_j(z+\pi\tau/2,\tau)\)\\
\midrule
 \(1\) & \(\theta_2(z,\tau)\) & \(\mathrm i\theta_4(z,\tau)\)\\
 \(2\) & \(-\theta_1(z,\tau)\) & \(\theta_3(z,\tau)\)\\
 \(3\) & \(\theta_4(z,\tau)\) & \(\theta_2(z,\tau)\)\\
 \(4\) & \(\theta_3(z,\tau)\) & \(\mathrm i\theta_1(z,\tau)\)\\
\bottomrule
\end{tabular}
\end{table}
\begin{proof}
对 \(z\mapsto z+\pi\)，整数频率的项不变，半整数频率的项乘以 \(-1\)，给出第一列公式。对 \(z\mapsto z+\pi\tau\)，把指数
\(u^2+2u\) 写成 \((u+1)^2-1\)，再把指标平移一位，提出共同因子 \(q^{-1}e^{-2\mathrm iz}\)。带 \((-1)^n\) 的级数在移指标后另产生一个负号，故得到第二列。
半周期表的第二列来自
\(e^{2\mathrm in(\pi/2)}=(-1)^n\) 与
\(e^{(2n+1)\mathrm i\pi/2}=\mathrm i(-1)^n\)。
对第三列，整数指标满足
\(n^2+n=(n+1/2)^2-1/4\)，半整数指标满足
\((n+1/2)^2+(n+1/2)=(n+1)^2-1/4\)。
例如对奇函数作这个移项，得到
\[
\begin{aligned}
 \theta_1(z+\pi\tau/2,\tau)
 &=-\mathrm i q^{-1/4}
       \sum_{n\in\mathbb Z}(-1)^nq^{(n+1)^2}e^{(2n+1)\mathrm iz}\\
 &=\mathrm i q^{-1/4}e^{-\mathrm iz}\theta_4(z,\tau).
\end{aligned}
\]
去掉交替符号就得到
\(\theta_2(z+\pi\tau/2,\tau)=q^{-1/4}e^{-\mathrm iz}\theta_3(z,\tau)\)。
整数指标的两行直接用第一项配方，分别得到
\(q^{-1/4}e^{-\mathrm iz}\theta_2(z,\tau)\) 与
\(\mathrm i q^{-1/4}e^{-\mathrm iz}\theta_1(z,\tau)\)，完成全表。
所有移项均在绝对收敛级数中进行。
\end{proof}

第二个方向的平移带有依赖 \(z\) 的乘法因子，这称为准周期性，不能把 \(\pi\tau\) 称作真正周期。与第二十二章的椭圆函数相比，这个差别很必要：非恒定整函数若有两个实线性无关的真正周期，就会在一个紧基本平行四边形上有界，继而在全平面有界，与 Liouville 定理矛盾。theta 函数允许大小在相邻格胞间按指数因子改变，因而能够同时保持整性与非平凡的格对称。取一次对数导数，乘法指数变成加法常数；再求导一次，常数消失，才会产生真正双周期的亚纯函数。

\subsection{与第二卷 Poisson 求和的联系}
\label{b3:subsec:240102}

第二卷从 Fourier 理论研究过 Poisson 求和。本章只需要高斯及其平移、调制的情形，下面直接证明这个版本，并把积分换序和 Fourier 展开成立的理由写出。这样，即使读者忘记了一般定理的函数空间条件，也可以沿着当前论证继续学习。

\begin{lemma}\label{b3:lem:24-gaussian-poisson}
设 \(t>0\)、\(a\in\mathbb R\)、\(b\in\mathbb C\)。则下列两边级数绝对收敛，并且
\begin{equation}\label{b3:eq:24-gaussian-poisson}
 \sum_{n\in\mathbb Z}e^{-\pi t(n+a)^2}e^{2\pi\mathrm ib(n+a)}
 =t^{-1/2}\sum_{m\in\mathbb Z}
       e^{-\pi(m-b)^2/t}e^{2\pi\mathrm ima}.
\end{equation}
特别地，\(a=b=0\) 时有
\[
 \sum_{n\in\mathbb Z}e^{-\pi tn^2}
     =t^{-1/2}\sum_{m\in\mathbb Z}e^{-\pi m^2/t}.
\]
\end{lemma}
\begin{proof}
固定 Fourier 变换约定
\(\widehat f(\xi)=\int_{\mathbb R}f(x)e^{-2\pi\mathrm ix\xi}\,\mathrm dx\)。
先计算 \(g_t(x)=e^{-\pi tx^2}\) 的变换。高斯积分为
\(\int_{\mathbb R}e^{-\pi tx^2}\,\mathrm dx=t^{-1/2}\)：在 \(t=1\)
时将积分平方化成平面上的非负积分，极坐标给出
\(2\pi\int_0^\infty e^{-\pi r^2}r\,\mathrm dr=1\)，一般情形再缩放。
令 \(F(\xi)=\widehat g_t(\xi)\)，对实 \(\xi\) 在积分下求导，并用
\(g_t'(x)=-2\pi txg_t(x)\) 分部积分，得到
\[
 F'(\xi)=\frac{\mathrm i}{t}
       \int_{\mathbb R}g_t'(x)e^{-2\pi\mathrm ix\xi}\,\mathrm dx
       =-\frac{2\pi\xi}{t}F(\xi).
\]
这里 \(g_t\) 与 \(xg_t\) 可积，且 \(g_t(x)\to0\) 保证边界项消失。
初值 \(F(0)=t^{-1/2}\) 因而给出
\[
 \widehat g_t(\xi)=t^{-1/2}e^{-\pi\xi^2/t}.
\]
对实数 \(b\)，令 \(f_b(x)=g_t(x)e^{2\pi\mathrm ibx}\)，于是
\(\widehat f_b(\xi)=t^{-1/2}e^{-\pi(\xi-b)^2/t}\)。

将 \(f_b\) 周期化为
\(P_b(x)=\sum_{n\in\mathbb Z}f_b(x+n)\)。当 \(x\in[0,1]\) 时，
\(|x+n|\ge |n|-1\)，而 \(f_b\) 的各阶导数都是多项式乘高斯再乘一个模为一的相位。故级数及其各阶导数一致收敛，\(P_b\) 为光滑的一周期函数。其第 \(m\) 个 Fourier 系数是
\[
\begin{aligned}
 \int_0^1P_b(x)e^{-2\pi\mathrm imx}\,\mathrm dx
 &=\sum_{n\in\mathbb Z}\int_0^1f_b(x+n)e^{-2\pi\mathrm imx}\,\mathrm dx\\
 &=\int_{\mathbb R}f_b(u)e^{-2\pi\mathrm imu}\,\mathrm du
   =t^{-1/2}e^{-\pi(m-b)^2/t}.
\end{aligned}
\]
第一步的绝对积分总和等于 \(\int_{\mathbb R}|f_b(u)|\,\mathrm du<\infty\)，
第二步用 \(m,n\in\mathbb Z\) 消去因子 \(e^{2\pi\mathrm imn}\)。
这些系数本身快速衰减，所以相应 Fourier 级数绝对一致收敛，定义一个连续函数 \(Q_b\)，且 \(P_b,Q_b\) 的全部 Fourier 系数相同。

为说明“系数相同”确实推出函数相同，可以对差 \(h=P_b-Q_b\) 用核函数
\[
 K_N(x)=\frac1N\left|\sum_{k=0}^{N-1}e^{2\pi\mathrm ikx}\right|^2.
\]
它非负、在 \([0,1]\) 上积分为一；离整数点固定距离时，
等比数列公式给 \(K_N(x)=O(N^{-1})\)。连续周期函数一致连续，把
\(\int_0^1K_N(u)(h(x-u)-h(x))\,\mathrm du\) 分成整数点附近与其余部分，前者由一致连续性控制，后者由上述界趋于零。因此
\(K_N*h\to h\) 一致收敛。另一方面，\(K_N\) 为有限三角多项式，而
\(h\) 的全部 Fourier 系数为零，故 \(K_N*h=0\)，从而 \(h=0\)。
在 \(P_b=Q_b\) 中代入 \(x=a\)，便得到实 \(b\) 时的求和公式。

最后将 \(b\) 延到复数。若 \(b\) 位于一个给定紧集，两侧项的模分别被
\(C e^{-\pi t(n+a)^2+C|n+a|}\) 和
\(C e^{-\pi m^2/t+C|m|}\) 控制。因此两边关于 \(b\) 都整，且在实轴相等。恒等定理给出全部 \(b\in\mathbb C\) 的等式，同时上述估计也证明绝对收敛。
\end{proof}

公式中的两个操作有清楚的对应：空间平移 \(a\) 变成右端的相位
\(e^{2\pi\mathrm ima}\)，空间调制 \(b\) 变成右端频率中心从 \(0\)
移到 \(b\)。取 \(a=1/2,b=0\) 时，左端在半整数格上求和，
右端出现 \((-1)^m\)；这正是 \(\theta_2\) 与 \(\theta_4\) 会在参数反演中交换的原因。

\subsection{模变换}
\label{b3:subsec:240103}

第二十三章的模群由 \(T:\tau\mapsto\tau+1\) 与
\(S:\tau\mapsto-1/\tau\) 生成。对 theta 函数，\(T\) 可以直接在级数中计算，\(S\) 则由刚证明的高斯求和公式产生。空间变量也须随 \(S\) 改为 \(z/\tau\)；只变参数而保持 \(z\) 不动，得不到下面的形式。

\begin{theorem}\label{b3:thm:24-theta-modular}
设 \(z\in\mathbb C\)、\(\tau\in\mathfrak H\)，并定义
\[
 A(\tau)=(-\mathrm i\tau)^{1/2}
   =\exp\left(\tfrac12\operatorname{Log}(-\mathrm i\tau)\right),
 \quad
 B(z,\tau)=A(\tau)e^{\mathrm i z^2/(\pi\tau)}.
\]
这里 \(\operatorname{Log}\) 是右半平面上辐角在
\((-\pi/2,\pi/2)\) 内的全纯对数，因此 \(A(\mathrm it)=\sqrt t>0\)
对所有 \(t>0\) 成立。则
\begin{equation}\label{b3:eq:24-theta-T}
\begin{aligned}
 \theta_1(z,\tau+1)&=e^{\pi\mathrm i/4}\theta_1(z,\tau),&
 \theta_2(z,\tau+1)&=e^{\pi\mathrm i/4}\theta_2(z,\tau),\\
 \theta_3(z,\tau+1)&=\theta_4(z,\tau),&
 \theta_4(z,\tau+1)&=\theta_3(z,\tau),
\end{aligned}
\end{equation}
以及
\begin{equation}\label{b3:eq:24-theta-S}
\begin{aligned}
 \theta_1(z/\tau,-1/\tau)&=-\mathrm iB(z,\tau)\theta_1(z,\tau),\\
 \theta_2(z/\tau,-1/\tau)&=B(z,\tau)\theta_4(z,\tau),\\
 \theta_3(z/\tau,-1/\tau)&=B(z,\tau)\theta_3(z,\tau),\\
 \theta_4(z/\tau,-1/\tau)&=B(z,\tau)\theta_2(z,\tau).
\end{aligned}
\end{equation}
\end{theorem}
\begin{proof}
对整数指标，\(e^{\pi\mathrm in^2}=(-1)^n\)，所以 \(T\) 交换
\(\theta_3,\theta_4\)。对半整数指标，
\[
 e^{\pi\mathrm i(n+1/2)^2}
   =e^{\pi\mathrm i/4}e^{\pi\mathrm in(n+1)}
   =e^{\pi\mathrm i/4},
\]
因为 \(n(n+1)\) 为偶数。这给出 \(\theta_1,\theta_2\) 的相位因子。

为同时算清 \(S\) 的下标和符号，暂记
\[
 \Phi_{a,b}(z,\tau)=
 \sum_{n\in\mathbb Z}e^{\pi\mathrm i\tau(n+a)^2+
                         2\mathrm i(n+a)(z+\pi b)},
 \quad a,b\in\mathbb R.
\]
这只是统一计算的记号；四个函数分别为
\[
 \theta_3=\Phi_{0,0},\quad
 \theta_2=\Phi_{1/2,0},\quad
 \theta_4=\Phi_{0,1/2},\quad
 \theta_1=-\Phi_{1/2,1/2}.
\]
先令 \(\tau=\mathrm it\)、\(t>0\)，在
\cref{b3:lem:24-gaussian-poisson}中把 \(t\) 换成 \(1/t\)，把调制参数
换成 \(b+z/(\pi\mathrm it)\)。该参数可以为复数，引理已经覆盖这种情形。
于是
\[
 \Phi_{a,b}\left(\frac z{\mathrm it},\frac{\mathrm i}t\right)
 =\sqrt t\sum_{m\in\mathbb Z}
    e^{-\pi t(m-b+\mathrm iz/(\pi t))^2}e^{2\pi\mathrm ima}.
\]
展开平方，提出 \(e^{z^2/(\pi t)}\)，再令 \(n=-m\)，得到
\[
 \Phi_{a,b}\left(\frac z{\mathrm it},\frac{\mathrm i}t\right)
 =\sqrt t\,e^{z^2/(\pi t)}e^{2\pi\mathrm iab}
                    \Phi_{b,-a}(z,\mathrm it).
\]
对每个固定 \(z\in\mathbb C\)，等式两边关于 \(\tau\in\mathfrak H\)
都全纯；这一点分别来自正规收敛、\(\tau\ne0\) 和已经选定的
\(A(\tau)\)。在虚轴的全部 \(\tau=\mathrm it\) 上相等，恒等定理便给
\[
 \Phi_{a,b}(z/\tau,-1/\tau)
 =B(z,\tau)e^{2\pi\mathrm iab}\Phi_{b,-a}(z,\tau)
\]
在整个上半平面成立。代入四组 \(a,b\)，即可得到反演公式。
在 \(a=b=1/2\) 的一行还要用
\(\Phi_{1/2,-1/2}=-\Phi_{1/2,1/2}\)，故总因子为
\(e^{\pi\mathrm i/2}(-1)=-\mathrm i\)，而不是 \(\mathrm i\)。
\end{proof}

在 \(\tau=\mathrm it\) 时，反演式中的高斯因子为
\(e^{z^2/(\pi t)}\)。若把等式反解成 \(\theta_3(z,\mathrm it)\)，才出现
衰减因子 \(e^{-z^2/(\pi t)}\)。这个方向区别可以防止记错指数的符号。
整个证明只对实高斯作 Fourier 变换，再作全纯延拓，没有未经估计就平移一条无穷积分直线。
\ProseParagraphBreak
平方根的选取也不能在连续两次变换时随意改变。事实上
\(A(-1/\tau)A(\tau)=1\)：乘积的平方为一，在虚轴取正值，而上半平面连通。
两次 \(S\) 把空间变量 \(z\) 变为 \(-z\)，高斯因子互相抵消。
三个偶函数的结果不变，\(\theta_1\) 的两个因子
\((-\mathrm i)^2=-1\) 则正好恢复它的奇性。这为根分支和相位提供了一个直接核对。

\subsection{\texorpdfstring{\(\theta\)}{\AnalysisPdfTheta} 常数}
\label{b3:subsec:240104}

\begin{definition}\label{b3:def:24-theta-constants}
对 \(\tau\in\mathfrak H\)，四个值
\(\theta_j(0,\tau)\)、\(1\le j\le4\)，称为\term[theta-changshu]{theta 常数}[theta constant]。
当空间变量已固定为零且参数明确时，简记为
\(\theta_j(\tau)\) 或 \(\theta_j\)。“常数”只表示不再依赖 \(z\)，并不表示它们不随 \(\tau\) 改变。
\end{definition}

其中 \(\theta_1(\tau)=0\)。另外三个常数的展开为
\begin{equation}\label{b3:eq:24-theta-constant-series}
\begin{aligned}
 \theta_2(\tau)&=2q^{1/4}(1+q^2+q^6+q^{12}+\cdots),\\
 \theta_3(\tau)&=1+2q+2q^4+2q^9+\cdots,\\
 \theta_4(\tau)&=1-2q+2q^4-2q^9+\cdots,
 \quad q=e^{\pi\mathrm i\tau}.
\end{aligned}
\end{equation}
第一行括号中的一般指数为 \(n(n+1)\)、\(n\ge0\)，后两行的指数为
\(n^2\)、\(n\ge1\)。这三者在上半平面都不为零；该事实将在下一节由乘积公式证明，目前尚不凭展开的首项排除其余项的复数抵消。
\ProseParagraphBreak
在\cref{b3:thm:24-theta-modular}中取 \(z=0\)，得到常数反演
\[
 \theta_2(-1/\tau)=A(\tau)\theta_4(\tau),\quad
 \theta_3(-1/\tau)=A(\tau)\theta_3(\tau),\quad
 \theta_4(-1/\tau)=A(\tau)\theta_2(\tau).
\]
这解释了为何 theta 常数与模形式有关，却不能直接把每个常数称作
第二十三章意义下的全模群标量模形式：变换因子出现平方根，而且
\(T\) 与 \(S\) 会交换不同函数。尤其 \(\theta_3(\tau+1)=\theta_4(\tau)\)，
二者的首个非恒定系数分别为 \(2q\) 和 \(-2q\)，所以并不相等。
把函数取幂并作适当组合，才可能消除这些相位和交换；只凭它含有一个
\((-\mathrm i\tau)^{1/2}\) 因子，不能省去其余变换条件。
\ProseParagraphBreak
分数幂的单值性同样可从 \(T\) 看清。虽然 \(\tau\mapsto\tau+2\) 保持
\(q=e^{\pi\mathrm i\tau}\) 不变，却使
\(\theta_2(\tau+2)=\mathrm i\theta_2(\tau)\)，因为
\(q^{1/4}=e^{\pi\mathrm i\tau/4}\) 随参数路径改变。后文的
\(q\partial_q\) 是已在\cref{b3:prop:24-theta-heat}说明的局部微分算子，
不会把这个全局分支信息丢掉。
\ProseParagraphBreak
参数反演还有一个可计算的用途。设 \(t>0\)，由于
\(n^2\ge1+3(n-1)\) 对 \(n\ge1\) 成立，
\[
 0<\theta_3(\mathrm it)-1
       =2\sum_{n\ge1}e^{-\pi tn^2}
       \le\frac{2e^{-\pi t}}{1-e^{-3\pi t}}.
\]
大 \(t\) 时原级数很快收敛，小 \(t\) 时则用
\[
 \theta_3(\mathrm it)=t^{-1/2}\theta_3(\mathrm i/t)
     =t^{-1/2}\left(1+O(e^{-\pi/t})\right)
     \quad(t\downarrow0).
\]
余项直接由刚给出的界把 \(t\) 换成 \(1/t\) 得到。这不只是一个形式对称：
它把收敛慢的参数变成收敛快的参数，同时精确给出主要大小
\(t^{-1/2}\)。从平方和生成函数走向 Mellin 变换时，零端点的这个大小也将决定积分最初的收敛范围。

% !TeX root = ../../Book3.tex
\section{Jacobi 三重积公式}
\label{b3:sec:2402}

级数把不同整数指标的贡献相加，适合研究变换和表示数；乘积把零点逐个列出，适合判断除法是否有意义。这里要把上一节的 theta 级数改写成乘积，并说明为什么无穷多个非零因子不会在极限中另外制造零点。

证明从有限多项式开始。先把有限乘积的系数算清楚，再把指标居中，使指数变成平方，最后在紧集上取极限。这条路线同时给出恒等式中的常数，避免只凭零点相同就猜测两个函数相等。

\subsection{三重积恒等式}
\label{b3:subsec:240201}

下面的有限恒等式起着普通二项式公式的作用。记号中的下标 \(Q\) 表示因子按等比数列变化；这里的系数一般是复数，不借用计数或正性来证明公式。

\Needspace{8\baselineskip}
\begin{lemma}[有限二项式恒等式]\label{b3:lem:24-finite-binomial}
设 \(Q\in\mathbb C\)、\(0<|Q|<1\)，并令
\[
 (Q;Q)_0=1,\qquad (Q;Q)_N=\prod_{j=1}^N(1-Q^j)
 \quad(N\ge1).
\]
对整数 \(N\ge0\) 及 \(0\le r\le N\)，定义
\[
 {N\brack r}_Q
 =\frac{(Q;Q)_N}{(Q;Q)_r(Q;Q)_{N-r}},
\]
并在 \(r<0\) 或 \(r>N\) 时规定 \({N\brack r}_Q=0\)。则对任意 \(t\in\mathbb C\)，
\begin{equation}\label{b3:eq:24-finite-binomial}
 \prod_{j=0}^{N-1}(1+tQ^j)
 =\sum_{r=0}^N Q^{r(r-1)/2}{N\brack r}_Q\,t^r.
\end{equation}
当 \(N=0\) 时，左边的空乘积和右边都等于一。
\end{lemma}
\begin{proof}
因 \(|Q|<1\)，所有分母因子都不为零。先证明递推关系
\[
 {N+1\brack r}_Q
 ={N\brack r}_Q+Q^{N+1-r}{N\brack r-1}_Q.
\]
在 \(1\le r\le N\) 时，把右边通分为分母
\((Q;Q)_r(Q;Q)_{N+1-r}\)，分子为
\[
 (Q;Q)_N\bigl((1-Q^{N+1-r})
              +Q^{N+1-r}(1-Q^r)\bigr)
 =(Q;Q)_{N+1}.
\]
在 \(r=0,N+1\) 时，越界系数为零的约定直接给出同一递推。

现在对 \(N\) 归纳。把 \(N\) 个因子的展开乘以 \(1+tQ^N\)，所得 \(t^r\) 的系数为
\[
 Q^{r(r-1)/2}{N\brack r}_Q
 +Q^N Q^{(r-1)(r-2)/2}{N\brack r-1}_Q.
\]
提出 \(Q^{r(r-1)/2}\) 后，括号中的式子就是上述递推的右边，因此得到 \(N+1\) 的公式。
\end{proof}

例如 \(N=2\) 时，公式是
\[
 (1+t)(1+tQ)=1+(1+Q)t+Qt^2.
\]
额外的因子 \(Q^{r(r-1)/2}\) 不能省去：选择两个因子中的 \(t\) 项时，所得系数是 \(Q\)，而不是一。

将有限公式居中后，就得到下面的\term[jacobi-sanchongji]{Jacobi 三重积公式}[Jacobi triple product identity]。

\Needspace{8\baselineskip}
\begin{theorem}[Jacobi 三重积]\label{b3:thm:24-triple-product}
设 \(q,w\in\mathbb C\)，且 \(0<|q|<1\)、\(w\ne0\)。则
\[
 \sum_{n\in\mathbb Z}q^{n^2}w^n
 =\prod_{m=1}^{\infty}
   (1-q^{2m})(1+wq^{2m-1})(1+w^{-1}q^{2m-1}).
\]
固定 \(q\) 后，级数绝对正规收敛，乘积正规收敛，二者都在
\(\mathbb C\setminus\{0\}\) 上定义关于 \(w\) 的全纯函数。
\end{theorem}
\begin{proof}
在式\eqref{b3:eq:24-finite-binomial}中，把因子个数取为 \(2N\)，再取
\[
 Q=q^2,\qquad t=wq^{-(2N-1)}.
\]
左边各因子中 \(q\) 的指数依次为
\(-(2N-1),-(2N-3),\ldots,2N-1\)。对负指数的那一半使用
\[
 1+wq^{-(2j-1)}
 =wq^{-(2j-1)}(1+w^{-1}q^{2j-1}).
\]
提出的全部单项式相乘为 \(w^Nq^{-N^2}\)，因为前 \(N\) 个正奇数的和为 \(N^2\)。

右边第 \(\ell\) 项含有
\[
 q^{\ell(\ell-1)}w^\ell q^{-(2N-1)\ell}
 =w^\ell q^{\ell^2-2N\ell}.
\]
因此两边乘以 \(w^{-N}q^{N^2}\)，并令 \(r=\ell-N\)，就得到
\begin{equation}\label{b3:eq:24-centered-product}
 \prod_{j=1}^N(1+wq^{2j-1})(1+w^{-1}q^{2j-1})
 =\sum_{r=-N}^N q^{r^2}w^r{2N\brack N+r}_{q^2}.
\end{equation}
平方指数来自 \(\ell^2-2N\ell+N^2=(\ell-N)^2\)，而双边指标来自把长度为 \(2N\) 的有限区间移到原点两侧。

接着说明极限为什么可以进入右边的和。由
\(\sum_{j\ge1}|q|^{2j}<\infty\) 和\cref{b3:prop:20-absolute-product}的非零乘积判据，
\[
 A=\prod_{j=1}^{\infty}(1-q^{2j})
\]
存在且非零。对每个固定整数 \(r\)，当 \(N\to\infty\) 时，商式中的三个有限乘积都趋于 \(A\)，故
\[
 {2N\brack N+r}_{q^2}\longrightarrow A^{-1}.
\]
这里还需要一个不随 \(N,r\) 变化的界。记
\[
 d=\prod_{j=1}^{\infty}(1-|q|^{2j})>0,\qquad
 M=\prod_{j=1}^{\infty}(1+|q|^{2j})<\infty.
\]
任一分母有限乘积的模不小于 \(d\)，分子有限乘积的模不大于 \(M\)，所以
\[
 \left|{2N\brack N+r}_{q^2}\right|
 \le Md^{-2}\qquad(-N\le r\le N).
\]

固定紧圆环 \(a\le|w|\le b\)，其中 \(0<a\le b<\infty\)，令
\(B=\max\{b,a^{-1},1\}\)。把区间外的系数补为零后，右边每一项的模都不超过
\[
 Md^{-2}|q|^{r^2}B^{|r|},\qquad r\in\mathbb Z.
\]
这个双边数列可和，因为平方衰减压过固定的指数增长。于是先把尾和统一压小，再对剩余有限项取极限，得到右边在紧圆环上一致趋于
\(A^{-1}\sum_{r\in\mathbb Z}q^{r^2}w^r\)。

左边也在紧圆环上一致收敛：两个因子与一的差分别被
\(b|q|^{2j-1}\) 和 \(a^{-1}|q|^{2j-1}\) 控制，控制数列均可和。
\cref{b3:prop:20-normal-product}因而适用。对式\eqref{b3:eq:24-centered-product}取极限，再乘以 \(A\)，即得所需公式及收敛结论。
\end{proof}

固定 \(r\) 的系数极限并不足以独自证明定理，因为求和区间还在扩大。证明中的共同上界专门控制那些随 \(N\) 向两端移动的项。另一方面，三重积的某个因子可以为零；正规收敛所保证的是充分远的尾部非零，并没有排除有限头部给出的零点。

\subsection{乘积展开}
\label{b3:subsec:240202}

现在恢复 \(\tau\in\mathfrak H\)、\(q=e^{\pi\mathrm i\tau}\)。本章的空间变量是 \(z\)，因此周期中出现 \(\pi\)，而不是一。以下所有
\(q^{1/4}\) 都表示 \(e^{\pi\mathrm i\tau/4}\)；把 \(q\) 单独看成穿孔圆盘中的变量时，不能另选一个全局的四次根来替换它。

\Needspace{8\baselineskip}
\begin{proposition}[四种 theta 函数的乘积]\label{b3:prop:24-theta-products}
设 \(\tau\in\mathfrak H\)、\(q=e^{\pi\mathrm i\tau}\)，并以
\(q^{1/4}=e^{\pi\mathrm i\tau/4}\) 规定分数幂。则上一节定义的四种 theta 函数对每个 \(z\in\mathbb C\) 满足
\begin{align}
 \theta_1(z,\tau)
 &=2q^{1/4}\sin z
   \prod_{n\ge1}(1-q^{2n})
   (1-2q^{2n}\cos2z+q^{4n}),
   \label{b3:eq:24-theta1-product}\\
 \theta_2(z,\tau)
 &=2q^{1/4}\cos z
   \prod_{n\ge1}(1-q^{2n})
   (1+2q^{2n}\cos2z+q^{4n}),
   \label{b3:eq:24-theta2-product}\\
 \theta_3(z,\tau)
 &=\prod_{n\ge1}(1-q^{2n})
   (1+2q^{2n-1}\cos2z+q^{4n-2}),
   \label{b3:eq:24-theta3-product}\\
 \theta_4(z,\tau)
 &=\prod_{n\ge1}(1-q^{2n})
   (1-2q^{2n-1}\cos2z+q^{4n-2}).
   \label{b3:eq:24-theta4-product}
\end{align}
固定 \(\tau\) 后，乘积关于 \(z\) 正规收敛，并可对 \(z\) 作任意固定阶数的逐项微分。其收敛在 \((z,\tau)\) 的紧集上也一致。
\end{proposition}
\begin{proof}
在三重积中分别取 \(w=e^{2\mathrm iz}\) 和 \(w=-e^{2\mathrm iz}\)，得到 \(\theta_3,\theta_4\) 的式子。其中两个含 \(w\) 的因子相乘，分别给出正号和负号的余弦项。

对半整数指标，先在级数中提取共同因子：
\[
 \begin{aligned}
 \theta_2(z,\tau)
 &=q^{1/4}e^{\mathrm iz}
   \sum_{n\in\mathbb Z}q^{n^2}(qe^{2\mathrm iz})^n,\\
 \theta_1(z,\tau)
 &=-\mathrm i q^{1/4}e^{\mathrm iz}
   \sum_{n\in\mathbb Z}q^{n^2}(-qe^{2\mathrm iz})^n.
 \end{aligned}
\]
第一行的两个指数因子成为
\(1+q^{2m}e^{2\mathrm iz}\) 和
\(1+q^{2m-2}e^{-2\mathrm iz}\)。
后一个乘积在 \(m=1\) 时提供 \(1+e^{-2\mathrm iz}\)；将剩余指标减一，便与第一个乘积配对。外面的因子化为
\[
 e^{\mathrm iz}(1+e^{-2\mathrm iz})=2\cos z.
\]
第二行同样给出负号的两个指数因子，而外面的因子为
\[
 -\mathrm i e^{\mathrm iz}(1-e^{-2\mathrm iz})=2\sin z.
\]
这就解释了四个公式中不同的首因子及符号。

若 \((z,\tau)\) 限制在一个紧集，则有 \(|q|\le r<1\) 及
\(|e^{\pm2\mathrm iz}|\le C\)。把每个二次因子拆为
\(1\pm q^m e^{2\mathrm iz}\) 与 \(1\pm q^m e^{-2\mathrm iz}\)，其偏离一的量被 \(Cr^m\) 控制。故乘积的正规收敛及有限分组均由\cref{b3:prop:20-normal-product}保证。固定 \(\tau\) 后，在稍大的 \(z\) 紧集上使用同一估计及 Cauchy 公式，就得到各阶导数的局部一致收敛。
\end{proof}

乘积表示给出的不只是一些显然存在的零点，还包括所有零点及其重数。后文把 theta 函数放在分母时，将用到这一完整信息。

\Needspace{8\baselineskip}
\begin{proposition}[theta 零点及常数非零性]\label{b3:prop:24-theta-zeros}
设 \(\tau\in\mathfrak H\)，令
\(\Lambda=\pi\mathbb Z+\pi\tau\mathbb Z\)。四个整函数
\(z\mapsto\theta_j(z,\tau)\) 的零点集分别为
\[
 \begin{aligned}
 Z(\theta_1)&=\Lambda,
 &Z(\theta_2)&=\frac{\pi}{2}+\Lambda,\\
 Z(\theta_3)&=\frac{\pi+\pi\tau}{2}+\Lambda,
 &Z(\theta_4)&=\frac{\pi\tau}{2}+\Lambda.
 \end{aligned}
\]
其中每一个零点都是单零点。特别地，
\(\theta_2(0,\tau)\)、\(\theta_3(0,\tau)\)、\(\theta_4(0,\tau)\)
均不为零，而 \(\theta_1(0,\tau)=0\) 且
\(\partial_z\theta_1(0,\tau)\ne0\)。
\end{proposition}
\begin{proof}
在 \(\theta_1\) 的乘积中，因子 \(1-q^{2n}\) 从不为零。
首因子 \(\sin z\) 的零点是 \(\pi\mathbb Z\)，且都单重。另外两个因子为
\[
 1-q^{2n}e^{2\mathrm iz},\qquad
 1-q^{2n}e^{-2\mathrm iz}.
\]
它们的零点分别是 \(-n\pi\tau+\pi\mathbb Z\) 和
\(n\pi\tau+\pi\mathbb Z\)。在各自的零点处，导数为
\(-2\mathrm i\) 或 \(2\mathrm i\)，所以也都是单零点。
因 \(\operatorname{Im}\tau>0\)，这些水平直线彼此不同，且不与
\(\sin z\) 的零点重叠。于是每个格点恰好来自一个单零因子。

对 \(\theta_2\)，首因子换成 \(\cos z\)，两个指数因子换成正号，零点的实部相应平移 \(\pi/2\)，得到 \(\pi/2+\Lambda\)。
对 \(\theta_3,\theta_4\)，指数为 \(2n-1\)，所以虚方向的平移为
\(\pm(n-\tfrac12)\pi\tau\)。正号因子再增加 \(\pi/2\) 的水平平移，负号因子不增加，分别得到所列的两个集合。各因子的导数仍非零，不同因子的零点仍不重叠。

最后，正规乘积在每个小邻域都可写成有限头部乘一个非零的全纯尾部；这是\cref{b3:prop:20-product-zero-orders}的具体适用情形。因此没有遗漏的额外零点，也没有因无限乘积而增加的重数。原点不属于后三个格的平移：比较虚部先确定 \(\tau\) 的系数，再比较实部即可验证。常数非零性随之成立。
\end{proof}

四个零点集在一个格胞中分别占据四个半周期位置。准周期平移会改变函数值的非零乘子，但不会改变零点的重数；因此这里描述的是整张平面上的零点排列，而不是只在一个基本区域内的局部结论。

\subsection{特殊值}
\label{b3:subsec:240203}

令 \(z=0\)。为减轻记号，以下把 \(\theta_j(0,\tau)\) 简写成
\(\theta_j\)，但撇号始终表示对空间变量 \(z\) 求导。由乘积公式，
\[
 \begin{aligned}
 \theta_2&=2q^{1/4}
       \prod_{n\ge1}(1-q^{2n})(1+q^{2n})^2,\\
 \theta_3&=\prod_{n\ge1}(1-q^{2n})(1+q^{2n-1})^2,\\
 \theta_4&=\prod_{n\ge1}(1-q^{2n})(1-q^{2n-1})^2.
 \end{aligned}
\]
这些是关于参数 \(\tau\) 的函数；“常数”只表示空间变量已经取零。

三个非零常数还决定奇函数在原点的导数，这一关系称为\term[jacobi-daoshu-hengdengshi]{Jacobi 导数恒等式}[Jacobi derivative identity]。

\Needspace{6\baselineskip}
\begin{proposition}[Jacobi 导数恒等式]\label{b3:prop:24-jacobi-derivative}
设 \(\tau\in\mathfrak H\)，四个 theta 函数采用本章的级数归一化，则
\begin{equation}\label{b3:eq:24-derivative-identity}
 \partial_z\theta_1(0,\tau)
 =\theta_2(0,\tau)\theta_3(0,\tau)\theta_4(0,\tau).
\end{equation}
两边均不为零。
\end{proposition}
\begin{proof}
令 \(q=e^{\pi\mathrm i\tau}\)，并规定
\(q^{1/4}=e^{\pi\mathrm i\tau/4}\)。在式\eqref{b3:eq:24-theta1-product}中除以 \(z\) 并令 \(z\to0\)，由于乘积局部一致收敛及 \(\sin z/z\to1\)，得到
\[
 \theta_1'(0,\tau)
 =2q^{1/4}\prod_{n\ge1}(1-q^{2n})^3.
\]
现在将三个 theta 常数的乘积式相乘。共同因子已经给出右边这个式子，剩下的因子是
\[
 \left(\prod_{n\ge1}
    (1+q^{2n})(1-q^{4n-2})\right)^2.
\]
记 \(A=\prod_{n\ge1}(1-q^{2n})\ne0\)。由
\(1+q^{2n}=(1-q^{4n})/(1-q^{2n})\)，括号内的乘积等于
\[
 \frac{\displaystyle
        \prod_{n\ge1}(1-q^{4n})
        \prod_{n\ge1}(1-q^{4n-2})}
      {\displaystyle\prod_{n\ge1}(1-q^{2n})}=1.
\]
分子两列指数 \(4,8,12,\ldots\) 与 \(2,6,10,\ldots\) 正好分割全部正偶数。这里的分组和约分合法，因为各乘积绝对收敛且非零；并没有把两个发散乘积相除。恒等式随之成立，非零性也从这些乘积直接得到。
\end{proof}

这个归一化没有额外的 \(\pi\) 因子，因为本章级数含
\(e^{2\mathrm inz}\)，基本水平准周期为 \(\pi\)。若改用周期为一的空间坐标，空间求导的换元因子也必须随之改变。

当 \(\tau=\mathrm it\)、\(t>0\) 时，\(q=e^{-\pi t}\in(0,1)\)。
三个常数乘积中的全部因子都是正数，因此
\(\theta_2,\theta_3,\theta_4>0\)。尤其是交替级数
\[
 \theta_4=1-2q+2q^4-2q^9+\cdots
\]
始终为正。单看前两项，在 \(q>1/2\) 时容易作出错误判断；乘积同时控制了被截去的无限尾部。对一般复参数，仍有非零性，但不再谈正负。

在 \(q\to0\) 时，级数还给出
\[
 \begin{aligned}
 \theta_2&=2q^{1/4}\bigl(1+q^2+O(q^6)\bigr),\\
 \theta_3&=1+2q+O(q^4),\qquad
 \theta_4=1-2q+O(q^4).
 \end{aligned}
\]
因此它们的乘积为 \(2q^{1/4}(1-3q^2+O(q^4))\)，与导数乘积的展开相符。这个检查能辨认归一化和低次系数，却不能代替对所有 \(\tau\in\mathfrak H\) 成立的恒等式证明。

\subsection{三重积的系数、零点与对数导数}
\label{b3:subsec:240204}

三重积允许比较系数，但要先说清楚比较的是哪个变量。固定
\(w\ne0\) 后，三重积两边关于 \(q\) 在 \(|q|<1\) 全纯，并在
\(q=0\) 都取值一；局部一致收敛保证可以使用幂级数的唯一性。
固定 \(q\) 后，关于 \(w\) 则是穿孔平面上的 Laurent 展开。两种比较都依赖解析收敛。靠近 \(|q|=1\) 时，上面的共同控制不再一致，不能未经新的估计就把 \(q\) 代成单位根。

还有一种需要格外区分的操作：对乘积求对数导数。含有零点的整函数没有全平面上的全纯对数，但在每个不含零点的小圆盘中都有局部对数；不同分支只差常数，所以对数的导数不依赖分支。

\Needspace{8\baselineskip}
\begin{proposition}[theta 的对数导数]\label{b3:prop:24-theta-log-derivative}
设 \(\tau\in\mathfrak H\)、\(q=e^{\pi\mathrm i\tau}\)，令
\(\Lambda=\pi\mathbb Z+\pi\tau\mathbb Z\) 及
\[
 L(z)=\frac{\partial_z\theta_1(z,\tau)}{\theta_1(z,\tau)}
 \qquad(z\in\mathbb C\setminus\Lambda).
\]
则
\begin{equation}\label{b3:eq:24-theta-logarithmic-series}
 L(z)=\cot z+
 4\sin(2z)\sum_{n\ge1}
 \frac{q^{2n}}{1-2q^{2n}\cos(2z)+q^{4n}}.
\end{equation}
右边的级数在 \(\mathbb C\setminus\Lambda\) 的紧集上绝对一致收敛，并可作任意固定阶数的逐项微分。在更小的范围
\(|\operatorname{Im}z|<\pi\operatorname{Im}\tau\)、\(z\notin\pi\mathbb Z\) 内，还有
\begin{equation}\label{b3:eq:24-theta-log-fourier}
 L(z)=\cot z+
 4\sum_{k\ge1}\frac{q^{2k}}{1-q^{2k}}\sin(2kz).
\end{equation}
函数 \(-L'\) 是以 \(\pi,\pi\tau\) 为周期的亚纯函数，其全部极点是
\(\Lambda\)，在每个 \(\lambda\in\Lambda\) 处的主部都是
\((z-\lambda)^{-2}\)。
\end{proposition}
\begin{proof}
在避开 \(\Lambda\) 的紧集上，theta 乘积的有限头部非零，尾因子一致接近一。把紧集稍微放大后，尾因子的偏差仍由几何级数控制。因此\cref{b3:prop:20-product-log-derivative}的条件成立：可对局部对数求导，并可继续逐项微分。由
\[
 \frac{\mathrm d}{\mathrm dz}
   \log(1-2q^{2n}\cos2z+q^{4n})
 =\frac{4q^{2n}\sin2z}
        {1-2q^{2n}\cos2z+q^{4n}},
\]
以及首因子 \(\sin z\) 的对数导数，得到第一式。这里的公式只要求避开零点，不要求 \(z\) 留在一个水平带内。

第二式需要额外的几何级数展开。若
\(|\operatorname{Im}z|<-\log|q|=\pi\operatorname{Im}\tau\)，则对每个 \(n\ge1\)，
\(|q^{2n}e^{\pm2\mathrm iz}|<1\)。将两个指数因子的对数导数分别展开后相加，得到
\[
 \frac{-2\mathrm i q^{2n}e^{2\mathrm iz}}
           {1-q^{2n}e^{2\mathrm iz}}
 +\frac{2\mathrm i q^{2n}e^{-2\mathrm iz}}
           {1-q^{2n}e^{-2\mathrm iz}}
 =4\sum_{k\ge1}q^{2nk}\sin(2kz).
\]
在闭子带 \(|\operatorname{Im}z|\le R<-\log|q|\) 上，相应双重级数的绝对值和不超过常数乘以
\[
 \sum_{k\ge1}
   \frac{|q|^{2k}e^{2kR}}{1-|q|^{2k}}<\infty.
\]
故可交换 \(n,k\) 的求和顺序，先对 \(n\) 求和，就得到式\eqref{b3:eq:24-theta-log-fourier}。

由上一节的准周期公式，
\[
 \theta_1(z+\pi,\tau)=-\theta_1(z,\tau),\qquad
 \theta_1(z+\pi\tau,\tau)
 =-q^{-1}e^{-2\mathrm iz}\theta_1(z,\tau).
\]
取对数导数即得 \(L(z+\pi)=L(z)\) 和
\(L(z+\pi\tau)=L(z)-2\mathrm i\)，再求导便消去加法常数。
另一方面，\cref{b3:prop:24-theta-zeros}保证每个格点均为单零点。
在 \(\lambda\) 附近可写
\(\theta_1(z,\tau)=(z-\lambda)h(z)\)，其中 \(h\) 全纯且非零。于是
\[
 L(z)=\frac1{z-\lambda}+\frac{h'(z)}{h(z)},\qquad
 -L'(z)=\frac1{(z-\lambda)^2}
          -\left(\frac{h'}h\right)'(z).
\]
第二项全纯，从而得到全部极点和主部。
\end{proof}

例如，在原点附近对式\eqref{b3:eq:24-theta-log-fourier}求导，并使用
\(\csc^2z=z^{-2}+\tfrac13+O(z^2)\)，得到
\begin{equation}\label{b3:eq:24-theta-local-part}
 -L'(z)=\frac1{z^2}+\frac13
       -8\sum_{k\ge1}\frac{kq^{2k}}{1-q^{2k}}+O(z^2).
\end{equation}
这里极点的系数一已经由单零点确定，常数项则由乘积归一化决定。后文与同一格的 \(\wp\) 函数比较时，二者之差只能剩下常数；不能只因二者都双周期就认定它们相等。

参数方向的对数微分也要有共同控制。例如
\[
 B(q)=\prod_{n\ge1}(1-q^{2n})\qquad(|q|<1)
\]
处处非零且全纯，并满足
\[
 q\frac{B'(q)}{B(q)}
 =-\sum_{n\ge1}\frac{2nq^{2n}}{1-q^{2n}}.
\]
在 \(|q|\le r<1\) 上，右边由
\(\sum_{n\ge1}2nr^{2n}/(1-r^{2n})<\infty\) 控制。
其他 theta 常数乘积的参数导数有同样的几何衰减。若出现
\(q^{1/4}\)，则在 \(\tau\) 平面使用
\[
 q\partial_q=\frac1{\pi\mathrm i}\partial_\tau,
 \qquad q\partial_q\log(q^{1/4})=\frac14.
\]
后一式是局部对数的导数陈述，不要求在整个穿孔 \(q\) 圆盘选定单值的 \(\log q\)。这些估计将使后文从乘积转到约数和的逐项运算有明确的适用范围。

% !TeX root = ../../Book3.tex
\section{\texorpdfstring{\(\theta\)}{\AnalysisPdfTheta} 级数作为生成函数}
\label{b3:sec:2403}

系数计数必须先约定是否区分顺序和符号。theta 乘法自动区分各坐标，而且正负整数分别出现，因此它计算的是有序带符号表示数。

\subsection{生成函数}
\label{b3:subsec:240301}

设 \(r_k(n)=\#\{(a_1,\ldots,a_k)\in\mathbb Z^k:\sum a_j^2=n\}\)。每个坐标绝对值不超过 \(\sqrt n\)，故数目有限。对 \(|q|<1\)，绝对收敛保证
\[
\theta_3(0,\tau)^k=\sum_{n\ge0}r_k(n)q^n.
\]
此处 \(r_k(0)=1\)。例如 \(r_2(1)=4\)、\(r_4(1)=8\)，分别选择非零坐标的位置和它的符号，不能把它们解释为不计顺序的分拆数。

\subsection{乘积与卷积}
\label{b3:subsec:240302}

乘法给 \(r_{k+\ell}(n)=\sum_{m=0}^n r_k(m)r_\ell(n-m)\)。这是把前 \(k\) 个坐标与后 \(\ell\) 个坐标分别计数的精确分解。

\ProseParagraphBreak

若改用 \(\theta_4\)，每个坐标带符号 \((-1)^{a_j}\)。由于 \(a_j^2\equiv a_j\pmod2\)，同一层 \(\sum a_j^2=n\) 上总符号恒为 \((-1)^n\)。因此 \(\theta_4^k=\sum(-1)^nr_k(n)q^n\)。它筛选奇偶层，却不单独筛选每个坐标的奇偶模式。

\subsection{系数比较}
\label{b3:subsec:240303}

在四个变量中还有一个有用的精确变换。对 \(a=(a_1,a_2,a_3,a_4)\) 用正交矩阵
\[
\frac12\begin{pmatrix}1&1&1&1\\1&1&-1&-1\\1&-1&1&-1\\1&-1&-1&1\end{pmatrix}
\]
变换为 \(b\)。当 \(\sum a_j\) 为奇数时，所有 \(b_j\in\mathbb Z+1/2\)；反向变换相同。半整数四元组又分为两类：写 \(b_j=m_j+1/2\)，按 \(\sum m_j\) 的奇偶分。只有其中一类在逆变换下给整数向量；将一个 \(b_j\) 改号交换两类且保持平方和。故全部半整数向量的计数恰为奇数平方和整数向量计数的两倍。比较级数得到
\begin{equation}\label{b3:eq:24-quartic-identity}
\theta_3^4=\theta_4^4+\theta_2^4.
\end{equation}
这个推导既说明常数二从何而来，也说明为什么四次方出现。它不是把前几项相同当成无穷恒等式，而是给出保持范数的对应。

\ProseParagraphBreak

生成函数还可以筛选同余条件。若 \(A(q)=\sum a_nq^n\)，\(\zeta_m=e^{2\pi\mathrm{i}/m}\)，则有限 Fourier 正交关系给
\[
 \frac1m\sum_{j=0}^{m-1}\zeta_m^{-bj}A(\zeta_m^j q)
 =\sum_{n\equiv b\ (m)}a_nq^n.
\]
交换有限和与级数不需要新的收敛定理；只要 \(|q|<1\)，每个旋转参数仍在原收敛圆盘内。对平方和而言，这个操作筛选总平方和的剩余类。若要每个坐标分别满足某些同余条件，则须先对每个一维 theta 因子筛选，再相乘。两个程序统计的集合不同，不能仅凭最后都含若干符号因子就互换。
\ProseParagraphBreak
四维正交变换中，半整数层的两类也可直接检查。写 \(b_j=m_j+1/2\)，逆变换的第一坐标为 \((\sum m_j+2)/2\)，其余坐标具有同样的整数奇偶条件。因此恰在 \(\sum m_j\) 偶数时，逆像为整数四元组；这时逆像的坐标和为 \(2b_1\)，必为奇数。把一个 \(b_j\) 改号，相应的整数 \(m_j\) 变成 \(-m_j-1\)，总和奇偶翻转，所以两类等势且逐层保范数。这样便把前述系数比较中的二倍关系落实为显式双射，而非只比较总的无穷集合大小。
\ProseParagraphBreak
平均表示数可以从几何体积作一个独立的数量级检验。令 \(A_k(R)\) 为 \(\mathbb Z^k\) 中长度不超过 \(R\) 的点数。以每个格点为中心放置边长一的闭立方体，它们内部不交，每个体积一。所取立方体的并包含半径 \(R-\sqrt{k}/2\) 的球，又包含在半径 \(R+\sqrt{k}/2\) 的球内。因此
\[
 A_k(R)=V_kR^k+O_k(R^{k-1}),
\]
其中 \(V_k\) 为单位球体积。这个估计不控制单个球壳上的点数，却控制从0累计到 \(R^2\) 的表示数。令 \(X=R^2\)，即为 \(\sum_{n\le X}r_k(n)=V_kX^{k/2}+O_k(X^{(k-1)/2})\)。平方和问题中，逐项精确公式与累计体积估计是不同层次的结果。

% !TeX root = ../../Book3.tex
\section{二平方和与四平方和}
\label{b3:sec:2404}

二平方和可以从 Gaussian 整数的唯一分解得到；四平方和的精确计数则可利用 theta 乘积的对数导数。后者需要一个半周期差值恒等式，下面将它从前章的椭圆函数推出来。

\subsection{二平方和定理}
\label{b3:subsec:240401}

\begin{theorem}[二平方和]\label{b3:thm:24-two-squares}
正整数 \(n\) 可表示为两个整数平方之和，当且仅当每个 \(p\equiv3\pmod4\) 的素因子在 \(n\) 中的指数为偶数。计数公式为
\(r_2(n)=4\sum_{d\mid n}\chi_4(d)\)，其中 \(\chi_4(d)=0\) 对偶数，\(\chi_4(d)=(-1)^{(d-1)/2}\) 对奇数。
\end{theorem}
\begin{proof}
先说明所需代数。Gaussian 整数 \(\mathbb Z[\mathrm{i}]\) 的范数为 \(a^2+b^2\)。对商的实虚部各取最近整数，可使除法余项的范数至多为除数范数的一半，因此存在 Euclid 算法、最大公因子和 Bézout 恒等式。若不可约元整除积，Bézout 恒等式证明它整除某个因子，继而按范数归纳得到唯一分解。
有理素数 \(p\equiv3\pmod4\) 仍不可约：若分解，其因子范数只能为 \(p\)，但平方模四不可能相加得三。对 \(p\equiv1\pmod4\)，Wilson 配对证明 \((p-1)!\equiv-1\pmod p\)，又有 \((p-1)!\equiv(-1)^{(p-1)/2}(((p-1)/2)!)^2\)，所以存在 \(u^2\equiv-1\pmod p\)。若 \(p\) 在 Gaussian 整数中不可约，它整除 \((u+\mathrm{i})(u-\mathrm{i})\) 就会整除一因子，这在虚部为 \(\pm1\) 时不可能。于是 \(p=\pi\overline\pi\)，其中 \(N\pi=p\)，两因子不相伴。素数二则为 \(-\mathrm{i}(1+\mathrm{i})^2\)。
设 \(n=2^a\prod_{p\equiv1(4)}p^{a_p}\prod_{q\equiv3(4)}q^{b_q}\)。若 \(N\alpha=n\)，每个惰性素数 \(q\) 在范数中的指数必须偶；反之这些指数偶时，可以给 \(\alpha\) 选 \(q^{b_q/2}\)、\((1+\mathrm{i})^a\)，并对每个分裂素数在 \(\pi\) 与 \(\overline\pi\) 间分配共 \(a_p\) 个因子。因此表示数为 \(4\prod(a_p+1)\)，前面的四为四个单位。约数和 \(\sum_{d\mid n}\chi_4(d)\) 按素数幂相乘：二贡献一，分裂素数贡献 \(a_p+1\)，惰性素数贡献一或零，恰好得到公式。
\end{proof}

\subsection{Jacobi 四平方和定理}
\label{b3:subsec:240402}

\begin{theorem}[Jacobi 四平方和]\label{b3:thm:24-four-squares}
对每个正整数 \(n\)，有
\[
r_4(n)=8\sum_{d\mid n,\ 4\nmid d}d.
\]
特别地，每个正整数均为四个整数平方之和。
\end{theorem}
证明将在紧接的两小节完成。这里最后一句来自约数一总在求和范围内，故表示数严格为正；不能把所有项为非负误当成已知总和不为零。该定理称为\term[jacobi-sipingfanghe]{Jacobi 四平方和定理}[Jacobi four-square theorem]。

\subsection{\texorpdfstring{\(\theta\)}{\AnalysisPdfTheta} 函数证明}
\label{b3:subsec:240403}

取格 \(\Lambda=\pi\mathbb Z+\pi\tau\mathbb Z\)，记 \(e_1=\wp(\pi/2)\)、\(e_3=\wp(\pi\tau/2)\)。由准周期性，\(-\mathrm{d}^2\log\theta_1(z)/\mathrm{d}z^2\) 双周期；在格点主部为 \(z^{-2}\)。与 \(\wp\) 相减是整椭圆函数，所以
\(\wp(z)=-(\log\theta_1(z))''+C\)。平移级数得 \(\theta_1(z+\pi/2)=\theta_2(z)\)，以及
\(\theta_1(z+\pi\tau/2)=\mathrm{i}q^{-1/4}e^{-\mathrm{i}z}\theta_4(z)\)。三个偶函数在原点导数为零，因此
\[
e_1-e_3=\frac{\theta_4''(0)}{\theta_4}-\frac{\theta_2''(0)}{\theta_2}.
\]
另一方面，函数
\(F(z)=(\theta_1'(0)/\theta_2)^2(\theta_2(z)/\theta_1(z))^2\) 双周期，在原点主部为 \(z^{-2}\)，并在 \(\pi/2\) 有二重零点；故 \(F=\wp-e_1\)。再代入 \(z=\pi\tau/2\)，利用
\(\theta_2(z+\pi\tau/2)=q^{-1/4}e^{-\mathrm{i}z}\theta_3(z)\) 和\eqref{b3:eq:24-derivative-identity}，得到 \(e_3-e_1=-\theta_3^4\)。于是
\[
\frac{\theta_4''(0)}{\theta_4}-\frac{\theta_2''(0)}{\theta_2}=\theta_3^4.
\]
每个 theta 级数还满足 \(q\partial_q\theta_j=-\tfrac14\partial_z^2\theta_j\)，直接逐项微分即可核对。因此
\begin{equation}\label{b3:eq:24-log-derivative}
4q\frac{\mathrm{d}}{\mathrm{d}q}\log(\theta_2/\theta_4)=\theta_3^4.
\end{equation}
整个推导中的除法由三重积保证分母非零，椭圆函数的比较则只比较同一格下的主部和一个零点，故不存在未确定的比例常数。

\subsection{生成函数方法}
\label{b3:subsec:240404}

\begin{proof}[\cref{b3:thm:24-four-squares}的证明]
对 \(\theta_2/\theta_4\) 的乘积取对数导数，由\eqref{b3:eq:24-log-derivative}得
\[
\theta_3^4=1+16\sum_{n\ge1}\frac{nq^{2n}}{1+q^{2n}}
+8\sum_{m\ge1,\ m\text{ odd}}\frac{mq^m}{1-q^m}.
\]
恒等式 \(x/(1+x)=x/(1-x)-2x^2/(1-x^2)\) 将偶数部分改写为
\(8\sum_{d\text{ even}}d q^d/(1-q^d)-32\sum_{n\ge1}nq^{4n}/(1-q^{4n})\)。与奇数部分相合，即为
\[
\theta_3^4=1+8\sum_{d\ge1,\ 4\nmid d}\frac{dq^d}{1-q^d}.
\]
在 \(|q|\le r<1\) 上所有对数导数级数和几何展开都绝对一致收敛，因为 \(\sum d r^d/(1-r^d)<\infty\)。故可把 \(q^d/(1-q^d)=\sum_{m\ge1}q^{dm}\) 代入并收集系数，得到所述约数和。最后与\secref{b3:sec:2403}中的生成函数比较。
\end{proof}
例如奇素数 \(p\) 的四平方和表示数为 \(8(p+1)\)；\(n=4\) 时为 \(8(1+2)=24\)，包括八个仅一个坐标为 \(\pm2\) 的表示与十六个四坐标均为 \(\pm1\) 的表示。这个核算检验了符号、顺序及排除四倍约数的三个约定。

\ProseParagraphBreak

半周期比较中的主部与常数值得再逐点核验。由于 \(\theta_1\) 在0单零，\(\theta_1(z)=\theta_1'(0)z+O(z^3)\)；偶函数 \(\theta_2(z)=\theta_2+O(z^2)\)。因此
\((\theta_1'(0)/\theta_2)^2(\theta_2(z)/\theta_1(z))^2=z^{-2}+O(1)\)，没有一阶极项。平方后的商在两周期下的符号和指数相消，确为椭圆函数。它与 \(\wp\) 的差因主部相同而整，故为常数；在 \(z=\pi/2\) 时分子为 \(\theta_1(\pi)=0\)，分母为 \(\theta_2\ne0\)，常数恰为 \(-e_1\)。这一步同时用到了零点的位置、单重性及常数非零，三者均由前面的级数或乘积给出。
\ProseParagraphBreak
在另一半周期，平移公式逐项给
\(\theta_2(\pi\tau/2)=q^{-1/4}\theta_3\) 与
\(\theta_1(\pi\tau/2)=\mathrm{i}q^{-1/4}\theta_4\)。代入平方商，先得到比值平方 \(-\theta_3^2/\theta_4^2\)，再乘
\((\theta_1'(0)/\theta_2)^2=\theta_3^2\theta_4^2\)，最后才得到 \(-\theta_3^4\)。负号来自 \(\mathrm{i}^2=-1\)，不是法向或周期方向的约定；四次方来自两个平方因子的相乘。明确这些步骤后，热方程的负四分之一因子也与最后的正号一致。
\ProseParagraphBreak
从乘积到 Lambert 级数同样可逐类追踪。\(\theta_2/\theta_4\) 中公共的 \(1-q^{2n}\) 全部消去，剩下常数 \(2q^{1/4}\)、偶幂的 \((1+q^{2n})^2\) 和奇幂的 \((1-q^{2n-1})^{-2}\)。乘上 \(4q\mathrm{d}/\mathrm{d}q\) 后，第一项给一，偶幂项给 \(16nq^{2n}/(1+q^{2n})\)，奇幂项给 \(8(2n-1)q^{2n-1}/(1-q^{2n-1})\)。把偶数项改写后，其中全部四的倍数恰被减掉一次。若忘记平方因子或 \(q^{1/4}\)，会分别造成表示数少一倍或漏掉 \(r_4(0)=1\)；前几项计数可以迅速检查这些错误。
\ProseParagraphBreak
例如 \(r_4(2)=24\)：选择两个非零位置有六种，各自符号有四种。\(r_4(3)=32\)：选择零坐标有四种，其余三个符号有八种。对 \(n=9\)，一类表示为 \((\pm3,0,0,0)\) 的排列，共八个；另一类是 \((\pm2,\pm2,\pm1,0)\)，零位置四种、绝对值一的位置三种、三个非零符号八种，共九十六个。总数一百零四，正好是 \(8(1+3+9)\)。相同绝对值坐标的置换不能再重复乘一个阶乘，否则会过计。
\ProseParagraphBreak
四平方公式还能给比体积法更细的累计误差。设 \(X\ge1\)，交换有限约数和得到
\[
 \sum_{1\le n\le X}r_4(n)=8\sum_{m\le X}\sum_{d\le X/m,\,4\nmid d}d.
\]
由 \(\sum_{d\le Y}d=Y^2/2+O(Y)\)，扣去四的倍数后内和为 \(3Y^2/8+O(Y)\)。于是总和为
\[
 3X^2\sum_{m\le X}m^{-2}+O\left(X\sum_{m\le X}m^{-1}\right)
 =\frac{\pi^2}{2}X^2+O(X\log(X+2)).
\]
最后使用正弦乘积比较所得 \(\sum m^{-2}=\pi^2/6\)，尾和为 \(O(1/X)\)。主项正是四维单位球体积 \(\pi^2/2\)，因此复分析推导的精确系数与实几何计数相符。此处取的是累计和；不能把其误差逐项作差就得到每个 \(r_4(n)\) 的同等精度。

% !TeX root = ../../Book3.tex
\OptionalSection{二次型的 \texorpdfstring{\(\theta\)}{\AnalysisPdfTheta} 级数}{b3:sec:2405}

把平方和换成一般正定二次型，收敛和生成函数意义仍保留，但参数反演会涉及对偶格，模群也通常缩小。

\subsection{正定二次型}
\label{b3:subsec:240501}

设 \(A\) 为实对称正定矩阵，\(Q(x)=\tfrac12x^TAx\)。正定性给 \(Q(x)\ge c|x|^2\)，因而 \(\sum_{m\in\mathbb Z^d}e^{2\pi\mathrm{i}\tau Q(m)}\) 在上半平面正规收敛。若 \(A\) 为整矩阵且对角线为偶数，则 \(Q(m)\) 为整数，级数对 \(\tau\mapsto\tau+1\) 不变。

\subsection{格 \texorpdfstring{\(\theta\)}{\AnalysisPdfTheta} 级数}
\label{b3:subsec:240502}

对欧氏格 \(L\)，记 \(L^*=\{y:\langle x,y\rangle\in\mathbb Z\ (x\in L)\}\)，基本区域体积为 \(v(L)\)。格 theta 级数为
\(\Theta_L(\tau)=\sum_{x\in L}e^{\pi\mathrm{i}\tau|x|^2}\)。把实高斯的 Fourier 变换按基矩阵换元，Poisson 求和给
\[
\Theta_L(-1/\tau)=\frac{(-\mathrm{i}\tau)^{d/2}}{v(L)}\Theta_{L^*}(\tau).
\]
高维求和公式可先对标准整数格逐个坐标作周期 Fourier 展开，再线性换元；高斯的共同快速衰减保证所有换序。

\subsection{模性}
\label{b3:subsec:240503}

若 \(L=L^*\)，并且每个格向量平方长度为偶整数，则 \(v(L)=1\)，平移与反演都封闭在同一个 theta 级数上。若另有 \(d\) 是八的倍数，则 \((-\mathrm{i})^{d/2}=1\)，两个生成元给出权 \(d/2\) 的全模群变换律；尖点常数项一及指数尾部给尖点全纯性。因此 \(\Theta_L\in M_{d/2}\)。这里把维数条件列为假设，并不在本节证明偶自对偶格的维数必为八的倍数。

\subsection{表示数}
\label{b3:subsec:240504}

若 \(Q\) 在整格上取非负整数值，则 Fourier 系数等于表示数 \(r_Q(n)=\#\{m:Q(m)=n\}\)。正定性保证每层有限。对非自对偶格，反演把它送到对偶格或若干陪集级数的组合，通常得到子群上的模形式或向量值变换律；不能从一个高斯公式直接宣称所有整二次型都给全模群标量模形式。\appref{b3:app:E}将说明有限层级的相应语言。

\ProseParagraphBreak

高维 Gaussian 公式中的体积因子可在矩阵形式下直接看到。若 \(A\) 正定，选正交坐标将它对角化，再逐坐标使用一维 Gaussian 变换，得到
\[
 \widehat{e^{-\pi t x^TAx}}(\xi)
 =t^{-d/2}(\det A)^{-1/2}e^{-\pi\xi^TA^{-1}\xi/t}.
\]
正定性确保特征值全正、Gaussian 可积；逆矩阵出现在频率的二次型中，行列式来自实变量换元。因此反演自然交换原格与对偶格，而不会对任意正定整矩阵保持同一个级数。若二次型只是半正定，例如 \(Q(x,y)=x^2\)，每个固定表示层中 \(y\) 可任意变化，系数已经无穷；若不定，例如 \(x^2-y^2\)，零层也无穷。正定性不仅用于估计，还保证生成函数本身的计数意义。
\ProseParagraphBreak
偶自对偶条件有一个具体的八维实例。令
\[
 D_8=\{x\in\mathbb Z^8:\textstyle\sum_jx_j\text{ 为偶数}\},\qquad
 L=D_8\cup\left(D_8+\left(\tfrac12,\ldots,\tfrac12\right)\right).
\]
两次半整数平移回到 \(D_8\)，故并集是加法群。整数部分的平方长度模二等于坐标和，因而偶；半整数部分 \(x+v\) 的平方长度为 \(|x|^2+\sum x_j+2\)，也为偶整数。由极化恒等式，任意两格向量的内积为整数，所以 \(L\subset L^*\)。\(D_8\) 在整数格中指标为二，体积二；\(L\) 又含它为指标二子格，故体积一。对偶格体积为倒数一，包含关系与相同体积迫使 \(L=L^*\)。这给出本节模性结论所需的全部假设。
\ProseParagraphBreak
还可独立数出其最短非零向量。整数部分长度平方为二的向量恰有两个坐标为 \(\pm1\)，共 \(\binom82\cdot4=112\) 个。半整数部分若有坐标绝对值至少 \(3/2\)，平方长度已至少四；所以长度平方二时八个坐标全为 \(\pm1/2\)。属于指定陪集要求负号数为偶数，恰有 \(2^7=128\) 个。总数为240，与权四 Eisenstein 级数的一次系数相符。这个几何核算检验了长度平方 \(2n\) 与模变量 \(e^{2\pi\mathrm{i}\tau}\) 中 \(q^n\) 的对应，也给抽象的偶自对偶条件一个可以逐坐标检查的模型。

% 本章小结（不计入编号）
\begin{chaptersummary}
二平方和的障碍来自素因子在 Gaussian 整数中的不同分解方式；四平方和的精确数量则来自 theta 乘积的对数导数。后者的关键连接是：先用乘积确定单零点和 Jacobi 导数常数，再比较同一空间格上的椭圆函数，在两个半周期取值，最后用热方程转成参数导数。这样得到的 Lambert 级数可以按约数整理，正文证明不依赖章末练习。

级数、乘积、空间格和模形式各自携带归一化。尤其本章 \(q=e^{\pi\mathrm{i}\tau}\)，而偶格的自然计数变量为 \(q_{\mathrm{mod}}=q^2\)。一般格始终有包含对偶格和基本区域体积的反演公式；只有进一步核对偶性、自对偶性及权数后，才能直接使用前章全模群模形式的维数结论。累计格点数的体积近似又说明，同一表示数既有解析结构，也有几何意义，但累计估计不能替代单个系数的精确公式。
\end{chaptersummary}

\BookExercises

% 八题重新组织：前六题使用基本路线，后两题使用选读格理论。
% 二平方和、三角数、热核、反演固定点、椭圆退化和倍参数恒等式均由正文工具推导；
% 末两题考察模变换的条件与模形式有限维性的应用，不承担正文证明依赖。
\begin{exercise}\label{b3:ex:24-1}
对正整数 \(n\)，令
\[
P_2(n)=\#\{(x,y)\in\mathbb Z^2:x^2+y^2=n,\ 
                                  \gcd(|x|,|y|)=1\}.
\]
规定 \(\gcd(a,0)=a\)（\(a>0\)），不同次序和符号分别计数。求 \(P_2(n)\) 的素因子公式，并比较 \(P_2(25)\) 与 \(r_2(25)\)。
\end{exercise}
\begin{solution}
若 \(n\) 含模四为三的素因子 \(\ell\)，则其指数为奇数时没有任何表示；指数为正偶数 \(2s\) 时，正文的分解公式使每个 \(\alpha=x+\mathrm{i}y\) 都含实因子 \(\ell^s\)，从而 \(\ell\mid x,y\)。所以原始表示中完全不能出现这种素因子。

写剩余的分解为
\[
n=2^a\prod_{j=1}^k p_j^{a_j},
\qquad p_j\equiv1\pmod4,\quad a_j\ge1.
\]
正文已证 \(\alpha\) 中 \(1+\mathrm{i}\) 的指数恰为 \(a\)。若 \(a\ge2\)，由于 \((1+\mathrm{i})^2=2\mathrm{i}\)，两个坐标均为偶数；若 \(a\le1\)，则二不能同时整除两个坐标。

对每个 \(p_j=\pi_j\bar\pi_j\)，\(\alpha\) 中的因子为
\(\pi_j^{r_j}\bar\pi_j^{a_j-r_j}\)。
实整数 \(p_j\) 整除 \(\alpha\)，当且仅当两个指数都至少为一，这又等价于 \(p_j\) 同时整除 \(x,y\)。故原始性要求
\(r_j=0\) 或 \(r_j=a_j\)，每个不同奇素因子给两种选择。反过来，若这些条件都满足，则任何可能同时整除两坐标的有理素数都已被排除，故最大公因子为一。四个单位分别计数，得到
\[
P_2(n)=
\begin{cases}
4\cdot2^k,&n=2^a\prod_{j=1}^k p_j^{a_j},\
                 a\in\{0,1\},\ p_j\equiv1\pmod4,\\
0,&\text{其余情形}.
\end{cases}
\]
空乘积允许 \(n=1,2\)，此时均有四个原始表示。对 \(25\)，
\(P_2(25)=8\)，对应绝对值型 \((3,4)\)；而 \(r_2(25)=12\)，另外四个为 \((\pm5,0),(0,\pm5)\)，它们的坐标有公因子五。
\end{solution}

\begin{exercise}\label{b3:ex:24-2}
令 \(T(a)=a(a+1)/2\)（\(a\in\mathbb Z_{\ge0}\)）。对整数 \(n\ge0\)，用 \(t_4(n)\) 表示满足
\(T(a_1)+T(a_2)+T(a_3)+T(a_4)=n\) 的有序非负整数四元组数。证明
\[
t_4(n)=\sigma_1(2n+1),
\qquad \sigma_1(m)=\sum_{\substack{d\mid m\\d>0}}d.
\]
说明这如何推出每个非负整数都是四个非负三角数之和。
\end{exercise}
\begin{solution}
恒等式 \(8T(a)+1=(2a+1)^2\) 将原问题化为
\[
(2a_1+1)^2+\cdots+(2a_4+1)^2=8n+4.
\]
令 \(m=2n+1\)，则 \(m\) 为正奇数。四个平方之和为 \(4m\) 时，平方模四的值为零或一，所以奇坐标的个数只能为零或四。因此全部表示恰分为全偶和全奇两类。全偶表示除以二，与 \(m\) 的四平方和表示一一对应，数量为 \(r_4(m)\)；全奇表示数就是
\(r_4(4m)-r_4(m)\)。

因为 \(m\) 为奇数，Jacobi 公式给
\[
r_4(m)=8\sigma_1(m),\qquad
r_4(4m)=8\!\sum_{d\mid m}(d+2d)=24\sigma_1(m).
\]
后式中允许的约数只有 \(d,2d\)，所有 \(4d\) 被排除，故全奇表示共有 \(16\sigma_1(m)\) 个。

每个全奇表示的四个坐标都非零。独立改变四个符号给恰好十六个不同表示，而取各坐标绝对值以后，
\(a_j=(|x_j|-1)/2\) 唯一确定。除以十六，便得所求公式。由于 \(\sigma_1(2n+1)\ge1\)，表示总存在；\(n=0\) 时公式给一，正对应四个 \(a_j\) 都为零。
\end{solution}

\begin{exercise}\label{b3:ex:24-3}
对 \(t>0\)、\(x\in\mathbb R\)，令
\[
K_t(x)=\frac1\pi\theta_3(x,\mathrm{i}t).
\]
证明 \(K_t\) 为正的 \(\pi\)-周期函数，
\(\int_0^\pi K_t(x)\,\mathrm dx=1\)，并满足
\[
\int_0^\pi K_s(x-y)K_t(y)\,\mathrm dy=K_{s+t}(x)
\quad(s,t>0).
\]
若 \(f\) 为连续的 \(\pi\)-周期函数，证明
\(\int_0^\pi K_t(x-y)f(y)\,\mathrm dy\) 在 \(t\downarrow0\) 时一致趋于 \(f(x)\)。
\end{exercise}
\begin{solution}
在高斯 Poisson 公式中把实参数换成 \(1/t\)，取平移参数 \(a=-x/\pi\)、调制参数 \(b=0\)，再用 \(\theta_3\) 关于空间变量的偶性，得
\[
K_t(x)=\frac1{\pi\sqrt t}
       \sum_{m\in\mathbb Z}
       \exp\!\left(-\frac{(x-\pi m)^2}{\pi t}\right).
\]
每一项都正，平移 \(x\mapsto x+\pi\) 只重排指标。对非负级数逐项积分，将各个长度为 \(\pi\) 的区间拼成实轴，得到
\[
\int_0^\pi K_t(x)\,\mathrm dx
=\frac1{\pi\sqrt t}\int_{\mathbb R}e^{-u^2/(\pi t)}\,\mathrm du=1.
\]
这也解释了前因子 \(1/\pi\)：原来的 theta 函数在一个空间周期上的积分为 \(\pi\)。

由原级数
\[
K_t(x)=\frac1\pi\sum_{n\in\mathbb Z}
               e^{-\pi n^2t}e^{2\mathrm{i}nx}.
\]
当 \(s,t>0\) 固定时，两个级数绝对一致收敛，其乘积可逐项积分。利用
\(\int_0^\pi e^{2\mathrm{i}(m-n)y}\,\mathrm dy
=\pi\)（\(m=n\)）或零（\(m\ne n\)），只留下相同频率，系数相乘成为
\(e^{-\pi n^2(s+t)}\)，即所述卷积公式。

最后把一个周期取为 \([-\pi/2,\pi/2]\)。对 \(0<\delta<\pi/2\)，高斯表达式给
\[
\int_{\delta\le |y|\le\pi/2}K_t(y)\,\mathrm dy
\le \frac1{\pi\sqrt t}
             \int_{|u|\ge\delta}e^{-u^2/(\pi t)}\,\mathrm du
\le \frac{\sqrt t}{\delta}e^{-\delta^2/(\pi t)}.
\]
第一个不等式把各平移区间合并，并扩大到 \(|u|\ge\delta\)；
第二个用 \(u/\delta\ge1\) 估计正半轴积分，再利用对称性。

连续周期函数在实轴上一致连续。给定 \(\varepsilon>0\)，可先选 \(\delta\)，使
\(|f(x-y)-f(x)|<\varepsilon\) 对所有 \(x\) 及 \(|y|<\delta\) 成立。利用质量为一，卷积与 \(f(x)\) 之差的绝对值至多为
\[
\varepsilon+
2\|f\|_\infty\frac{\sqrt t}{\delta}e^{-\delta^2/(\pi t)}.
\]
右侧与 \(x\) 无关，再令 \(t\downarrow0\)，即得一致收敛。这同时给出了本章热方程中周期初值的一个具体构造。
\end{solution}

\begin{exercise}\label{b3:ex:24-4}
设 \(T(t)=\sum_{n\in\mathbb Z}e^{-\pi n^2t}\)（\(t>0\)）。
利用 \(T(t)=t^{-1/2}T(1/t)\) 证明
\[
\sum_{n\in\mathbb Z}n^2e^{-\pi n^2}
=\frac1{4\pi}\sum_{n\in\mathbb Z}e^{-\pi n^2}.
\]
再说明函数 \(H(s)=e^{s/4}T(e^s)\)（\(s\in\mathbb R\)）为何是偶函数，并用它解释上式中的固定参数 \(t=1\)。
\end{exercise}
\begin{solution}
在包含一的紧区间上，导数级数由常数倍的
\(\sum n^2e^{-c n^2}\) 控制，故可逐项求导。
反演式微分为
\[
T'(t)=-\tfrac12t^{-3/2}T(1/t)-t^{-5/2}T'(1/t).
\]
取 \(t=1\)，得到 \(2T'(1)=-T(1)/2\)。
代入 \(T'(1)=-\pi\sum n^2e^{-\pi n^2}\)，就是所求恒等式。

另一方面，反演式在 \(t=e^s\) 处给
\(T(e^s)=e^{-s/2}T(e^{-s})\)，因此
\[
H(s)=e^{-s/4}T(e^{-s})=H(-s).
\]
参数反演 \(t\mapsto1/t\) 在对数坐标中成为反射
\(s\mapsto-s\)，其固定点为 \(s=0\)，即 \(t=1\)。
偶性要求 \(H'(0)=0\)，而
\(H'(0)=T(1)/4+T'(1)\)，再次给出同一矩恒等式。
\end{solution}

\begin{exercise}\label{b3:ex:24-5}
设 \(\tau\in\mathfrak H\)、\(q=e^{\pi\mathrm{i}\tau}\)，
\(\Lambda_\tau=\pi\mathbb Z+\pi\tau\mathbb Z\)。
在 \(|\operatorname{Im}z|<\pi\operatorname{Im}\tau\)、\(z\notin\pi\mathbb Z\) 内证明
\[
\wp(z;\Lambda_\tau)=\frac1{\sin^2z}-\frac13
 +8\sum_{m\ge1}\frac{m q^{2m}}{1-q^{2m}}
                         \bigl(1-\cos(2mz)\bigr).
\]
进而证明，当 \(t\to\infty\) 时，
\[
\wp(z;\pi\mathbb Z+\pi\mathrm{i}t\mathbb Z)
\longrightarrow\frac1{\sin^2z}-\frac13
\]
在 \(\mathbb C\setminus\pi\mathbb Z\) 的每个紧集上一致成立。
\end{exercise}
\begin{solution}
正文的对数导数公式在所述条带内为
\[
\frac{\theta_1'(z)}{\theta_1(z)}
=\cot z+4\sum_{m\ge1}
              \frac{q^{2m}}{1-q^{2m}}\sin(2mz).
\]
在条带的紧子集上，各阶导数都受带多项式因子的几何级数控制，所以可以求导，也可以在原点比较 Laurent 展开的系数。
由 \(\cot z=z^{-1}-z/3+O(z^3)\)，上式中 \(z\) 的系数为
\[
-\frac13+8\sum_{m\ge1}\frac{m q^{2m}}{1-q^{2m}}.
\]
另一方面，正文展开
\(\theta_1(z)=Az+Bz^3+O(z^5)\) 给出该系数为
\(2B/A=\theta_1'''(0)/(3\theta_1'(0))\)。
将此值代入\cref{b3:prop:24-theta-wp}，并对上面的对数导数求导，便得到所求公式。常数 \(-1/3\) 正是使原点 Laurent 展开的常数项为零所必需的修正。

设紧集 \(K\subset\mathbb C\setminus\pi\mathbb Z\)，并令
\(R=\max_{z\in K}|\operatorname{Im}z|\)。当 \(t>R/\pi\) 时，\(K\) 位于公式的条带内。取 \(\tau=\mathrm{i}t\)，则
\(q=e^{-\pi t}\)，且
\[
|1-\cos(2mz)|\le1+e^{2mR}\quad(z\in K).
\]
对于充分大的 \(t\)，分母 \(1-e^{-2\pi mt}\) 一致远离零，所以级数尾项的总和由
\[
C\sum_{m\ge1}m e^{-2\pi mt}(1+e^{2mR})
=O_K(e^{-2\pi t})
\]
控制。这给出紧集上一致收敛，甚至给出指数误差。几何上，第二个周期向虚方向变长，有限区域内留下的极点排列趋于实轴上的 \(\pi\mathbb Z\)。
\end{solution}

\begin{exercise}\label{b3:ex:24-6}
对 \(\tau\in\mathfrak H\)，记 \(\theta_j(\tau)=\theta_j(0,\tau)\)。
直接重排绝对收敛的双重级数，证明倍参数恒等式
\[
\begin{aligned}
\theta_3(\tau)^2&=\theta_3(2\tau)^2+\theta_2(2\tau)^2,\\
\theta_4(\tau)^2&=\theta_3(2\tau)^2-\theta_2(2\tau)^2,\\
\theta_2(\tau)^2&=2\theta_2(2\tau)\theta_3(2\tau).
\end{aligned}
\]
据此另行推出 \(\theta_3(\tau)^4=\theta_2(\tau)^4+\theta_4(\tau)^4\)。
\end{exercise}
\begin{solution}
在 \(\theta_3(\tau)^2\) 中，令
\(u=m+n,\ v=m-n\)。整数对 \((m,n)\) 与同奇偶的整数对 \((u,v)\) 一一对应，逆变换为
\(m=(u+v)/2,\ n=(u-v)/2\)，且
\(m^2+n^2=(u^2+v^2)/2\)。
当 \(u=2a,v=2b\) 时，贡献为
\(\sum_{a,b}q^{2a^2+2b^2}=\theta_3(2\tau)^2\)；
当 \(u=2a+1,v=2b+1\) 时，贡献为
\(\sum_{a,b}q^{2(a+1/2)^2+2(b+1/2)^2}=\theta_2(2\tau)^2\)。
这证明第一式。

\(\theta_4(\tau)^2\) 多出符号
\((-1)^{m+n}=(-1)^u\)，在全偶部分为正、全奇部分为负，故得到第二式。
对 \(\theta_2(\tau)^2\)，改用
\[
u=(m+\tfrac12)+(n+\tfrac12)=m+n+1,\qquad
v=(m+\tfrac12)-(n+\tfrac12)=m-n.
\]
此时 \(u,v\) 为异奇偶整数，逆变换仍唯一。
偶奇和奇偶两部分各给一个
\(\theta_3(2\tau)\theta_2(2\tau)\)，于是得到第三式。
所有这些分组都由高斯型双重级数的绝对收敛保证，不存在条件收敛下的重排问题。

最后将前两式相减、相加，再相乘：
\[
\theta_3(\tau)^4-\theta_4(\tau)^4
=4\theta_2(2\tau)^2\theta_3(2\tau)^2
=\theta_2(\tau)^4.
\]
这提供了四次恒等式的另一条路线，并说明倍参数公式与正文的奇偶格点变换有共同的计数基础。
\end{solution}

\begin{optionalexercise}\label{b3:ex:24-7}
在标准欧氏平面中取
\[
L=\operatorname{diag}(\sqrt2,1/\sqrt2)\mathbb Z^2.
\]
求其基本区域体积与对偶格，证明 \(L\ne L^*\)，但
\(\Theta_L=\Theta_{L^*}\)。写出 \(\Theta_L(-1/\tau)\) 的变换律，并证明 \(\Theta_L\) 不满足 \(\tau\mapsto\tau+1\) 的不变性。由此解释，为何体积为一，甚至反演公式在同一个函数上闭合，仍不足以得到全模群模形式。
\end{optionalexercise}
\begin{solution}
基矩阵的行列式为一，所以 \(v(L)=1\)，而逆转置给
\[
L^*=\operatorname{diag}(1/\sqrt2,\sqrt2)\mathbb Z^2.
\]
向量 \((0,1/\sqrt2)\) 属于 \(L\) 而不属于 \(L^*\)，故两格不同；交换两个坐标却将 \(L\) 等距映到 \(L^*\)，所以长度计数相同。由定义，
\[
\Theta_L(\tau)
=\sum_{m,n\in\mathbb Z}e^{\pi\mathrm{i}\tau(2m^2+n^2/2)}
=\theta_3(0,2\tau)\theta_3(0,\tau/2).
\]
一般格反演在维数二给
\(\Theta_L(-1/\tau)=(-\mathrm{i}\tau)\Theta_{L^*}(\tau)
=-\mathrm{i}\tau\Theta_L(\tau)\)。

但 \(L\) 的最短非零向量为 \((0,\pm1/\sqrt2)\)，平方长度是 \(1/2\)，所以当 \(\operatorname{Im}\tau\to\infty\) 时，
\[
\Theta_L(\tau)=1+2e^{\pi\mathrm{i}\tau/2}
                 +O(e^{-2\pi\operatorname{Im}\tau}).
\]
这个估计对 \(\operatorname{Re}\tau=0,1\) 均成立，由绝对收敛级数中剩余平方长度至少为二得到。
平移 \(\tau\mapsto\tau+1\) 将显示出的第二项乘以 \(\mathrm{i}\)，因而
\[
\Theta_L(1+\mathrm{i}t)-\Theta_L(\mathrm{i}t)
=2(\mathrm{i}-1)e^{-\pi t/2}+O(e^{-2\pi t})\ne0
\]
对充分大的 \(t\) 成立。
这里体积为一没有保证自对偶，也没有保证长度平方为偶整数。反演中的额外常数 \(-\mathrm{i}\) 和平移失败都表明，不能直接套用前章全模群的整数权变换律。
\end{solution}

\begin{optionalexercise}\label{b3:ex:24-8}
设 \(L\) 是标准欧氏空间中的正定偶自对偶格，记
\[
R_L(n)=\#\{v\in L:|v|^2=2n\},\qquad
Q=e^{2\pi\mathrm{i}\tau}.
\]
当 \(\dim L=16\) 时，求 \(R_L(1)\) 与 \(R_L(2)\)。
当 \(\dim L=24\) 时，用 \(R_L(1)\) 表示整个 theta 级数，并证明
\(R_L(2)=196560-24R_L(1)\)。
特别地，若某个二十四维偶自对偶格没有平方长度为二的向量，它有多少个平方长度为四的向量？本题不要求证明这样的格存在。
\end{optionalexercise}
\begin{solution}
末节已证 \(\Theta_L\in M_{\dim L/2}\)，而正定性使常数项为一，偶性使展开为
\(\Theta_L=1+\sum_{n\ge1}R_L(n)Q^n\)。
这里使用 \(Q=q^2\)，所以 \(Q^n\) 计平方长度 \(2n\)，不能读成平方长度 \(n\)。

由\cref{b3:thm:23-modular-ring}，\(M_8=\mathbb C E_4^2\)。
十六维情形比较常数项便有 \(\Theta_L=E_4^2\)。
又由\eqref{b3:eq:23-e4-e6-fourier}，
\[
E_4=1+240Q+2160Q^2+O(Q^3).
\]
平方后得
\[
E_4^2=1+480Q+61920Q^2+O(Q^3).
\]
因此 \(R_L(1)=480,\ R_L(2)=61920\)，这些数量与具体选择哪个十六维偶自对偶格无关。

二十四维时，\(M_{12}=\mathbb C E_4^3\oplus\mathbb C\Delta\)。
常数项一先确定 \(E_4^3\) 的系数为一，再比较一次项，得到
\[
\Theta_L=E_4^3+(R_L(1)-720)\Delta.
\]
为求第二个系数，使用正文已算出的
\[
E_4^3=1+720Q+179280Q^2+O(Q^3),\qquad
\Delta=Q-24Q^2+O(Q^3).
\]
于是
\[
R_L(2)=179280-24(R_L(1)-720)
       =196560-24R_L(1).
\]
若 \(R_L(1)=0\)，就有恰好 \(196560\) 个平方长度为四的向量。这个推论从给定格的条件出发；它计算可能的长度分布，并不构造或分类这样的格。
\end{solution}

